% Le vandalisme et le terrorisme — Allemand 1ère A — Cours signé OnBuch+
% Programme officiel MINESEC. Compiler avec : tectonic ALA25-le-vandalisme-et-le-terrorisme.tex
\newcommand{\DOCMATIERE}{Allemand}
\newcommand{\DOCNIVEAU}{1ère A}
\newcommand{\DOCMODULE}{Chapitre 5 : Citoyenneté}
\newcommand{\DOCLECON}{ALA25}
\newcommand{\DOCTITRE}{Le vandalisme et le terrorisme}
% OnBuch+ — préambule « cours signés » ENRICHI (compilé avec tectonic / XeLaTeX)
% Système visuel « Pop » d'OnBuch+ : fond crème (jamais blanc), bordures encre
% épaisses, ombres « dures » décalées, titres Archivo Black en capitales.
% Version MATHÉMATIQUES : graphes (pgfplots/tikz), boîtes pédagogiques
% (définition, propriété, méthode, attention, le savais-tu, exercice résolu,
% exercice, TP, bilan), page de garde et signature OnBuch+ ancrée en bas.
%
% Macros attendues dans main.tex AVANT \input{preamble} :
%   \DOCMATIERE  \DOCNIVEAU  \DOCMODULE  \DOCLECON (numéro)  \DOCTITRE
\documentclass[11pt]{article}

\usepackage{fontspec}
\usepackage[ngerman,french]{babel} % français = langue principale ; l'allemand via \all{..} / textebox
\usepackage[a4paper,margin=2cm,top=3.4cm,bottom=2.8cm,headheight=1.4cm,footskip=1.3cm]{geometry}
\usepackage{amsmath,amssymb,amsfonts}
\usepackage{mathtools} % \xmapsto, \coloneqq... souvent utilisés par les modèles
\usepackage{cancel} % simplifications barrées (bilans de forces)
% \abs{z} (module, valeur absolue) et \norm{u} (norme), utilisés par les modèles
\providecommand{\abs}{}\let\abs\relax\DeclarePairedDelimiter{\abs}{\lvert}{\rvert}
\providecommand{\norm}{}\let\norm\relax\DeclarePairedDelimiter{\norm}{\lVert}{\rVert}
\usepackage[table]{xcolor} % [table] : fournit aussi \rowcolors (alternance de lignes)
\usepackage{pagecolor}
\usepackage{tikz}
\usetikzlibrary{arrows.meta,positioning,calc,shapes.geometric,shapes.misc,shapes.arrows,decorations.pathmorphing,decorations.pathreplacing,decorations.markings,patterns,shadows,mindmap,trees,fit,backgrounds,matrix,intersections,angles,quotes}
\usepackage{pgfplots}
\usepackage[version=4]{mhchem} % équations nucléaires \ce{^{235}_{92}U}
\usepackage[european,siunitx]{circuitikz} % schémas électriques
\pgfplotsset{compat=1.18}
\usepgfplotslibrary{fillbetween,groupplots} % aires entre courbes (calcul intégral)
\usepackage{enumitem}
\usepackage{fancyhdr}
% [most] charge toutes les bibliothèques tcolorbox (listings, minted, raster,
% documentation...) et gonfle inutilement la mémoire TeX sur un document long
% (~300 boîtes) : on ne charge que ce qui est réellement utilisé.
\usepackage[skins,breakable]{tcolorbox}
\usepackage{graphicx}
\usepackage{colortbl}
\usepackage{booktabs}
\usepackage{tabularx}
\usepackage{diagbox} % case d'en-tête diagonale des tableaux à double entrée
% Colonnes extensibles centrée / gauche / droite (C, L, R), souvent utilisées par les modèles
\newcolumntype{C}{>{\centering\arraybackslash}X}
\newcolumntype{L}{>{\raggedright\arraybackslash}X}
\newcolumntype{R}{>{\raggedleft\arraybackslash}X}
\usepackage{array}
\usepackage{multicol}
\usepackage{multirow}
\usepackage{siunitx}
% parse-numbers=false : accepte aussi \SI{2,5 \times 10^{-3}}{...}
\sisetup{locale=FR,output-decimal-marker={,},group-minimum-digits=4,parse-numbers=false}
\DeclareSIUnit\ohm{\ensuremath{\symup{\Omega}}} % U+2126 absent de la police
\DeclareSIUnit\division{div} % réglages d'oscilloscope (ms/div, V/div)
\DeclareSIUnit\tr{tr} % tours (vitesse de rotation, ex. \SI{1500}{\tr\per\minute})
\DeclareSIUnit\molar{M} % concentration molaire (ex. \SI{3}{\molar}, \SI{1}{\micro\molar})
\DeclareSIUnit\an{an} % année (le modèle confond parfois avec \year, un compteur TeX) — ex. \SI{5}{\kilo\meter\per\an}
\DeclareSIUnit\FCFA{FCFA} % franc CFA (ex. \SI{500}{\FCFA\per\kilo\gram})
\DeclareSIUnit\degreeBrix{\textdegree Brix} % teneur en sucre d'un sirop/jus (réfractométrie)
\DeclareSIUnit\month{mois} % le modèle utilise parfois \month, un compteur TeX, comme unité de durée
\usepackage{qrcode}
% \degree : commande fréquemment utilisée par les modèles pour un angle ou une
% température en dehors de \SI{}{} (non fournie par les paquets chargés).
\providecommand{\degree}{^{\circ}}
% \qed (fin de démonstration), utilisé par les modèles sans amsthm
\providecommand{\qed}{\hfill$\square$}
% \barre : double barre des valeurs interdites dans un tableau de variations (tkz-tab)
\providecommand{\barre}{\ensuremath{\|}}
% environnement proof (amsthm non chargé) utilisé par les modèles
\ifdefined\proof\else\newenvironment{proof}[1][Démonstration]{\par\noindent\textit{#1.}\ }{\qed\par}\fi
\usepackage{lastpage}
\usepackage{hyperref}
\hypersetup{hidelinks}

% ============================================================
% Palette « Pop » OnBuch+ — mêmes hex que la classe P (plus_ui.dart)
% ============================================================
\definecolor{poporange}{HTML}{F59321}
\definecolor{popdark}{HTML}{E07A0C}
\definecolor{popcream}{HTML}{FFF6E8}
\definecolor{popink}{HTML}{1C1714}
\definecolor{poppurple}{HTML}{7A5AE0}
\definecolor{popgreen}{HTML}{2E9E6A}
\definecolor{poppink}{HTML}{E0457B}
\definecolor{popblue}{HTML}{2F7DE1}
\definecolor{popgold}{HTML}{C98A12}
\definecolor{popmuted}{HTML}{8A7F76}
\definecolor{popline}{HTML}{EADFD2}
\definecolor{popgray}{HTML}{8A7F76}
\colorlet{popgrey}{popgray}
% Alias tolérants (le modèle nomme parfois les couleurs différemment)
\colorlet{popwhite}{white}
\colorlet{popblack}{popink}
\colorlet{popred}{poppink}
\colorlet{popyellow}{popgold}
% Teintes claires (fonds de tableaux/encadrés) — le modèle nomme parfois
% les couleurs avec un suffixe L (ex. popblueL) sans les avoir définies.
\colorlet{poporangeL}{poporange!20}
\colorlet{popdarkL}{popdark!20}
\colorlet{poppurpleL}{poppurple!20}
\colorlet{popgreenL}{popgreen!20}
\colorlet{poppinkL}{poppink!20}
\colorlet{popblueL}{popblue!20}
\colorlet{popgoldL}{popgold!20}
\colorlet{popredL}{popred!20}
\colorlet{popgrayL}{popgray!20}
% Teintes claires pour fonds de tableaux / zones
\colorlet{popblueL}{popblue!10!white}
\colorlet{poporangeL}{poporange!14!white}
\colorlet{popgreenL}{popgreen!12!white}
\colorlet{poppurpleL}{poppurple!12!white}
\colorlet{poppinkL}{poppink!12!white}
\colorlet{popgoldL}{popgold!14!white}

\pagecolor{popcream}

% ============================================================
% Typographie
% ============================================================
\defaultfontfeatures{Ligatures=TeX}
\setmainfont{PlusJakartaSans-Medium.ttf}[Path=../../../fonts/,BoldFont=PlusJakartaSans-Bold.ttf]
% Maths en sans-serif (Fira Math) avec chiffres et lettres droites de Plus
% Jakarta, pour que les formules se fondent dans le texte.
\usepackage[math-style=ISO,bold-style=ISO]{unicode-math}
\setmathfont{FiraMath-Regular.otf}[Path=../../../fonts/]
\setmathfont{PlusJakartaSans-Medium.ttf}[Path=../../../fonts/,range={up/num,up/{Latin,latin},\mathup}]
\usepackage{icomma} % virgule décimale sans espace en mode math
% Caractères mathématiques Unicode absents des polices de texte (Plus Jakarta,
% Archivo Black) : rendus par leur équivalent mathématique, en texte comme en
% formule — sinon ils s'affichent en carré vide (ex. « Arithmétique dans ℤ »).
\usepackage{newunicodechar}
\newunicodechar{ℝ}{\ensuremath{\mathbb{R}}}
\newunicodechar{ℤ}{\ensuremath{\mathbb{Z}}}
\newunicodechar{ℂ}{\ensuremath{\mathbb{C}}}
\newunicodechar{ℕ}{\ensuremath{\mathbb{N}}}
\newunicodechar{ℚ}{\ensuremath{\mathbb{Q}}}
\newunicodechar{𝔻}{\ensuremath{\mathbb{D}}}
\newunicodechar{α}{\ensuremath{\alpha}}
\newunicodechar{β}{\ensuremath{\beta}}
\newunicodechar{γ}{\ensuremath{\gamma}}
\newunicodechar{δ}{\ensuremath{\delta}}
\newunicodechar{ε}{\ensuremath{\varepsilon}}
\newunicodechar{θ}{\ensuremath{\theta}}
\newunicodechar{λ}{\ensuremath{\lambda}}
\newunicodechar{μ}{\ensuremath{\mu}}
\newunicodechar{σ}{\ensuremath{\sigma}}
\newunicodechar{τ}{\ensuremath{\tau}}
\newunicodechar{φ}{\ensuremath{\varphi}}
\newunicodechar{ψ}{\ensuremath{\psi}}
\newunicodechar{ω}{\ensuremath{\omega}}
\newunicodechar{Σ}{\ensuremath{\Sigma}}
% majuscules produites par \MakeUppercase dans les titres (α-aminés -> Α)
\newunicodechar{Α}{\ensuremath{\alpha}}
\newunicodechar{Β}{\ensuremath{\beta}}
\newunicodechar{Γ}{\ensuremath{\Gamma}}
\newunicodechar{Δ}{\ensuremath{\Delta}}
\newunicodechar{Ε}{\ensuremath{\varepsilon}}
\newunicodechar{Θ}{\ensuremath{\Theta}}
\newunicodechar{Λ}{\ensuremath{\Lambda}}
\newunicodechar{Μ}{\ensuremath{\mu}}
\newunicodechar{Τ}{\ensuremath{\tau}}
\newunicodechar{Φ}{\ensuremath{\Phi}}
\newunicodechar{Ψ}{\ensuremath{\Psi}}
\newunicodechar{Ω}{\ensuremath{\Omega}}
\newunicodechar{↦}{\ensuremath{\mapsto}}
\newunicodechar{∈}{\ensuremath{\in}}
\newunicodechar{∉}{\ensuremath{\notin}}
\newunicodechar{∧}{\ensuremath{\wedge}}
\newunicodechar{⇔}{\ensuremath{\Leftrightarrow}}
\newunicodechar{⟺}{\ensuremath{\Longleftrightarrow}}
\newunicodechar{⇒}{\ensuremath{\Rightarrow}}
\newunicodechar{⟹}{\ensuremath{\Longrightarrow}}
\newunicodechar{∘}{\ensuremath{\circ}}
\newunicodechar{∪}{\ensuremath{\cup}}
\newunicodechar{∩}{\ensuremath{\cap}}
\newunicodechar{≡}{\ensuremath{\equiv}}
\newunicodechar{⊂}{\ensuremath{\subset}}
\newunicodechar{⊥}{\ensuremath{\perp}}
\newunicodechar{∃}{\ensuremath{\exists}}
\newunicodechar{∀}{\ensuremath{\forall}}
\newunicodechar{∛}{\ensuremath{\sqrt[3]{\,}}}
\newunicodechar{∜}{\ensuremath{\sqrt[4]{\,}}}
\newunicodechar{⃗}{\ensuremath{{}^{\to}}}
\newunicodechar{ⁿ}{\ifmmode^{n}\else\textsuperscript{\ensuremath{n}}\fi}
\newunicodechar{ˣ}{\ifmmode^{x}\else\textsuperscript{\ensuremath{x}}\fi}
\newunicodechar{ᵇ}{\ifmmode^{b}\else\textsuperscript{\ensuremath{b}}\fi}
\newunicodechar{ʳ}{\ifmmode^{r}\else\textsuperscript{\ensuremath{r}}\fi}
\newunicodechar{ᵘ}{\ifmmode^{u}\else\textsuperscript{\ensuremath{u}}\fi}
\newunicodechar{ʸ}{\ifmmode^{y}\else\textsuperscript{\ensuremath{y}}\fi}
\newunicodechar{ᵃ}{\ifmmode^{a}\else\textsuperscript{\ensuremath{a}}\fi}
\newunicodechar{ᵉ}{\ifmmode^{e}\else\textsuperscript{\ensuremath{e}}\fi}
\newunicodechar{ᵏ}{\ifmmode^{k}\else\textsuperscript{\ensuremath{k}}\fi}
\newunicodechar{ᵖ}{\ifmmode^{p}\else\textsuperscript{\ensuremath{p}}\fi}
\newunicodechar{ᴺ}{\ifmmode^{N}\else\textsuperscript{\ensuremath{N}}\fi}
\newunicodechar{ᶜ}{\ifmmode^{c}\else\textsuperscript{\ensuremath{c}}\fi}
\newunicodechar{ʰ}{\ifmmode^{h}\else\textsuperscript{\ensuremath{h}}\fi}
\newunicodechar{ᵠ}{\ifmmode^{q}\else\textsuperscript{\ensuremath{q}}\fi}
\newunicodechar{ₙ}{\ifmmode_{n}\else\textsubscript{\ensuremath{n}}\fi}
\newunicodechar{ₐ}{\ifmmode_{a}\else\textsubscript{\ensuremath{a}}\fi}
\newunicodechar{ᵢ}{\ifmmode_{i}\else\textsubscript{\ensuremath{i}}\fi}
\newunicodechar{ᵣ}{\ifmmode_{r}\else\textsubscript{\ensuremath{r}}\fi}
\newunicodechar{ₖ}{\ifmmode_{k}\else\textsubscript{\ensuremath{k}}\fi}
\newunicodechar{ᵦ}{\ifmmode_{\beta}\else\textsubscript{\ensuremath{\beta}}\fi}
\newunicodechar{ₓ}{\ifmmode_{x}\else\textsubscript{\ensuremath{x}}\fi}
\newfontfamily\popdisplay{ArchivoBlack-Regular.ttf}[Path=../../../fonts/]
\newfontfamily\poplogo{SpaceGrotesk-Bold.ttf}[Path=../../../fonts/]
\newfontfamily\popsemi{PlusJakartaSans-SemiBold.ttf}[Path=../../../fonts/]
\newfontfamily\popxbold{PlusJakartaSans-ExtraBold.ttf}[Path=../../../fonts/]
\newfontfamily\popxboldls{PlusJakartaSans-ExtraBold.ttf}[Path=../../../fonts/,LetterSpace=3.0]

\setlist[enumerate]{leftmargin=1.7em,itemsep=5pt,topsep=4pt}
\setlist[itemize]{leftmargin=1.7em,itemsep=4pt,topsep=4pt,label={\color{poporange}\textbullet}}
\setlength{\parindent}{0pt}
\setlength{\parskip}{5pt}
\renewcommand{\arraystretch}{1.25}

% pgfplots : style Pop par défaut
\pgfplotsset{
  every axis/.append style={axis line style={popink,line width=1pt},tick style={popink},
    grid style={popline,line width=0.5pt},label style={font=\small},tick label style={font=\footnotesize}},
  popcurve/.style={poporange,line width=1.8pt,no markers,smooth},
}
% Même style disponible dans un \draw TikZ simple (hors pgfplots)
\tikzset{popcurve/.style={poporange,line width=1.8pt}}
\tikzset{popgrid/.style={popline,very thin}}
\colorlet{popcurve}{poporange} % le nom du style est parfois utilisé comme couleur

% ============================================================
% Pill / squiggle
% ============================================================
\newcommand{\poppill}[3]{%
  \tikz[baseline=-0.55ex]\node[draw=popink,line width=1.2pt,rounded corners=2.6pt,%
    fill=#2,inner xsep=6.5pt,inner ysep=3pt]{%
    \popxboldls\fontsize{7}{7}\selectfont\color{#3}\MakeUppercase{#1}};%
}
\newcommand{\popsquiggle}[1]{%
  \tikz[baseline=-0.4ex]{
    \draw[#1,line width=2.6pt,line cap=round]
      plot [smooth] coordinates {(0,0.09) (0.34,0.01) (0.68,0.21) (1.02,0.01) (1.36,0.10)};
  }%
}
% Mot-clé surligné (marqueur orange sous le texte)
\newcommand{\cle}[1]{{\popxbold\color{popdark}#1}}

% ============================================================
% En-tête / pied
% ============================================================
\pagestyle{fancy}
\fancyhf{}
\renewcommand{\headrulewidth}{0pt}
\renewcommand{\footrulewidth}{0pt}
\lhead{\raisebox{-0.15ex}{\poplogo\fontsize{13}{13}\selectfont\color{popink}OnBuch\color{poporange}+}}
\rhead{\poppill{\DOCMATIERE\ \textbullet\ \DOCNIVEAU\ \textbullet\ Leçon \DOCLECON}{white}{popink}}
\lfoot{\popxbold\fontsize{7.6}{7.6}\selectfont\color{popmuted}\MakeUppercase{Cours signé OnBuch+}\ \textbullet\ {\popsemi\DOCTITRE}}
\rfoot{\tikz[baseline=-0.6ex]\node[rounded corners=8.5pt,fill=popink,minimum height=17pt,minimum width=17pt,inner xsep=5pt,inner sep=0]{\popxbold\fontsize{8}{8}\selectfont\color{popcream}\thepage\,/\,\pageref*{LastPage}};}

% ============================================================
% Titres
% ============================================================
\newcommand{\chaptitre}[2]{%
  \begin{tcolorbox}[enhanced,colback=poporange,colframe=popink,boxrule=2.6pt,arc=11pt,
      drop shadow=popink,left=15pt,right=15pt,top=12pt,bottom=11pt,before skip=3pt,after skip=18pt]
    \poppill{#2}{popink}{popcream}\par\vspace{9pt}
    {\raggedright\hyphenpenalty=10000\exhyphenpenalty=10000\popdisplay\fontsize{22}{24}\selectfont\color{popcream}\MakeUppercase{#1}\par}
  \end{tcolorbox}
}
% Section numérotée : pastille orange avec le numéro + titre Archivo
\newcounter{popsec}
\newcommand{\coursec}[1]{%
  \stepcounter{popsec}\setcounter{popsub}{0}%
  \par\vspace{16pt}\noindent
  \begin{minipage}{\linewidth}
  \tikz[baseline=-0.7ex]\node[draw=popink,line width=1.6pt,fill=poporange,rounded corners=4pt,minimum size=22pt,inner sep=0]{\popdisplay\fontsize{12}{12}\selectfont\color{popcream}\thepopsec};%
  \hspace{8pt}{\popdisplay\fontsize{15}{17}\selectfont\color{popink}\MakeUppercase{#1}}\par
  \vspace{4pt}\noindent{\color{poporange}\rule{36pt}{3.6pt}}
  \end{minipage}\par\nopagebreak\vspace{8pt}
}
\newcounter{popsub}
\newcommand{\courssub}[1]{%
  \stepcounter{popsub}%
  \par\vspace{9pt}\noindent{\popxbold\fontsize{11.5}{13}\selectfont\color{popink}{\color{poporange}\thepopsec.\thepopsub}\hspace{6pt}#1}\par\nopagebreak\vspace{4pt}
}

% ============================================================
% Boîtes pédagogiques — même recette : fond blanc, bordure épaisse colorée,
% ombre dure encre, badge plein. Titre optionnel [#1].
% ============================================================
\tcbset{popbase/.style={breakable,enhanced,colback=white,boxrule=2.3pt,arc=9pt,
  shadow={3pt}{-3pt}{0pt}{popink},left=14pt,right=14pt,top=11pt,bottom=11pt,before skip=11pt,after skip=15pt,
  fontupper=\popsemi\fontsize{10.6}{14.5}\selectfont\color{popink},
  fontlower=\popsemi\fontsize{10.6}{14.5}\selectfont\color{popink}}}
% Coche dessinée (le glyphe ✓ n'existe pas dans les polices du document)
\newcommand{\popcheck}[1]{\tikz[baseline=-0.5ex]\draw[#1,line width=1.6pt,line cap=round,line join=round](-0.13,0.0)--(-0.04,-0.09)--(0.13,0.1);}
\providecommand{\coche}{\popcheck{popgreen}}
\providecommand{\resultat}[1]{\par\smallskip\noindent\fcolorbox{popgreen}{popgreenL}{\parbox{\dimexpr\linewidth-2\fboxsep-2\fboxrule\relax}{\textbf{Résultat.} #1}}\par\smallskip}
\newcommand{\popboxhead}[3]{\poppill{#1}{#2}{white}\ifx\relax#3\relax\else\hspace{6pt}{\popxbold\fontsize{10.6}{12}\selectfont\color{popink}#3}\fi\par\vspace{7pt}}

\newtcolorbox{aretenir}[1][]{popbase,colframe=popgreen,before upper={\popboxhead{À retenir}{popgreen}{#1}}}
% Texte en allemand (document de lecture, dialogue, extrait) : cadre
% d'étude. Un titre optionnel (source, auteur) s'affiche dans l'en-tête.
\newtcolorbox{textebox}[1][]{popbase,colframe=popblue,before upper={\popboxhead{Text}{popblue}{#1}\begin{otherlanguage}{ngerman}},after upper={\end{otherlanguage}}}
% Marques ✓ / ✗ (forme correcte / fautive) souvent utilisées par les modèles
\providecommand{\cross}{\textcolor{poppink}{\ensuremath{\times}}}
\providecommand{\xmark}{\cross}\providecommand{\crossmark}{\cross}\providecommand{\cmark}{\ensuremath{\checkmark}}
% Ligne de réponse / justification à compléter (inventée par les modèles)
\providecommand{\justification}[1]{\par\noindent\textit{Justification~:} \dotfill\par}
% \textipa (package tipa non chargé) : la police ne couvre pas l'API -> simple passage
\providecommand{\textipa}[1]{#1}
% Soulignement (ulem non chargé) : \uline, \ul
\providecommand{\uline}[1]{\underline{#1}}\providecommand{\ul}[1]{\underline{#1}}
% \glossaire{\item ...} inventée par les modèles : liste de glossaire
\providecommand{\glossaire}[1]{\begin{itemize}#1\end{itemize}}
% \alert (beamer) inventée par les modèles : texte mis en évidence
\providecommand{\alert}[1]{\textbf{\textcolor{poppink}{#1}}}
% variantes ASCII d'ordinaux écrites en macro
\providecommand{\iere}{\textsuperscript{re}}\providecommand{\ieme}{\textsuperscript{e}}\providecommand{\ere}{\textsuperscript{re}}\providecommand{\eme}{\textsuperscript{e}}\providecommand{\er}{\textsuperscript{er}}\providecommand{\ie}{\textsuperscript{e}}
% Ordinaux français écrits en macro par les modèles : 1\ière, 3\ième
\newcommand{\ière}{\textsuperscript{re}}\newcommand{\ième}{\textsuperscript{e}}
% \exemple (inventée par les modèles) : étiquette « Exemple : »
\providecommand{\exemple}{\textbf{Exemple~:}\ }
% \square utilisable aussi hors mode math (cases à cocher)
\AtBeginDocument{\let\squareMath\square\renewcommand{\square}{\ensuremath{\squareMath}}}
% Mot ou phrase en allemand à l'intérieur d'un paragraphe français
\newcommand{\all}[1]{\ifmmode\text{\textit{#1}}\else\begingroup\selectlanguage{ngerman}\itshape #1\endgroup\fi}
\let\esp\all % alias toléré
\newtcolorbox{exemplebox}[1][]{popbase,colframe=poporange,before upper={\popboxhead{Exemple}{poporange}{#1}}}
\newtcolorbox{definition}[1][]{popbase,colframe=popblue,before upper={\popboxhead{Définition}{popblue}{#1}}}
\newtcolorbox{propriete}[1][]{popbase,colframe=poppurple,before upper={\popboxhead{Propriété}{poppurple}{#1}}}
\newtcolorbox{methode}[1][]{popbase,colframe=popgold,before upper={\popboxhead{Méthode}{popgold}{#1}}}
\newtcolorbox{attention}[1][]{popbase,colframe=poppink,before upper={\popboxhead{Attention}{poppink}{#1}}}
\newtcolorbox{experience}[1][]{popbase,colframe=popblue,colback=popblueL,before upper={\popboxhead{Expérience}{popblue}{#1}}}
% « Le savais-tu ? » : carte violette pleine (contexte, culture, Cameroun)
\newtcolorbox{savaistu}[1][]{popbase,colback=poppurple,colframe=popink,
  fontupper=\popsemi\fontsize{10.4}{14.2}\selectfont\color{white},
  before upper={\poppill{Le savais-tu\,?}{white}{poppurple}\ifx\relax#1\relax\else\hspace{6pt}{\popxbold\color{white}#1}\fi\par\vspace{7pt}}}
% Exercice résolu : énoncé en haut, \tcblower puis la solution sur fond crème
\newcounter{popexo}\newcounter{popexr}
\newtcolorbox{exoresolu}[1][]{popbase,colframe=poporange,colbacklower=popcream,
  segmentation style={popink,line width=1.2pt,dashed},
  before upper={\stepcounter{popexr}\popboxhead{Exercice résolu \thepopexr}{poporange}{#1}},
  before lower={\poppill{Solution}{popink}{popcream}\par\vspace{6pt}}}
% Exercice (non résolu en place) — la correction est dans la partie Corrigés
\newtcolorbox{exercice}[1][]{popbase,colframe=popink,
  before upper={\stepcounter{popexo}\popboxhead{Exercice \thepopexo}{popink}{#1}}}
% Corrigé
\newtcolorbox{corrige}[1][]{popbase,colframe=popgreen,colback=popgreenL,
  before upper={\popboxhead{Corrigé}{popgreen}{#1}}}
% Objectifs / prérequis / bilan
\newtcolorbox{objectifs}{popbase,colframe=popink,colback=poporangeL,
  before upper={\popboxhead{Objectifs de la leçon}{popink}{}}}
\newtcolorbox{prerequis}{popbase,colframe=popink,colback=popblueL,
  before upper={\popboxhead{Prérequis}{popblue}{}}}
\newtcolorbox{bilan}{popbase,colframe=popink,colback=white,boxrule=2.8pt,
  before upper={\popboxhead{Fiche bilan}{poporange}{L'essentiel à connaître pour l'examen}}}
% Figure encadrée avec légende
\newtcolorbox{popfigure}[1][]{enhanced,breakable=false,colback=white,colframe=popline,boxrule=1.4pt,arc=8pt,
  left=10pt,right=10pt,top=10pt,bottom=8pt,before skip=10pt,after skip=12pt,halign=center,
  lower separated=false,halign lower=center,
  fontlower=\popsemi\fontsize{9}{11.5}\selectfont\color{popmuted},
  #1}
\newcommand{\legende}[1]{\par\vspace{6pt}{\popsemi\fontsize{9}{11.5}\selectfont\color{popmuted}#1}}

% ============================================================
% Signature OnBuch+
% ============================================================
\newcommand{\poplogoinline}[1]{{\poplogo\fontsize{#1}{#1}\selectfont\color{popink}OnBuch\color{poporange}+}}
\newcommand{\signatureonbuch}{%
  \begin{tcolorbox}[enhanced,colback=popink,colframe=popink,boxrule=2.2pt,arc=10pt,
      drop shadow=poporange,left=14pt,right=14pt,top=10pt,bottom=10pt,before skip=0pt,after skip=0pt]
    \tikz[baseline=-0.7ex]\node[circle,fill=popgreen,draw=popcream,line width=1.3pt,minimum size=22pt,inner sep=0]{\popcheck{white}};%
    \hspace{9pt}\begin{minipage}[c]{0.62\linewidth}
      {\popxbold\fontsize{12}{13}\selectfont\color{popcream}Signé\ \textbullet\ {\poplogo OnBuch\color{poporange}+}}\par\vspace{2pt}
      {\raggedright\popsemi\fontsize{8}{10.5}\selectfont\color{popline}\MakeUppercase{Équipe pédagogique OnBuch+}\\\MakeUppercase{Conforme au programme officiel MINESEC}\par}
    \end{minipage}\hfill
    {\poplogo\fontsize{22}{22}\selectfont\color{popcream}OnBuch\color{poporange}+}
  \end{tcolorbox}%
}

% ============================================================
% Page de garde
% #1 titre · #2 matière · #3 niveau · #4 module · #5 n° leçon · #6 durée ·
% #7 description / objectifs (texte ou itemize)
% ============================================================
\newcommand{\pagegarde}[7]{%
  \thispagestyle{empty}
  \newgeometry{a4paper,margin=1.7cm,top=1.6cm,bottom=1.4cm}%
  \begin{tikzpicture}[remember picture,overlay]
    \fill[poporange!18] ([xshift=-3.2cm,yshift=-3.6cm]current page.north east) circle (4.6cm);
    \fill[poppurple!14] ([xshift=2.4cm,yshift=7.6cm]current page.south west) circle (3.8cm);
  \end{tikzpicture}%
  {\poplogo\fontsize{30}{30}\selectfont\color{popink}OnBuch\color{poporange}+}\hfill
  \raisebox{0.6ex}{\poppill{Cours signé \textbullet\ Premium}{popink}{popcream}}\par
  \vspace{1pt}\popsquiggle{poppurple}\par
  \vspace{8pt}
  \begin{tcolorbox}[enhanced,colback=poporange,colframe=popink,boxrule=2.8pt,arc=13pt,
      drop shadow=popink,left=18pt,right=18pt,top=12pt,bottom=13pt,after skip=10pt]
    \poppill{#2 \textbullet\ #3}{popink}{popcream}\hspace{5pt}\poppill{#4}{white}{popink}\par\vspace{8pt}
    {\popdisplay\fontsize{14}{14}\selectfont\color{popink}LEÇON #5}\par\vspace{4pt}
    {\raggedright\hyphenpenalty=10000\exhyphenpenalty=10000\popdisplay\fontsize{25}{27}\selectfont\color{popcream}\MakeUppercase{#1}\par}
    \vspace{8pt}
    {\popxbold\fontsize{9.5}{9.5}\selectfont\color{popink}\MakeUppercase{Durée conseillée : #6}}
  \end{tcolorbox}
  \begin{tcolorbox}[enhanced,colback=white,colframe=popink,boxrule=2.2pt,arc=10pt,
      drop shadow=popink,left=16pt,right=16pt,top=10pt,bottom=10pt,after skip=8pt]
    \poppill{Dans cette leçon}{poppurple}{white}\par\vspace{5pt}
    \popsemi\fontsize{9.3}{12.6}\selectfont\color{popink}#7
  \end{tcolorbox}
  \noindent\begin{minipage}[t]{0.487\linewidth}
    \begin{tcolorbox}[enhanced,colback=popgreen,colframe=popink,boxrule=2.2pt,arc=10pt,
        drop shadow=popink,left=11pt,right=11pt,top=10pt,bottom=10pt,height=3.1cm,valign=center]
      \tcbox[colback=white,colframe=popink,boxrule=1.3pt,arc=5pt,boxsep=0pt,nobeforeafter,
        left=3pt,right=3pt,top=3pt,bottom=3pt,valign=center]{\qrcode[height=1.9cm]{https://play.google.com/store/apps/details?id=com.luvvix.onbuch&hl=fr}}%
      \hspace{8pt}\begin{minipage}[c]{0.5\linewidth}
        {\popdisplay\fontsize{11}{12.5}\selectfont\color{white}\MakeUppercase{Télécharge OnBuch}}\par\vspace{4pt}
        {\raggedright\popsemi\fontsize{8.4}{11}\selectfont\color{white}Gratuit \textbullet\ Google Play\\Tuteur IA, annales, concours, résultats\par}
      \end{minipage}
    \end{tcolorbox}
  \end{minipage}\hfill
  \begin{minipage}[t]{0.487\linewidth}
    \begin{tcolorbox}[enhanced,colback=poppurple,colframe=popink,boxrule=2.2pt,arc=10pt,
        drop shadow=popink,left=11pt,right=11pt,top=10pt,bottom=10pt,height=3.1cm,valign=center]
      \tikz[baseline=-0.7ex]\node[circle,fill=white,draw=popink,line width=1.3pt,minimum size=1.6cm,inner sep=0]{\popdisplay\fontsize{24}{24}\selectfont\color{poppurple}+};%
      \hspace{8pt}\begin{minipage}[c]{0.6\linewidth}
        {\popdisplay\fontsize{11}{12.5}\selectfont\color{white}\MakeUppercase{Passe à OnBuch+}}\par\vspace{4pt}
        {\raggedright\popsemi\fontsize{8.4}{11}\selectfont\color{white}Tous les cours signés, Tuteur IA illimité, fiches PDF — onglet OnBuch+ de l'app\par}
      \end{minipage}
    \end{tcolorbox}
  \end{minipage}
  \vspace{8pt}
  \popcontenu
  \vfill
  \signatureonbuch
  \restoregeometry
  \newpage
}

% Rangée « Ce cours contient » (page de garde)
\newcommand{\poptile}[3]{% #1 couleur #2 gros texte #3 légende
  \begin{tcolorbox}[enhanced,colback=#1,colframe=popink,boxrule=2pt,arc=9pt,shadow={2.5pt}{-2.5pt}{0pt}{popink},
      left=6pt,right=6pt,top=9pt,bottom=9pt,halign=center,valign=center,height=1.75cm,width=\linewidth,nobeforeafter]
    {\popdisplay\fontsize{12.5}{13}\selectfont\color{white}\MakeUppercase{#2}}\par\vspace{3pt}
    {\centering\popsemi\fontsize{7.8}{9.5}\selectfont\color{white}#3\par}
  \end{tcolorbox}}
\newcommand{\popcontenu}{%
  \par\noindent{\popxboldls\fontsize{7.5}{7.5}\selectfont\color{popmuted}CE COURS SIGNÉ CONTIENT}\par\vspace{7pt}
  \noindent\begin{minipage}[t]{0.235\linewidth}\poptile{popblue}{Cours}{complet, illustré, pas à pas}\end{minipage}\hfill
  \begin{minipage}[t]{0.235\linewidth}\poptile{poporange}{Méthodes}{et exercices résolus}\end{minipage}\hfill
  \begin{minipage}[t]{0.235\linewidth}\poptile{poppink}{Exercices}{type examen + corrigés}\end{minipage}\hfill
  \begin{minipage}[t]{0.235\linewidth}\poptile{popgreen}{Fiche bilan}{l'essentiel à retenir}\end{minipage}\par}

% Sommaire
\newtcolorbox{sommaire}{enhanced,colback=white,colframe=popink,boxrule=2.2pt,arc=10pt,
  drop shadow=popink,left=18pt,right=18pt,top=13pt,bottom=13pt,before skip=6pt,after skip=20pt,
  before upper={\poppill{Sommaire}{poppurple}{white}\par\vspace{9pt}\popsemi\fontsize{10.6}{14.5}\selectfont\color{popink}}}

% Signature de fin de cours — poussée tout en bas de la dernière page
\newcommand{\findecours}{%
  \par\vfill\nopagebreak
  \begin{center}\popsquiggle{poporange}\end{center}\vspace{4pt}
  \signatureonbuch
}
\AtBeginDocument{\def\llbracket{\mathopen{[\mkern-3.5mu[}}\def\rrbracket{\mathclose{]\mkern-3.5mu]}}} % intervalles d'entiers (glyphe ⟦ absent de la police)
% Glyphes absents de Fira Math : redéfinis à partir de glyphes disponibles
\AtBeginDocument{%
  \def\setminus{\mathbin{\backslash}}\let\smallsetminus\setminus
  \def\vdots{\mathord{\vbox{\baselineskip3.2pt\lineskiplimit0pt\kern4pt\hbox{.}\hbox{.}\hbox{.}}}}%
  \def\ddots{\mathinner{\mkern1mu\raise6.5pt\hbox{.}\mkern2mu\raise3.6pt\hbox{.}\mkern2mu\raise0.7pt\hbox{.}\mkern1mu}}%
  \def\triangle{\mathord{\scalebox{0.9}{$\symup{\Delta}$}}}%
  \def\nabla{\mathord{\raisebox{\depth}{\scalebox{1}[-1]{$\symup{\Delta}$}}}}%
  \def\checkmark{\popcheck{popdark}}%
  \def\star{\mathbin{\ast}}\def\bigstar{\mathord{\ast}}%
  \def\diamond{\mathbin{\scalebox{0.6}{$\Diamond$}}}%
  \def\bigcup{\mathop{\mathchoice{\raisebox{-0.3em}{\scalebox{1.7}{$\cup$}}}{\scalebox{1.25}{$\cup$}}{\cup}{\cup}}\displaylimits}%
  \def\bigcap{\mathop{\mathchoice{\raisebox{-0.3em}{\scalebox{1.7}{$\cap$}}}{\scalebox{1.25}{$\cap$}}{\cap}{\cap}}\displaylimits}%
  \def\lesseqqgtr{\mathrel{\lessgtr}}\def\bigoplus{\mathop{\oplus}\displaylimits}%
}
\providecommand{\ssi}{si et seulement si} % abréviation fréquente
\AtBeginDocument{\providecommand{\wideparen}{\overparen}} % arc de cercle

\begin{document}

\pagegarde{Le vandalisme et le terrorisme}{Allemand}{1ère A}{Chapitre 5 : Citoyenneté}{ALA25}{2 h}{Cette leçon propose la lecture et le commentaire d'un texte d'actualité sur le vandalisme et le terrorisme. Elle permet d'acquérir le vocabulaire thématique (der Terror, der Anschlag, die Gewalt, beschädigen, die Polizei, die Sicherheit) et de maîtriser le passif simple ainsi que l'expression de la cause et de la conséquence.
\vspace{4pt}
\begin{multicols}{2}\small
\begin{itemize}
\item À la fin de cette leçon, je sais comprendre globalement puis en détail un texte allemand sur le vandalisme et le terrorisme.
\item Je sais employer correctement le vocabulaire du thème avec les articles et les pluriels.
\item Je sais reconnaître et former le passif simple au présent et au prétérit.
\item Je sais exprimer la cause (weil, denn, wegen) et la conséquence (deshalb, so dass, darum).
\item Je sais répondre à des questions de compréhension en phrases complètes.
\item Je sais résumer un texte en quelques phrases au passif.
\end{itemize}\end{multicols}}

\begin{sommaire}
\begin{multicols}{2}
\begin{enumerate}
\item Texte support : lecture et première compréhension
\item Compréhension détaillée : questions et méthode
\item Lexique thématique : le vandalisme et le terrorisme
\item Grammaire 1 : le passif simple
\item Grammaire 2 : la cause et la conséquence
\item Méthode du commentaire et production guidée
\item Activité d'intégration
\item Méthodes et exercices résolus
\item Exercices
\item Corrigés des exercices
\item Fiche bilan
\end{enumerate}
\end{multicols}
\end{sommaire}

\begin{objectifs}
\begin{itemize}
\item À la fin de cette leçon, je sais comprendre globalement puis en détail un texte allemand sur le vandalisme et le terrorisme.
\item Je sais employer correctement le vocabulaire du thème avec les articles et les pluriels.
\item Je sais reconnaître et former le passif simple au présent et au prétérit.
\item Je sais exprimer la cause (weil, denn, wegen) et la conséquence (deshalb, so dass, darum).
\item Je sais répondre à des questions de compréhension en phrases complètes.
\item Je sais résumer un texte en quelques phrases au passif.
\item Je sais rédiger un court paragraphe argumenté sur un fait de violence urbaine.
\item Je sais éviter les pièges fréquents des francophones (passif avec 'sein', place de 'werden', 'wegen' + génitif/datif).
\end{itemize}
\end{objectifs}

\begin{prerequis}
\begin{itemize}
\item Le présent et le prétérit des verbes forts et faibles (Seconde)
\item Les auxiliaires haben/sein et les participes passés (Seconde)
\item Les cas : nominatif, accusatif, datif (Seconde)
\item Les connecteurs simples : und, aber, denn, weil (début d'année)
\item Le vocabulaire de la ville et de la vie quotidienne (Seconde)
\end{itemize}
\end{prerequis}

\newpage

\coursec{Texte support : lecture et première compréhension}

\courssub{Situation d'accroche et problématique}

Imagine la scène : en rentrant du lycée, ce soir, tu découvres que le banc public devant ton quartier a été \emph{cassé} pendant la nuit et que le mur de ton école est couvert de graffitis. Le lendemain, aux informations, on parle d'un \emph{attentat} dans une capitale européenne : des contrôles renforcés dans les gares, des caméras partout, des citoyens inquiets. Ces deux faits, en apparence très différents, appartiennent au même grand thème : \all{die Gewalt} (la violence).

Comment parler de ces faits en allemand ? Comment expliquer \emph{pourquoi} ils arrivent et \emph{quelles en sont les conséquences} ? C'est toute la problématique de cette leçon : te donner les mots et les structures pour \cle{comprendre} un texte d'actualité sur le vandalisme et le terrorisme, et pour en \cle{débattre} avec des arguments.

Dans cette première étape, nous allons lire un article de presse allemand, le comprendre globalement, puis apprendre à le reformuler avec nos propres mots — une compétence essentielle à l'examen.

\courssub{Première lecture : découverte du texte}

\begin{methode}[Comment aborder un texte de presse allemand ?]
Avant même de lire en détail, suis toujours ces étapes :
\begin{enumerate}
\item \textbf{Lecture silencieuse} (2 minutes) : lis le texte entier sans dictionnaire. Ne t'arrête pas sur les mots inconnus.
\item \textbf{Repérage des indices} : titre, longueur, mots en gras, noms propres. Demande-toi : de quel type de texte s'agit-il ?
\item \textbf{Lecture à voix haute} : prononce chaque phrase à voix haute. En allemand, presque tout se prononce comme cela s'écrit.
\item \textbf{Compréhension globale} : réponds mentalement aux cinq questions : \all{Wer? Was? Wo? Wann? Warum?} (Qui ? Quoi ? Où ? Quand ? Pourquoi ?)
\item \textbf{Lecture détaillée} : seulement maintenant, utilise le glossaire pour les mots nouveaux.
\end{enumerate}
\end{methode}

Lis maintenant le texte suivant, d'abord en silence, puis à voix haute.

\begin{textebox}[Gewalt in der Stadt — article de presse]
\textbf{Gewalt in der Stadt}

In vielen deutschen Städten werden nachts Bushaltestellen beschädigt und Wände mit Graffiti beschmiert. Die Polizei spricht von Vandalismus. „Jede Woche müssen zerbrochene Scheiben und zerstörte Sitzbänke ersetzt werden“, erklärt ein Polizist aus München. „Das kostet die Stadt viel Geld.“

Viel schlimmer sind Terroranschläge: Dabei werden Menschen verletzt oder sogar getötet. Wegen dieser Anschläge wird die Sicherheit an Bahnhöfen und Flughäfen verstärkt. Taschen werden kontrolliert, Kameras werden installiert, und mehr Polizisten werden eingesetzt.

Viele Bürger haben Angst, deshalb fordern sie mehr Kameras und mehr Polizisten. „Ich fühle mich im Zug nicht mehr sicher“, sagt eine Studentin aus Berlin. „Aber ich weiß, dass die Polizei alles tut, um uns zu schützen.“

Experten sagen: Die Jugend muss über Gewalt aufgeklärt werden, damit solche Taten verhindert werden. In vielen Schulen werden deshalb Projekte gegen Gewalt organisiert. „Prävention ist besser als Strafe“, meint ein Lehrer. „Wenn junge Menschen verstehen, warum Gewalt falsch ist, werden weniger Straftaten begangen.“
\end{textebox}

\begin{definition}[Glossaire : les mots nouveaux du texte]
\begin{itemize}
\item \all{die Gewalt} (pas de pluriel courant) : la violence
\item \all{die Bushaltestelle, -n} : l'arrêt de bus
\item \all{beschädigen} : endommager, dégrader
\item \all{beschmieren} : salir, barbouiller
\item \all{die Scheibe, -n} : la vitre ; \all{zerbrochen} : cassé, brisé
\item \all{die Sitzbank, -bänke} : le banc (public)
\item \all{ersetzen} : remplacer
\item \all{der Terroranschlag, -anschläge} : l'attentat terroriste
\item \all{verletzen} : blesser ; \all{töten} : tuer
\item \all{die Sicherheit} : la sécurité ; \all{verstärken} : renforcer
\item \all{der Bahnhof, -höfe} : la gare ; \all{der Flughafen, -häfen} : l'aéroport
\item \all{der Bürger, -} : le citoyen ; \all{fordern} : exiger, réclamer
\item \all{sich sicher fühlen} : se sentir en sécurité
\item \all{schützen} : protéger
\item \all{aufklären} : sensibiliser, informer
\item \all{die Tat, -en} : l'acte, le fait ; \all{verhindern} : empêcher
\item \all{die Prävention} : la prévention ; \all{die Strafe, -n} : la punition, la peine
\item \all{die Straftat, -en} : le délit, le crime ; \all{begehen} (irrégulier) : commettre
\end{itemize}
\end{definition}

\begin{attention}[Prononciation : quelques pièges]
\begin{itemize}
\item \all{Gewalt} se prononce « gue-valt » (le \all{w} allemand = notre « v »).
\item \all{Sicherheit} se prononce « zi-heur-haït ».
\item \all{Bürger} : le \all{ü} se prononce comme le « u » français : « bur-gueur ».
\item \all{Flughafen} se prononce « flouk-ha-fen ».
\end{itemize}
\end{attention}

\courssub{Repérage : thème, type de texte, source, destinataire}

Avant de répondre aux questions, identifions ensemble les caractéristiques du texte. C'est la première chose à faire dans tout exercice de compréhension écrite.

\begin{center}
\begin{tabularx}{\linewidth}{|X|X|}\hline
\rowcolor{popblueL} \textbf{Critère} & \textbf{Réponse} \\ \hline
Thème (\all{das Thema}) & La violence en ville : vandalisme et terrorisme \\ \hline
Type de texte (\all{der Texttyp}) & Un article de presse (reportage d'actualité) \\ \hline
Source probable (\all{die Quelle}) & Un journal allemand (quotidien régional ou national) \\ \hline
Destinataire (\all{der Empfänger}) & Le grand public, les lecteurs d'un journal \\ \hline
Indices qui le prouvent & Titre informatif, citations de témoins (policier, étudiante, expert), ton objectif, faits récents \\ \hline
\end{tabularx}
\end{center}

\begin{aretenir}[Comment reconnaître un article de presse ?]
Un article de presse se reconnaît à : un \cle{titre} court et informatif, des \cle{faits} récents et vérifiables, des \cle{citations} entre guillemets („…“), un ton plutôt \cle{objectif}, et des paragraphes courts. Ici, les citations du policier, de l'étudiante et du professeur sont des indices typiques.
\end{aretenir}

\courssub{Compréhension globale : Wer? Was? Wo? Wann? Warum?}

Réponds maintenant aux cinq questions fondamentales. C'est la clé de la compréhension globale : si tu peux y répondre, tu as compris l'essentiel du texte.

\begin{itemize}
\item \all{Wer?} (Qui ?) — Des vandales inconnus, des terroristes, la police, les citoyens, les experts.
\item \all{Was?} (Quoi ?) — Des dégradations (arrêts de bus, graffitis) et des attentats ; des mesures de sécurité renforcées.
\item \all{Wo?} (Où ?) — Dans de nombreuses villes allemandes : gares, aéroports, rues.
\item \all{Wann?} (Quand ?) — La nuit pour le vandalisme ; aujourd'hui, à notre époque, pour les attentats.
\item \all{Warum?} (Pourquoi ?) — Le texte ne donne pas les motivations des auteurs, mais il explique les conséquences : peur des citoyens, sécurité renforcée, prévention à l'école.
\end{itemize}

\begin{exemplebox}[Premières questions de compréhension globale (oral)]
Entraîne-toi à répondre à voix haute, par une phrase complète :
\begin{enumerate}
\item \all{Worum geht es in diesem Text?} — \all{Es geht um Gewalt in der Stadt: Vandalismus und Terrorismus.} « De quoi parle ce texte ? — Il parle de la violence en ville : le vandalisme et le terrorisme. »
\item \all{Was wird nachts beschädigt?} — \all{Bushaltestellen und Wände werden beschädigt.} « Qu'est-ce qui est endommagé la nuit ? — Les arrêts de bus et les murs. »
\item \all{Was fordern viele Bürger?} — \all{Sie fordern mehr Kameras und mehr Polizisten.} « Que réclament beaucoup de citoyens ? — Plus de caméras et plus de policiers. »
\end{enumerate}
\end{exemplebox}

\courssub{Deux notions clés : Vandalismus et Terrorismus}

Le texte oppose deux formes de violence. Il est essentiel de bien les distinguer.

\begin{definition}[der Vandalismus]
Le \cle{vandalisme} (\all{der Vandalismus}, sans pluriel) est la \emph{destruction volontaire de biens publics ou privés} : \all{die absichtliche Zerstörung fremden Eigentums}. Exemples : casser un banc, briser une vitre, couvrir un mur de graffitis. Il s'agit de dégâts matériels, sans intention de tuer.
\end{definition}

\begin{definition}[der Terrorismus]
Le \cle{terrorisme} (\all{der Terrorismus}, sans pluriel) est \emph{l'usage de la violence pour des motifs politiques ou idéologiques afin de semer la peur} (\all{der Terror}). Contrairement au vandalisme, il vise des personnes : \all{Dabei werden Menschen verletzt oder sogar getötet.} « Des personnes sont blessées ou même tuées. »
\end{definition}

\begin{center}
\begin{tabularx}{\linewidth}{|X|X|X|}\hline
\rowcolor{popblueL} \textbf{Critère} & \textbf{Vandalismus} & \textbf{Terrorismus} \\ \hline
Cible & Les biens (murs, bancs, vitres) & Les personnes \\ \hline
But & Dégrader, marquer son territoire & Semer la peur (der Terror) \\ \hline
Motivation & Souvent sans motif politique & Politique ou idéologique \\ \hline
Exemple du texte & Graffitis, arrêts de bus cassés & Attentats dans les gares et aéroports \\ \hline
Conséquence & Coût pour la ville & Sécurité renforcée, peur des citoyens \\ \hline
\end{tabularx}
\end{center}

\begin{savaistu}[Un mot, deux sens : der Anschlag]
En Allemagne, le mot \all{der Anschlag} désigne à la fois un \emph{attentat} et une \emph{affiche placardée} au mur ! C'est le contexte qui décide du sens : \all{ein Terroranschlag} = un attentat terroriste ; \all{ein Anschlag an der Wand} = une affiche au mur. D'ailleurs, le panneau d'affichage s'appelle \all{der Anschlagkasten}. Attention donc à toujours lire la phrase entière avant de traduire !
\end{savaistu}

\courssub{Carte mentale du vocabulaire}

Pour mémoriser le vocabulaire du texte, rien de tel qu'une carte mentale. Observe comment les mots s'organisent autour du mot central \all{Gewalt}.

\begin{popfigure}
\begin{tikzpicture}[mindmap, grow cyclic, every node/.style={concept, align=center, font=\small},
root concept/.append style={concept color=popblue, font=\bfseries},
level 1/.append style={level distance=3.2cm, sibling angle=90},
level 2/.append style={level distance=2.2cm, sibling angle=50, font=\footnotesize}]
\node[root concept]{Gewalt}
  child[concept color=poporange]{ node {Vandalismus}
    child { node {beschädigen} }
    child { node {Graffiti} }
  }
  child[concept color=poppink]{ node {Terrorismus}
    child { node {der Anschlag} }
    child { node {der Terror} }
  }
  child[concept color=popgreen]{ node {Polizei}
    child { node {die Sicherheit} }
    child { node {Kameras} }
  }
  child[concept color=poppurple]{ node {Folgen}
    child { node {Angst} }
    child { node {Prävention} }
  };
\end{tikzpicture}
\legende{Carte mentale : le champ lexical de la violence (die Gewalt)}
\end{popfigure}

\begin{methode}[Construire ta propre carte mentale]
\begin{enumerate}
\item Écris le mot-clé au centre (ici : \all{Gewalt}).
\item Trace une branche par grande idée du texte (ici : 4 branches).
\item Ajoute sur chaque branche les mots du texte qui s'y rapportent, avec leur \cle{article} et leur \cle{pluriel}.
\item Reproduis cette carte dans ton cahier et complète-la au fur et à mesure de la leçon.
\end{enumerate}
\end{methode}

\courssub{Reformulation : dire le texte avec ses propres mots}

Dernière étape de la compréhension globale : résumer le texte en \emph{une seule phrase}, avec tes propres mots. C'est un exercice classique à l'examen.

\begin{exoresolu}[Reformuler le contenu du texte en une phrase]
Énoncé : \all{Fassen Sie den Text in einem Satz zusammen.} « Résumez le texte en une phrase. »
\tcblower
Solution détaillée : on reprend les cinq réponses (Wer? Was? Wo? Wann? Warum?) et on les condense.

\all{In vielen deutschen Städten gibt es Vandalismus und Terroranschläge, deshalb wird die Sicherheit verstärkt und die Jugend wird über Gewalt aufgeklärt.}

« Dans de nombreuses villes allemandes, il y a du vandalisme et des attentats terroristes ; c'est pourquoi la sécurité est renforcée et la jeunesse est sensibilisée à la violence. »

Méthode : sujet (\all{es gibt...}) + les deux problèmes (\all{Vandalismus und Terroranschläge}) + les conséquences introduites par \all{deshalb} et \all{und}. Une seule phrase, toutes les idées essentielles.
\end{exoresolu}

\begin{attention}[Pièges fréquents des francophones]
\begin{itemize}
\item Ne traduis pas « la violence » par \all{die Violenz} : ce mot n'existe pas ! On dit \all{die Gewalt}.
\item \all{beschädigen} signifie « endommager », pas « blesser » : pour les personnes, on utilise \all{verletzen}.
\item Attention au genre : on dit \all{der Terrorismus}, \all{der Vandalismus} (masculins, comme tous les mots en \all{-ismus}), mais \all{die Gewalt}, \all{die Polizei}, \all{die Sicherheit} (féminins).
\item \all{die Polizei} est toujours au singulier en allemand : \all{Die Polizei spricht von Vandalismus.} « La police parle de vandalisme. »
\end{itemize}
\end{attention}

\begin{exercice}[À toi de jouer !]
\begin{enumerate}
\item Réponds oralement : \all{Was ist schlimmer: Vandalismus oder Terrorismus? Warum?}
\item Retrouve dans le texte trois mots de la famille de \all{sicher}.
\item Résume le texte en une phrase, sans regarder la correction de l'exercice résolu.
\item Recopie la carte mentale et ajoute deux mots du glossaire sur chaque branche.
\end{enumerate}
\end{exercice}

\begin{corrige}[Exercice : À toi de jouer !]
\begin{enumerate}
\item Réponse possible : \all{Terrorismus ist schlimmer, weil Menschen verletzt oder getötet werden.} « Le terrorisme est pire, parce que des personnes sont blessées ou tuées. »
\item \all{sicher} (sûr), \all{die Sicherheit} (la sécurité), \all{sich sicher fühlen} (se sentir en sécurité).
\item Exemple : \all{Der Text handelt von Gewalt in deutschen Städten und von den Folgen für die Bürger.} « Le texte traite de la violence dans les villes allemandes et de ses conséquences pour les citoyens. »
\item Exemples : Vandalismus : \all{die Sitzbank, beschmieren} ; Terrorismus : \all{töten, verletzen} ; Polizei : \all{schützen, kontrollieren} ; Folgen : \all{die Strafe, aufklären}.
\end{enumerate}
\end{corrige}

\begin{aretenir}[L'essentiel de cette première étape]
\begin{itemize}
\item Le texte \all{Gewalt in der Stadt} est un \cle{article de presse} destiné au grand public.
\item Il oppose deux formes de violence : le \cle{vandalisme} (dégâts matériels) et le \cle{terrorisme} (violence contre des personnes, pour semer la peur).
\item Pour comprendre globalement un texte : réponds toujours à \all{Wer? Was? Wo? Wann? Warum?}
\item Pour mémoriser le vocabulaire : construis une \cle{carte mentale} avec articles et pluriels.
\item Pour résumer : condense les cinq réponses en \cle{une phrase} avec tes propres mots.
\end{itemize}
\end{aretenir}

\coursec{Compréhension détaillée : questions et méthode}

Dans la section précédente, nous avons lu pour la première fois le texte \all{„Vandalismus in der Stadt“} et nous en avons dégagé le sens global : de jeunes gens dégradent des biens publics la nuit, et la police ainsi que les habitants cherchent des solutions. Cette première lecture, rapide et globale, ne suffit pas : à l'examen comme dans la vie, il faut savoir retrouver une information précise, prouver ce que l'on affirme et exprimer son avis. C'est l'objet de cette deuxième section : apprendre à répondre, de manière rigoureuse et méthodique, à des questions de compréhension détaillée portant sur notre texte.

\courssub{Le texte de référence (rappel)}

Pour travailler, relisons attentivement le texte support. C'est sur lui que porteront toutes les questions de cette section.

\begin{textebox}[Vandalismus in der Stadt]
In unserer Stadt gibt es ein großes Problem: Vandalismus. Fast jede Nacht werden Bushaltestellen beschädigt, Wände werden mit Farbe beschmiert und Blumen im Park werden zerstört. Die Täter sind meistens Jugendliche zwischen 14 und 18 Jahren. Warum machen sie das? Ein Sozialarbeiter erklärt: „Viele Jugendliche haben keine Arbeit und keine Perspektive. Sie sind frustriert, und deshalb zerstören sie die Sachen der anderen.“

Die Polizei kontrolliert jetzt öfter die Parks und die Schulen. Trotzdem ist die Situation schwierig, denn die Täter kommen immer nachts, wenn niemand sie sehen kann. Die Stadt hat deshalb Kameras installiert, aber viele Bürger sind dagegen. Sie sagen: „Wir wollen Sicherheit, aber wir wollen auch unsere Freiheit.“

Ein Lehrer hat eine andere Idee: „Die Jugendlichen brauchen Aktivitäten: Sport, Musik, Theater. Wenn sie am Wochenende etwas zu tun haben, haben sie keine Zeit für Vandalismus.“ Die Stadt will jetzt ein Jugendzentrum bauen. Viele Einwohner finden diese Idee gut, weil sie das Problem an der Wurzel löst. Aber das Zentrum kostet viel Geld, und das Projekt dauert mindestens zwei Jahre.
\end{textebox}

\textbf{Glossaire :} die Bushaltestelle, -n (l'arrêt de bus) ; beschädigen (endommager) ; beschmieren (barbouiller, salir) ; zerstören (détruire) ; der Täter, - (l'auteur, le coupable) ; der Sozialarbeiter, - (le travailleur social) ; die Perspektive, -n (la perspective, l'avenir) ; frustriert (frustré) ; die Sicherheit (la sécurité) ; die Freiheit (la liberté) ; der Einwohner, - (l'habitant) ; die Wurzel, -n (la racine) ; das Jugendzentrum, die Jugendzentren (le centre de jeunes).

\courssub{Les questions de compréhension : les W-Fragen}

\begin{definition}[Les W-Fragen (questions ouvertes)]
En allemand, les questions de compréhension commencent presque toujours par un mot interrogatif qui commence par la lettre \textbf{W} : on les appelle les \cle{W-Fragen}. Chaque mot interrogatif cherche un type d'information précis.
\end{definition}

\begin{center}
\begin{tabularx}{\linewidth}{|l|X|X|}
\hline
\rowcolor{popblueL} \textbf{Mot interrogatif} & \textbf{Français} & \textbf{Information cherchée} \\
\hline
Wer? & Qui ? & une personne, un groupe de personnes \\
\hline
Was? & Quoi ? Que ? & une chose, une action \\
\hline
Wo? & Où ? & un lieu \\
\hline
Wann? & Quand ? & un moment, une date \\
\hline
Warum? & Pourquoi ? & une cause, une raison \\
\hline
Wie? & Comment ? & une manière, un moyen \\
\hline
Wie oft? & Combien de fois ? & une fréquence \\
\hline
Wie viel(e)? & Combien (de) ? & une quantité \\
\hline
\end{tabularx}
\end{center}

\begin{exemplebox}[Observer avant de comprendre]
Observons une question et sa réponse, tirées de notre texte :

\all{Was wird nachts beschädigt?} — « Qu'est-ce qui est endommagé la nuit ? »

Réponse : \all{Nachts werden Bushaltestellen beschädigt.} — « La nuit, les arrêts de bus sont endommagés. »

Remarquons deux choses : d'abord, la réponse reprend les mots de la question (\all{nachts}, \all{beschädigt}) ; ensuite, elle forme une \textbf{phrase complète}, avec un sujet, un verbe conjugué et un point final. C'est exactement ce qu'attend l'examinateur.
\end{exemplebox}

\courssub{La méthode en quatre étapes}

\begin{methode}[Répondre à une question de compréhension]
Pour répondre correctement à une question de compréhension, suis toujours ces quatre étapes, dans l'ordre :
\begin{enumerate}
\item \textbf{Lire la question} attentivement et identifier le mot interrogatif (\all{Wer? Was? Wo? Warum? Wie?}) : il annonce le type d'information à chercher.
\item \textbf{Repérer les mots-clés} de la question (le verbe, le nom principal) : ce sont eux que l'on va chercher dans le texte.
\item \textbf{Retrouver et souligner la phrase du texte} qui contient la réponse. Ne jamais répondre de mémoire : la preuve est toujours dans le texte.
\item \textbf{Rédiger une phrase complète} en reprenant les mots de la question, sans recopier tout le paragraphe : une seule phrase claire suffit.
\end{enumerate}
\end{methode}

\begin{popfigure}
\begin{tikzpicture}[
  node distance=0.55cm and 0.9cm,
  box/.style={draw=popblue, fill=popblueL, rounded corners=3pt, align=center, minimum height=0.9cm, text width=3.1cm, font=\small},
  arr/.style={-{Stealth[length=2.5mm]}, thick, popdark}
]
\node[box, align=center] (q){\textbf{1. Question}\\ \all{Was wird nachts beschädigt?}};
\node[box, right=of q, align=center] (r){\textbf{2. Repérage}\\ souligner la phrase dans le texte};
\node[box, right=of r, align=center] (f){\textbf{3. Reformulation}\\ reprendre les mots de la question};
\node[box, right=of f, align=center] (p){\textbf{4. Phrase complète}\\ \all{Nachts werden Bushaltestellen beschädigt.}};
\draw[arr] (q) -- (r);
\draw[arr] (r) -- (f);
\draw[arr] (f) -- (p);
\end{tikzpicture}
\legende{De la question à la réponse : le parcours en quatre étapes.}
\end{popfigure}

\begin{exoresolu}[Appliquer la méthode à trois questions]
Énoncé : réponds aux questions suivantes par des phrases complètes, en allemand.

a) \all{Wer beschädigt die Bushaltestellen?}

b) \all{Warum zerstören die Jugendlichen die Sachen der anderen?}

c) \all{Was will die Stadt bauen?}
\tcblower
Solution détaillée :

a) Mot interrogatif : \all{Wer?} → on cherche une personne. Mots-clés : \all{beschädigt}, \all{Bushaltestellen}. Dans le texte : \all{„Die Täter sind meistens Jugendliche zwischen 14 und 18 Jahren.“} Réponse : \all{Jugendliche zwischen 14 und 18 Jahren beschädigen die Bushaltestellen.} (Des jeunes de 14 à 18 ans endommagent les arrêts de bus.)

b) Mot interrogatif : \all{Warum?} → on cherche une cause. Dans le texte, le Sozialarbeiter explique : \all{„Sie sind frustriert, und deshalb zerstören sie die Sachen der anderen.“} Réponse : \all{Die Jugendlichen zerstören die Sachen der anderen, weil sie frustriert sind und keine Perspektive haben.} (Les jeunes détruisent les biens des autres parce qu'ils sont frustrés et n'ont pas d'avenir.)

c) Mot interrogatif : \all{Was?} → on cherche une chose. Dans le texte : \all{„Die Stadt will jetzt ein Jugendzentrum bauen.“} Réponse : \all{Die Stadt will ein Jugendzentrum bauen.} (La ville veut construire un centre de jeunes.)
\end{exoresolu}

\begin{attention}[Les erreurs qui coûtent des points]
\begin{itemize}
\item \textbf{Ne recopie pas tout le paragraphe.} Si la question porte sur une information précise, une phrase recopiée de dix lignes montre que tu n'as pas compris. Sélectionne uniquement l'information demandée.
\item \textbf{Ne réponds pas hors texte.} Même si tu connais le sujet, ta réponse doit venir du texte, pas de tes connaissances personnelles.
\item \textbf{Respecte l'ordre des mots allemand} dans ta réponse : le verbe conjugué en deuxième position. Mauvais : \all{Nachts Bushaltestellen werden beschädigt.} Correct : \all{Nachts werden Bushaltestellen beschädigt.}
\item \textbf{Accorde le verbe avec le nouveau sujet.} Si la question est au singulier et la réponse au pluriel, adapte la conjugaison.
\end{itemize}
\end{attention}

\courssub{Les questions vrai ou faux : justifier par une citation}

Un deuxième type de question très fréquent est le \cle{vrai/faux} (\all{Richtig oder falsch?}). La règle d'or est absolue : il ne suffit jamais de répondre \all{richtig} ou \all{falsch}. Il faut \textbf{justifier} sa réponse par une phrase tirée du texte, que l'on recopie entre guillemets.

\begin{propriete}[La justification par la citation]
Pour justifier une réponse vrai/faux, on cite la phrase du texte qui prouve la réponse. On utilise les guillemets allemands „ … “ et on peut introduire la citation par \all{Im Text steht:} (« Dans le texte, il est dit : ») ou \all{Der Text sagt:} (« Le texte dit : »).

Modèle de réponse : \all{Falsch. Im Text steht: „Die Täter kommen immer nachts.“} — « Faux. Dans le texte, il est dit : „Les auteurs viennent toujours la nuit.“ »
\end{propriete}

\begin{exemplebox}[Vrai ou faux : deux exemples corrigés]
a) \all{Die Täter kommen am Tag.} (Richtig oder falsch?)

Réponse : \all{Falsch. Im Text steht: „Die Täter kommen immer nachts, wenn niemand sie sehen kann.“} — « Faux. Dans le texte, il est dit : „Les auteurs viennent toujours la nuit, quand personne ne peut les voir.“ »

b) \all{Alle Bürger sind für die Kameras.} (Richtig oder falsch?)

Réponse : \all{Falsch. Der Text sagt: „Viele Bürger sind dagegen.“} — « Faux. Le texte dit : „Beaucoup de citoyens sont contre.“ »
\end{exemplebox}

\begin{attention}[Ne jamais répondre par un simple « ja » ou « nein »]
Au baccalauréat comme à l'examen de fin d'année, une réponse par \all{ja}, \all{nein}, \all{richtig} ou \all{falsch} sans justification ne rapporte \textbf{aucun point}, même si elle est correcte. L'examinateur veut voir la \textbf{preuve} : une phrase complète ou une citation tirée du texte. Retiens la formule : \emph{réponse + justification = point complet}.
\end{attention}

\courssub{La question d'opinion : « Was denkst du? »}

Le troisième type de question ne cherche plus une information dans le texte, mais \textbf{ton avis personnel}. On la reconnaît à des formules comme : \all{Was denkst du?} (« Qu'en penses-tu ? »), \all{Was meinst du?} (« Quel est ton avis ? »), \all{Findest du die Idee gut?} (« Trouves-tu l'idée bonne ? »).

\begin{propriete}[Exprimer son opinion avec « ich meine, dass... »]
Pour donner son avis en allemand, on utilise des verbes d'opinion suivis de \all{dass} + subordonnée (verbe conjugué à la fin) :

\all{Ich meine, dass das Jugendzentrum eine gute Idee ist.} — « Je pense que le centre de jeunes est une bonne idée. »

\all{Ich finde, dass Kameras wichtig für die Sicherheit sind.} — « Je trouve que les caméras sont importantes pour la sécurité. »

\all{Ich glaube, dass Sport den Jugendlichen hilft.} — « Je crois que le sport aide les jeunes. »

Autres formules utiles : \all{Meiner Meinung nach...} (« À mon avis... », verbe en deuxième position), \all{Ich bin der Meinung, dass...} (« Je suis de l'avis que... »).
\end{propriete}

\begin{center}
\begin{tabularx}{\linewidth}{|X|X|X|}
\hline
\rowcolor{popblueL} \textbf{Français} & \textbf{Deutsch} & \textbf{Remarque} \\
\hline
Je pense que... & Ich meine, dass... & verbe à la fin après \all{dass} \\
\hline
Je trouve que... & Ich finde, dass... & verbe à la fin \\
\hline
Je crois que... & Ich glaube, dass... & verbe à la fin \\
\hline
À mon avis,... & Meiner Meinung nach... & verbe en 2\textsuperscript{e} position \\
\hline
Je suis d'accord. & Ich bin einverstanden. & faux ami : \all{agréer} n'existe pas \\
\hline
Je ne suis pas d'accord. & Ich bin nicht einverstanden. & négation avec \all{nicht} \\
\hline
\end{tabularx}
\end{center}

\begin{attention}[Le verbe à la fin après « dass »]
L'erreur la plus fréquente des francophones est de garder l'ordre français après \all{dass}. En français : « Je pense que le centre est une bonne idée » (sujet + verbe). En allemand, le verbe conjugué va \textbf{à la fin} de la subordonnée :

Mauvais : \all{Ich meine, dass das Zentrum ist eine gute Idee.}

Correct : \all{Ich meine, dass das Zentrum eine gute Idee \textbf{ist}.}
\end{attention}

\courssub{Le résumé guidé en 3 à 4 phrases}

Dernière tâche de compréhension : résumer le texte (\all{der Text zusammenfassen}). Un bon résumé au niveau Première contient 3 à 4 phrases et suit la logique du texte : le problème, les causes, les solutions.

\begin{methode}[Rédiger un résumé guidé]
\begin{enumerate}
\item Relis le texte et repère \textbf{trois idées principales} : ici, le problème (le vandalisme), la cause (la frustration des jeunes), les solutions (caméras, centre de jeunes).
\item Écris \textbf{une phrase par idée}, avec tes propres mots si possible, sans détails secondaires (pas de chiffres, pas de citations).
\item Relie les phrases avec des connecteurs simples : \all{denn} (car), \all{deshalb} (c'est pourquoi), \all{aber} (mais).
\item Vérifie : verbe en deuxième position, majuscules aux noms, point final.
\end{enumerate}
\end{methode}

\begin{exemplebox}[Modèle de résumé]
\all{In der Stadt gibt es viel Vandalismus: Jugendliche beschädigen nachts Bushaltestellen und Wände. Sie machen das, weil sie frustriert sind und keine Perspektive haben. Die Stadt installiert deshalb Kameras, aber viele Bürger sind dagegen. Ein Jugendzentrum mit Sport und Musik könnte das Problem an der Wurzel lösen.}

« Dans la ville, il y a beaucoup de vandalisme : des jeunes endommagent la nuit des arrêts de bus et des murs. Ils font cela parce qu'ils sont frustrés et n'ont pas d'avenir. La ville installe donc des caméras, mais beaucoup de citoyens sont contre. Un centre de jeunes avec du sport et de la musique pourrait résoudre le problème à la racine. »
\end{exemplebox}

\courssub{Un point de grammaire utile pour les réponses : l'agent au passif}

Beaucoup de questions sur ce texte portent sur des phrases au passif (\all{werden beschädigt}, \all{werden beschmiert}). Pour répondre à la question \all{Von wem?} (« Par qui ? »), il faut connaître l'expression de l'agent.

\begin{propriete}[Le complément d'agent au passif : « von » + datif]
Au passif allemand, le complément d'agent (en français : « par... ») se rend par la préposition \all{von} suivie du \textbf{datif} :

\all{Die Wände werden \textbf{von Jugendlichen} beschmiert.} — « Les murs sont barbouillés \textbf{par des jeunes}. »

\all{Die Parks werden \textbf{von der Polizei} kontrolliert.} — « Les parcs sont contrôlés \textbf{par la police}. »

Attention au datif après \all{von} : \all{von der Polizei} (féminin), \all{von dem / vom Täter} (masculin), \all{von Jugendlichen} (pluriel avec -\textbf{n}).
\end{propriete}

\begin{exoresolu}[Question avec complément d'agent]
Énoncé : réponds par une phrase complète.

\all{Von wem werden die Parks kontrolliert?} — « Par qui les parcs sont-ils contrôlés ? »
\tcblower
Solution : on repère dans le texte \all{„Die Polizei kontrolliert jetzt öfter die Parks und die Schulen.“} L'agent est donc \all{die Polizei}. On construit la réponse au passif avec \all{von} + datif :

\all{Die Parks werden von der Polizei kontrolliert.} — « Les parcs sont contrôlés par la police. »

Piège évité : on n'écrit pas \all{von die Polizei} : après \all{von}, l'article féminin passe au datif \all{der}.
\end{exoresolu}

\courssub{Entraînement}

\begin{exercice}[À ton tour]
Réponds aux questions suivantes par des phrases complètes, en t'appuyant sur le texte.

a) \all{Wo werden Blumen zerstört?}

b) \all{Viele Bürger sind gegen die Kameras.} (Richtig oder falsch? Justifie ta réponse par une citation.)

c) \all{Was denkst du: Ist ein Jugendzentrum eine gute Lösung?} (Réponds avec \all{Ich meine, dass...}.)

d) Résume le texte en 3 phrases.
\end{exercice}

\begin{corrige}[Exercice « À ton tour »]
a) \all{Die Blumen werden im Park zerstört.} (« Les fleurs sont détruites dans le parc. ») — Repérage : \all{„Blumen im Park werden zerstört.“}

b) \all{Richtig, viele Bürger sind gegen die Kameras. Im Text steht: „Wir wollen Sicherheit, aber wir wollen auch unsere Freiheit.“} (« Vrai, beaucoup de citoyens sont contre les caméras. Dans le texte, il est dit : „Nous voulons la sécurité, mais nous voulons aussi notre liberté.“ »)

c) Réponse possible : \all{Ich meine, dass ein Jugendzentrum eine gute Lösung ist, weil die Jugendlichen dann Aktivitäten haben.} (« Je pense qu'un centre de jeunes est une bonne solution, parce que les jeunes ont alors des activités. ») — Toute opinion justifiée est acceptée, à condition que le verbe soit à la fin après \all{dass} et \all{weil}.

d) Résumé possible : \all{Jugendliche zerstören nachts Sachen in der Stadt, weil sie frustriert sind. Die Stadt installiert Kameras, aber viele Bürger sind dagegen. Ein Jugendzentrum soll das Problem lösen.} (« Des jeunes détruisent des choses la nuit dans la ville parce qu'ils sont frustrés. La ville installe des caméras, mais beaucoup de citoyens sont contre. Un centre de jeunes doit résoudre le problème. »)
\end{corrige}

\begin{aretenir}[L'essentiel de la section]
\begin{itemize}
\item Les \cle{W-Fragen} (\all{Wer? Was? Wo? Wann? Warum? Wie?}) annoncent le type d'information à chercher.
\item Méthode en quatre étapes : lire la question, repérer les mots-clés, souligner la phrase du texte, rédiger une phrase complète.
\item Au vrai/faux, toujours \textbf{justifier par une citation} entre guillemets allemands „ … “ ; jamais de simple \all{ja} ou \all{nein}.
\item Pour l'opinion : \all{Ich meine / finde / glaube, dass...} avec le verbe conjugué \textbf{à la fin}.
\item Le résumé : 3 à 4 phrases, une idée par phrase, reliées par \all{denn}, \all{deshalb}, \all{aber}.
\item Au passif, le complément d'agent se construit avec \all{von} + datif : \all{von der Polizei}, \all{von Jugendlichen}.
\end{itemize}
\end{aretenir}

Tu sais désormais répondre méthodiquement à toutes les questions de compréhension. La section suivante te donnera les outils lexicaux pour parler toi-même, avec précision, du vandalisme et du terrorisme.

\coursec{Lexique thématique : le vandalisme et le terrorisme}

Après avoir lu le texte support et répondu aux questions de compréhension détaillée, nous disposons maintenant d'une bonne vision globale du sujet. Mais pour parler et écrire correctement sur le vandalisme et le terrorisme, il faut maîtriser le vocabulaire précis du thème. C'est l'objectif de cette section : construire, champ lexical par champ lexical, un vocabulaire solide, avec l'article et le pluriel de chaque nom, les familles de mots et les expressions utiles. Retenez cette règle d'or de l'apprentissage de l'allemand : on n'apprend jamais un nom seul, on l'apprend toujours avec son \cle{article} (\all{der}, \all{die}, \all{das}) et son \cle{pluriel}.

\courssub{Observation : le vocabulaire dans son contexte}

Avant toute liste de mots, observons d'abord le vocabulaire en action dans un court texte original. Lisez-le une première fois sans dictionnaire, puis repérez les mots que vous comprenez déjà grâce au français.

\begin{textebox}[Gewalt in der Stadt — ein Zeitungsartikel]
Am Montagabend haben unbekannte Täter in der Innenstadt mehrere Busse beschädigt und die Wände des Bahnhofs beschmiert. Die Zerstörung ist groß: Fenster sind kaputt, Sitze sind zerschnitten. Die Polizei hat sofort reagiert. „Wir verstärken die Sicherheit in der ganzen Stadt“, sagt ein Polizist. „Neue Kameras kontrollieren jetzt den Bahnhof und den Marktplatz.“

Die Bürger der Stadt haben Angst. „Ich habe Angst vor der Gewalt“, sagt eine Frau am Bahnhof. „Früher war unsere Stadt sicher, heute nicht mehr.“ Ein junger Mann meint: „Der Staat muss mehr Kontrollen machen und die Täter schneller finden.“

Auch der Terrorismus ist ein Thema in den Nachrichten. Letztes Jahr konnte die Polizei einen Anschlag verhindern: Der Anschlag wurde von der Polizei verhindert, bevor die Terroristen ihn verüben konnten. Viele Menschen waren damals Opfer von Angst und Panik. Die Regierung erklärt: „Wir klären jeden Angriff auf. Der Schutz der Bürger ist unsere wichtigste Aufgabe.“

Experten sagen: Vandalismus beginnt oft klein — ein beschmierter Tisch, ein kaputtes Fenster. Deshalb ist Prävention in den Schulen so wichtig. Wer die Gewalt versteht, kann sie verhindern.
\end{textebox}

\textbf{Glossaire}

\begin{center}
\begin{tabularx}{\linewidth}{|X|X|}\hline
\rowcolor{popblueL} \textbf{Deutsch} & \textbf{Français} \\ \hline
unbekannt & inconnu \\ \hline
der Bahnhof, die Bahnhöfe & la gare \\ \hline
zerschneiden (zerschnitten) & découper, lacérer \\ \hline
die Nachrichten (Pl.) & les informations, les nouvelles \\ \hline
die Panik (Sg.) & la panique \\ \hline
die Regierung, -en & le gouvernement \\ \hline
die Aufgabe, -n & la tâche, la mission \\ \hline
die Prävention (Sg.) & la prévention \\ \hline
\end{tabularx}
\end{center}

\textbf{Questions rapides}

\begin{enumerate}
\item Was haben die Täter beschädigt? (Qu'ont endommagé les auteurs des dégradations ?)
\item Wie reagiert die Polizei? (Comment réagit la police ?)
\item Wer hat Angst, und wovor? (Qui a peur, et de quoi ?)
\item Was konnte die Polizei letztes Jahr verhindern? (Qu'est-ce que la police a pu empêcher l'année dernière ?)
\end{enumerate}

\begin{corrige}[Questions rapides]
\begin{enumerate}
\item Die Täter haben mehrere Busse beschädigt und die Wände des Bahnhofs beschmiert.
\item Die Polizei verstärkt die Sicherheit und installiert neue Kameras.
\item Die Bürger haben Angst vor der Gewalt.
\item Die Polizei konnte einen Anschlag (der Terroristen) verhindern.
\end{enumerate}
\end{corrige}

\courssub{Champ lexical de la violence}

Observons d'abord les noms qui désignent la violence elle-même. Notez bien l'article et le pluriel de chaque mot.

\begin{center}
\begin{tabularx}{\linewidth}{|X|X|}\hline
\rowcolor{popblueL} \textbf{Deutsch} & \textbf{Français} \\ \hline
der Terror (Sg.) & la terreur \\ \hline
der Terrorismus (Sg.) & le terrorisme \\ \hline
der Vandalismus (Sg.) & le vandalisme \\ \hline
die Gewalt (Sg.) & la violence \\ \hline
der Anschlag, die Anschläge & l'attentat \\ \hline
der Angriff, die Angriffe & l'attaque \\ \hline
die Zerstörung, die Zerstörungen & la destruction \\ \hline
\end{tabularx}
\end{center}

\begin{exemplebox}[Exemples traduits]
\all{Der Anschlag wurde von der Polizei verhindert.} « L'attentat a été empêché par la police. »

\all{Der Vandalismus in den Schulen nimmt zu.} « Le vandalisme dans les écoles augmente. »

\all{Nach dem Angriff gab es viele Opfer.} « Après l'attaque, il y a eu de nombreuses victimes. »

\all{Die Zerstörung des Bahnhofs kostet viel Geld.} « La destruction de la gare coûte beaucoup d'argent. »
\end{exemplebox}

\begin{attention}[Faux ami : die Gewalt]
\all{Die Gewalt} signifie « la violence » ou « la force », et jamais « la gêne » ! Ne confondez pas avec le français « gêner » : « gêner, déranger » se dit \all{stören} en allemand. De même, « avec gêne » (embarrassé) se dit \all{verlegen}. Retenez : \all{Gewalt} = violence ; \all{Störung} = dérangement, gêne.
\end{attention}

\begin{propriete}[Noms abstraits sans pluriel]
Les noms abstraits comme \all{der Terrorismus}, \all{der Vandalismus}, \all{die Gewalt}, \all{der Terror} s'emploient presque toujours au singulier, sans article défini quand on parle du phénomène en général : \all{Terrorismus ist ein Problem unserer Zeit.} « Le terrorisme est un problème de notre époque. » En revanche, \all{der Anschlag} et \all{der Angriff} désignent des événements concrets et ont un pluriel : \all{die Anschläge}, \all{die Angriffe}.
\end{propriete}

\begin{attention}[der Anschlag ou der Angriff ?]
Les deux mots se traduisent par « attaque », mais ils ne sont pas interchangeables. \all{Der Anschlag} désigne un \cle{attentat}, une attaque criminelle ou terroriste : \all{ein Anschlag auf einen Polizisten} « un attentat contre un policier ». \all{Der Angriff} est plus général : attaque militaire, sportive ou verbale : \all{ein Angriff im Fußballspiel} « une attaque dans un match de football ». Dans le contexte du terrorisme, on dit presque toujours \all{der Anschlag}.
\end{attention}

\courssub{Champ lexical des acteurs}

Qui fait quoi ? Distinguons les auteurs des actes, les forces de l'ordre et les victimes.

\begin{center}
\begin{tabularx}{\linewidth}{|X|X|}\hline
\rowcolor{popblueL} \textbf{Deutsch} & \textbf{Français} \\ \hline
der Täter, die Täter & l'auteur (du délit, du crime) \\ \hline
der Terrorist, die Terroristen & le terroriste \\ \hline
die Polizei (Sg.) & la police \\ \hline
der Polizist, die Polizisten & le policier, l'agent de police \\ \hline
das Opfer, die Opfer & la victime \\ \hline
der Bürger, die Bürger & le citoyen \\ \hline
\end{tabularx}
\end{center}

\begin{exemplebox}[Exemples traduits]
\all{Die Täter sind noch unbekannt.} « Les auteurs sont encore inconnus. »

\all{Die Polizei sucht die Terroristen.} « La police recherche les terroristes. »

\all{Ein Polizist kontrolliert den Bahnhof.} « Un agent de police surveille la gare. »

\all{Die Opfer bekommen Hilfe vom Staat.} « Les victimes reçoivent l'aide de l'État. »

\all{Die Bürger informieren die Polizei.} « Les citoyens informent la police. »
\end{exemplebox}

\begin{attention}[die Polizei : toujours singulier]
En allemand, \all{die Polizei} est un nom collectif toujours au \cle{singulier} : \all{Die Polizei kommt.} « La police arrive. » (et non « die Polizei kommen »). Pour parler d'un agent précis, utilisez \all{der Polizist} (pluriel : \all{die Polizisten}) ; pour une femme agent : \all{die Polizistin, die Polizistinnen}. Ne dites jamais « ein Polizei » !
\end{attention}

\begin{propriete}[Le masculin faible : der Terrorist, der Polizist]
\all{Der Terrorist} et \all{der Polizist} appartiennent à la déclinaison dite « faible » : ils prennent la terminaison \cle{-en} à tous les cas sauf au nominatif singulier : \all{Ich sehe den Terroristen.} (accusatif), \all{mit dem Polizisten} (datif). Pluriel : \all{die Terroristen}, \all{die Polizisten}. Même déclinaison pour \all{der Student}, \all{der Junge}, \all{der Mensch}.
\end{propriete}

\courssub{Champ lexical des actions}

Voici les verbes essentiels du thème. Ce sont tous des verbes \cle{faibles} (réguliers), sauf mention contraire.

\begin{center}
\begin{tabularx}{\linewidth}{|X|X|}\hline
\rowcolor{popblueL} \textbf{Deutsch} & \textbf{Français} \\ \hline
beschädigen & endommager, dégrader \\ \hline
zerstören & détruire \\ \hline
beschmieren & barbouiller, couvrir de graffitis \\ \hline
verletzen & blesser \\ \hline
töten & tuer \\ \hline
verhindern & empêcher \\ \hline
verstärken & renforcer \\ \hline
aufklären & élucider, éclaircir (une affaire) \\ \hline
\end{tabularx}
\end{center}

\begin{exemplebox}[Exemples traduits]
\all{Jugendliche beschädigen die Bänke im Park.} « Des jeunes dégradent les bancs du parc. »

\all{Die Täter beschmieren die Wände mit Farbe.} « Les auteurs barbouillent les murs avec de la peinture. »

\all{Bei dem Angriff wurden drei Menschen verletzt.} « Lors de l'attaque, trois personnes ont été blessées. »

\all{Die Polizei will jeden Fall aufklären.} « La police veut élucider chaque affaire. »

\all{Die Sicherheit wird verstärkt.} « La sécurité est renforcée. »
\end{exemplebox}

\begin{attention}[beschädigen ou zerstören ?]
\all{Beschädigen} signifie « endommager, dégrader » : l'objet est abîmé mais existe encore (\all{ein beschädigtes Auto} « une voiture endommagée »). \all{Zerstören} signifie « détruire » : l'objet n'existe plus (\all{Die Bombe zerstört das Haus.} « La bombe détruit la maison. »). Pour le vandalisme « léger » (bancs, murs, vitres), on utilise plutôt \all{beschädigen}.
\end{attention}

\begin{attention}[verletzen : deux sens]
\all{Verletzen} signifie « blesser » (physiquement) mais aussi « violer » une règle : \all{das Gesetz verletzen} « violer la loi ». Attention aussi à ne pas confondre avec \all{verlieren} (perdre) : \all{Er verliert sein Geld.} « Il perd son argent. »
\end{attention}

\courssub{Champ lexical de la sécurité}

Face à la violence, la société répond par la sécurité. Voici le vocabulaire de cette réponse.

\begin{center}
\begin{tabularx}{\linewidth}{|X|X|}\hline
\rowcolor{popblueL} \textbf{Deutsch} & \textbf{Français} \\ \hline
die Sicherheit (Sg.) & la sécurité \\ \hline
die Kamera, die Kameras & la caméra (de surveillance) \\ \hline
die Kontrolle, die Kontrollen & le contrôle \\ \hline
der Schutz (Sg.) & la protection \\ \hline
die Angst, die Ängste & la peur, l'angoisse \\ \hline
\end{tabularx}
\end{center}

\begin{exemplebox}[Exemples traduits]
\all{Die Kameras kontrollieren den Eingang der Schule.} « Les caméras surveillent l'entrée de l'école. »

\all{Der Schutz der Bürger ist wichtig.} « La protection des citoyens est importante. »

\all{Nach dem Anschlag haben viele Menschen Angst.} « Après l'attentat, beaucoup de gens ont peur. »

\all{Am Flughafen gibt es strenge Kontrollen.} « À l'aéroport, il y a des contrôles stricts. »
\end{exemplebox}

\begin{attention}[Faux ami : die Kontrolle]
\all{Die Kontrolle} signifie « le contrôle » au sens de « vérification, surveillance » (\all{die Passkontrolle} « le contrôle des passeports »), mais pas « une interro, un devoir surveillé » : à l'école, on dit \all{die Klassenarbeit} ou \all{die Prüfung}. De même, \all{kontrollieren} = « vérifier, surveiller », et non « contrôler » au sens de « maîtriser » (qui se dit \all{beherrschen}).
\end{attention}

\courssub{Familles de mots : un mot en cache plusieurs}

L'allemand construit des familles de mots autour d'une même racine. Apprendre la famille, c'est apprendre quatre ou cinq mots pour le prix d'un.

\begin{definition}[Familles de mots du thème]
\begin{itemize}
\item \all{die Gewalt} (la violence) → \all{gewalttätig} (violent, adj.) → \all{der Gewalttäter} (l'auteur de violences)
\item \all{der Terror} (la terreur) → \all{der Terrorist} (le terroriste) → \all{terroristisch} (terroriste, adj.) → \all{der Terrorismus} (le terrorisme)
\item \all{sicher} (sûr, adj.) → \all{die Sicherheit} (la sécurité) → \all{unsicher} (incertain) → \all{sichern} (sécuriser)
\item \all{zerstören} (détruire) → \all{die Zerstörung} (la destruction)
\item \all{die Angst} (la peur) → \all{ängstlich} (craintif) → \all{Angst haben vor + Dativ} (avoir peur de)
\end{itemize}
\end{definition}

\begin{popfigure}
\begin{center}
\begin{tikzpicture}[
  node distance=1.1cm,
  every node/.style={font=\small},
  racine/.style={rectangle, rounded corners, draw=popblue, fill=popblueL, thick, align=center, minimum width=2.2cm},
  enfant/.style={rectangle, rounded corners, draw=popgreen, fill=popgreenL, thick, align=center, minimum width=2.2cm},
  fleche/.style={-{Stealth[length=2.5mm]}, thick, popdark}
]
\node[racine, align=center] (sicher){\all{sicher}\\ (sûr, adj.)};
\node[enfant, above left=1.1cm and 1.6cm of sicher, align=center] (sicherheit){\all{die Sicherheit}\\ (la sécurité)};
\node[enfant, above right=1.1cm and 1.6cm of sicher, align=center] (unsicher){\all{unsicher}\\ (incertain)};
\node[enfant, below=1.1cm of sicher, align=center] (sichern){\all{sichern}\\ (sécuriser)};
\node[enfant, below=1.1cm of sichern, align=center] (versichern){\all{die Versicherung}\\ (l'assurance)};
\draw[fleche] (sicher) -- (sicherheit);
\draw[fleche] (sicher) -- (unsicher);
\draw[fleche] (sicher) -- (sichern);
\draw[fleche] (sichern) -- (versichern);
\end{tikzpicture}
\end{center}
\legende{Arbre de la famille de mots autour de \all{sicher} : la racine adjective donne un nom (\all{die Sicherheit}), un contraire avec le préfixe \all{un-} (\all{unsicher}), un verbe (\all{sichern}) et un nom dérivé (\all{die Versicherung}).}
\end{popfigure}

\begin{savaistu}[D'où vient le mot « Terror » ?]
Le mot \all{Terror} vient du latin \emph{terror}, qui signifie « terreur, grande frayeur ». L'allemand et le français partagent cette racine latine : \all{der Terror} / « la terreur », \all{der Terrorist} / « le terroriste », \all{terroristisch} / « terroriste ». Cette parenté facilite la mémorisation : quand vous voyez un mot allemand en \all{-ismus} (\all{Terrorismus}, \all{Vandalismus}, \all{Optimismus}), il correspond presque toujours au mot français en « -isme ». Bonne nouvelle supplémentaire : ces noms en \all{-ismus} sont tous masculins en allemand (\all{der}), comme leurs équivalents français — un cas rare où les deux langues sont d'accord sur le genre ! Le mot \all{Vandalismus}, lui, rappelle les Vandales, un peuple germanique de l'Antiquité tardive dont le nom est devenu synonyme de destruction gratuite.
\end{savaistu}

\courssub{Expressions utiles à mémoriser}

\begin{aretenir}[Trois expressions clés]
\begin{itemize}
\item \all{einen Anschlag verüben} : commettre un attentat. \all{Die Terroristen verüben einen Anschlag.} « Les terroristes commettent un attentat. »
\item \all{Angst haben vor + Dativ} : avoir peur de. \all{Die Bürger haben Angst vor der Gewalt.} « Les citoyens ont peur de la violence. » Attention au datif : \all{vor dem Terror}, \all{vor den Anschlägen}.
\item \all{die Sicherheit verstärken} : renforcer la sécurité. \all{Nach dem Angriff verstärkt die Stadt die Sicherheit.} « Après l'attaque, la ville renforce la sécurité. »
\end{itemize}
\end{aretenir}

\begin{center}
\begin{tabularx}{\linewidth}{|X|X|X|}\hline
\rowcolor{popblueL} \textbf{Français} & \textbf{Deutsch} & \textbf{Remarque} \\ \hline
commettre un attentat & einen Anschlag verüben & verbe faible : verübt, verübte, hat verübt \\ \hline
avoir peur de la violence & Angst vor der Gewalt haben & \all{vor} + datif \\ \hline
renforcer la sécurité & die Sicherheit verstärken & verbe faible \\ \hline
empêcher un attentat & einen Anschlag verhindern & verbe faible \\ \hline
élucider une affaire & einen Fall aufklären & particule séparable : \all{klärt ... auf} \\ \hline
\end{tabularx}
\end{center}

\begin{attention}[Verbes à particule séparable ou inséparable]
\all{Aufklären} est un verbe à particule \cle{séparable} : à la deuxième position de la phrase, la particule \all{auf-} va à la fin : \all{Die Polizei klärt den Fall auf.} « La police élucide l'affaire. » Mais à l'infinitif avec \all{zu} ou après un modal, le verbe reste entier : \all{Der Fall muss aufgeklärt werden.} « L'affaire doit être élucidée. » En revanche, \all{verhindern}, \all{verstärken}, \all{zerstören}, \all{beschädigen}, \all{verletzen} sont \cle{inséparables} (préfixes \all{ver-}, \all{zer-}, \all{be-}) : leur participe passé ne prend pas de \all{ge-} : \all{beschädigt}, \all{verhindert}, \all{zerstört}.
\end{attention}

\courssub{Exercice résolu et entraînement}

\begin{exoresolu}[Complétez avec le mot qui convient]
Complétez les phrases avec : \all{der Anschlag}, \all{die Polizei}, \all{das Opfer}, \all{die Sicherheit}, \all{beschädigt}, \all{verhindern}.

1. Unbekannte Täter haben das Schulgebäude \ldots\ 
2. \ldots\ sucht die Täter.
3. Die Stadt will einen Angriff \ldots\ 
4. Nach dem \ldots\ haben die Bürger Angst.
5. Die Kameras verstärken \ldots\ am Bahnhof.
6. \ldots\ des Angriffs liegt im Krankenhaus.
\tcblower
\textbf{Solution détaillée}

1. \all{Unbekannte Täter haben das Schulgebäude beschädigt.} « Des auteurs inconnus ont endommagé le bâtiment scolaire. » — participe passé de \all{beschädigen} (verbe faible à préfixe inséparable : pas de \all{ge-}).

2. \all{Die Polizei sucht die Täter.} « La police recherche les auteurs. » — singulier obligatoire : \all{sucht}, pas « suchen ».

3. \all{Die Stadt will einen Angriff verhindern.} « La ville veut empêcher une attaque. » — infinitif après \all{wollen}.

4. \all{Nach dem Anschlag haben die Bürger Angst.} « Après l'attentat, les citoyens ont peur. » — \all{nach} + datif : \all{dem Anschlag}.

5. \all{Die Kameras verstärken die Sicherheit am Bahnhof.} « Les caméras renforcent la sécurité à la gare. » — accusatif : \all{die Sicherheit} (invariable au féminin singulier).

6. \all{Das Opfer des Angriffs liegt im Krankenhaus.} « La victime de l'attaque est à l'hôpital. » — nominatif neutre : \all{das Opfer} ; \all{des Angriffs} est un génitif.
\end{exoresolu}

\begin{exercice}[À vous de jouer]
1. Traduisez en allemand : a) « L'attentat a été empêché par la police. » b) « Les citoyens ont peur du terrorisme. » c) « Des jeunes ont barbouillé les murs du lycée. »

2. Formez le nom ou l'adjectif de la même famille : a) \all{zerstören} → (nom) \ldots\ b) \all{sicher} → (nom) \ldots\ c) \all{der Terror} → (personne) \ldots\ d) \all{die Gewalt} → (adjectif) \ldots\

3. Rédigez trois phrases sur la sécurité dans votre ville en utilisant : \all{die Sicherheit verstärken}, \all{die Kamera}, \all{Angst haben vor}.
\end{exercice}

\begin{corrige}[Exercice « À vous de jouer »]
1. a) \all{Der Anschlag wurde von der Polizei verhindert.} b) \all{Die Bürger haben Angst vor dem Terrorismus.} c) \all{Jugendliche haben die Wände des Gymnasiums beschmiert.} (En Allemagne, \all{das Gymnasium} désigne le lycée ; au Cameroun, on peut aussi dire \all{die Schule} ou \all{das Lyzeum}.)

2. a) \all{die Zerstörung} ; b) \all{die Sicherheit} ; c) \all{der Terrorist} ; d) \all{gewalttätig}.

3. Exemple de production : \all{In meiner Stadt verstärkt die Polizei die Sicherheit. Neue Kameras kontrollieren den Markt. Viele Bürger haben Angst vor der Gewalt, aber die Polizei schützt uns.}
\end{corrige}

\begin{aretenir}[L'essentiel du lexique]
\begin{itemize}
\item La violence : \all{der Terrorismus} (Sg.), \all{der Vandalismus} (Sg.), \all{die Gewalt} (Sg.), \all{der Anschlag} (Pl. die Anschläge), \all{der Angriff} (Pl. die Angriffe), \all{die Zerstörung} (Pl. die Zerstörungen).
\item Les acteurs : \all{der Täter} (Pl. die Täter), \all{der Terrorist} (Pl. die Terroristen), \all{die Polizei} (singulier !), \all{der Polizist} (Pl. die Polizisten), \all{das Opfer} (Pl. die Opfer), \all{der Bürger} (Pl. die Bürger).
\item Les actions : \all{beschädigen}, \all{zerstören}, \all{beschmieren}, \all{verletzen}, \all{töten}, \all{verhindern}, \all{verstärken}, \all{aufklären} (séparable !).
\item La sécurité : \all{die Sicherheit} (Sg.), \all{die Kamera} (Pl. die Kameras), \all{die Kontrolle} (Pl. die Kontrollen), \all{der Schutz} (Sg.), \all{die Angst} (Pl. die Ängste).
\item Expressions : \all{einen Anschlag verüben}, \all{Angst haben vor + Dativ}, \all{die Sicherheit verstärken}.
\end{itemize}
\end{aretenir}

Ce vocabulaire est désormais votre boîte à outils. Dans la section suivante, nous verrons comment ces mots s'organisent dans la phrase grâce à une structure grammaticale essentielle du thème : le \cle{passif}, qui permet de dire « l'attentat \emph{a été empêché} » sans nommer l'auteur.

\coursec{Grammaire 1 : le passif simple}

Dans la section précédente, nous avons découvert le vocabulaire du vandalisme et du terrorisme : \all{der Anschlag}, \all{die Gewalt}, \all{beschädigen}, \all{die Sicherheit}. Or, dans les textes journalistiques qui traitent de ces sujets, une construction grammaticale revient sans cesse : \all{Bushaltestellen werden beschädigt}, \all{Die Sicherheit wird verstärkt}. Cette construction s'appelle le \cle{passif}. C'est la première règle de grammaire de notre leçon, et elle est indispensable pour comprendre et produire des textes d'actualité.

\courssub{Observation : découvrir le passif dans le texte}

\begin{textebox}[Vandalismus in der Stadt — ein Zeitungsartikel]
In der Nacht von Samstag auf Sonntag wurden in unserer Stadt mehrere Bushaltestellen beschädigt. Scheiben wurden eingeschlagen, und Wände wurden mit Farbe beschmiert. Auch zwei Autos wurden zerstört. Die Polizei wurde am Morgen informiert. Ein Zeuge wurde befragt, aber der Täter wurde noch nicht gefunden. „Die Sicherheit wird jetzt verstärkt“, sagte ein Polizist. Kameras werden an den Haltestellen installiert, und mehr Polizisten werden in der Nacht eingesetzt. Viele Bürger sind schockiert: „Warum wird so etwas gemacht?“, fragt eine Mutter. Die Schäden werden von der Stadt bezahlt. Der Bürgermeister erklärt: „Solche Taten werden bestraft. Wer gesehen hat, muss die Polizei informieren.“ In den Schulen wird jetzt über Gewalt und Respekt diskutiert. Plakate werden gemacht, und ein Projekt gegen Vandalismus wird organisiert. Die Stadt hofft: Solche Taten werden in Zukunft verhindert.
\end{textebox}

\textbf{Glossaire :} die Bushaltestelle, -n (l'arrêt de bus) ; die Scheibe, -n (la vitre) ; einschlagen (briser, part. passé : eingeschlagen) ; beschmieren (salir, barbouiller) ; der Zeuge, -n (le témoin) ; befragen (interroger) ; der Bürger, - (le citoyen) ; der Schaden, ¨- (le dégât) ; bestrafen (punir) ; das Plakat, -e (l'affiche) ; verhindern (empêcher) ; einsetzen (déployer, engager ; part. passé : eingesetzt).

\textbf{Questions de compréhension :}
\begin{enumerate}
\item Was wurde in der Nacht beschädigt?
\item Wer wurde am Morgen informiert?
\item Was wird an den Haltestellen installiert?
\end{enumerate}

\begin{aretenir}[Ce que nous observons]
Dans ce texte, presque toutes les phrases ont la même structure :
\begin{center}
\all{werden} (conjugué) \quad + \quad \textbf{Partizip II à la fin de la phrase}
\end{center}
\all{Bushaltestellen \textbf{werden} beschädigt.} — \all{Die Sicherheit \textbf{wird} verstärkt.} — \all{Wände \textbf{wurden} beschmiert.}
\end{aretenir}

\courssub{La règle de formation}

\begin{propriete}[Formation du passif : werden + Partizip II]
Le passif allemand (das \cle{Passiv}, plus précisément le \cle{Vorgangspassif} ou passif de processus) se forme avec l'auxiliaire \all{werden}, conjugué à la personne du sujet, suivi du \cle{participe passé} (Partizip II) du verbe principal, placé \textbf{toujours à la fin de la phrase}.
\begin{center}
Sujet (Nominativ) + \all{werden} conjugué + (compléments) + Partizip II
\end{center}
\all{Der Täter \textbf{wird} verhaftet.} « Le coupable est arrêté. »
\end{propriete}

\textbf{Conjugaison de \all{werden} au présent :}

\begin{center}
\begin{tabular}{|l|l|l|}
\hline
\rowcolor{popblueL} Personne & werden + Partizip II & Exemple \\ \hline
ich & werde & Ich \textbf{werde} befragt. \\ \hline
du & wirst & Du \textbf{wirst} informiert. \\ \hline
er/sie/es & wird & Der Täter \textbf{wird} verhaftet. \\ \hline
wir & werden & Wir \textbf{werden} geschützt. \\ \hline
ihr & werdet & Ihr \textbf{werdet} gewarnt. \\ \hline
sie/Sie & werden & Die Wände \textbf{werden} beschmiert. \\ \hline
\end{tabular}
\end{center}

\begin{exemplebox}[Exemples traduits — Présent]
\begin{itemize}
\item \all{Der Täter wird verhaftet.} « Le coupable est arrêté. »
\item \all{Die Bushaltestelle wird repariert.} « L'arrêt de bus est réparé. »
\item \all{Kameras werden installiert.} « Des caméras sont installées. »
\item \all{Die Sicherheit wird verstärkt.} « La sécurité est renforcée. »
\item \all{Ein Zeuge wird befragt.} « Un témoin est interrogé. »
\end{itemize}
\end{exemplebox}

\courssub{Le passif au prétérit}

Pour raconter des événements passés (comme dans notre article de presse), on utilise le passif au \cle{prétérit} : \all{wurde / wurden} + Partizip II.

\begin{center}
\begin{tabular}{|l|l|l|}
\hline
\rowcolor{popblueL} Personne & Präteritum de werden & Exemple \\ \hline
ich & wurde & Ich \textbf{wurde} befragt. \\ \hline
du & wurdest & Du \textbf{wurdest} informiert. \\ \hline
er/sie/es & wurde & Das Auto \textbf{wurde} zerstört. \\ \hline
wir & wurden & Wir \textbf{wurden} gewarnt. \\ \hline
ihr & wurdet & Ihr \textbf{wurdet} geschützt. \\ \hline
sie/Sie & wurden & Die Wände \textbf{wurden} beschmiert. \\ \hline
\end{tabular}
\end{center}

\begin{exemplebox}[Exemples traduits — Prétérit]
\begin{itemize}
\item \all{Die Wände wurden beschmiert.} « Les murs ont été barbouillés. »
\item \all{Viele Menschen wurden verletzt.} « Beaucoup de personnes ont été blessées. »
\item \all{Die Polizei wurde informiert.} « La police a été informée. »
\item \all{Zwei Autos wurden zerstört.} « Deux voitures ont été détruites. »
\end{itemize}
\end{exemplebox}

\courssub{Transformation actif $\rightarrow$ passif}

\begin{methode}[Transformer une phrase active en phrase passive]
\begin{enumerate}
\item \textbf{Repérer le complément d'objet direct} (Akkusativ) de la phrase active : \all{Der Täter beschädigt \textbf{das Auto}.}
\item \textbf{Le transformer en sujet} : l'accusatif devient le nominatif du passif : \all{\textbf{Das Auto} ...}
\item \textbf{Conjuguer \all{werden}} à la personne de ce nouveau sujet : \all{Das Auto \textbf{wird} ...}
\item \textbf{Placer le participe passé à la fin} : \all{Das Auto wird \textbf{beschädigt}.}
\item (Facultatif) L'ancien sujet devient \all{von} + datif, ou disparaît : \all{Das Auto wird \textbf{vom Täter} beschädigt.}
\end{enumerate}
\end{methode}

\begin{popfigure}
\begin{center}
\begin{tikzpicture}[node distance=0.6cm, >=Stealth]
\node[draw=popblue, fill=popblueL, rounded corners, align=center, text width=5.2cm] (aktiv) {\textbf{Aktiv}\\[2pt] \all{Der Täter} (Nom.)\\ \all{beschädigt} (Verb)\\ \all{das Auto} (Akk.)};
\node[draw=popgreen, fill=popgreenL, rounded corners, align=center, text width=5.2cm, right=3.2cm of aktiv] (passiv) {\textbf{Passiv}\\[2pt] \all{Das Auto} (Nom.)\\ \all{wird} (werden)\\ \all{beschädigt} (Partizip II)};
\draw[->, very thick, poporange] (aktiv) -- node[above, align=center, font=\small]{Akkusativ $\rightarrow$ Nominativ} (passiv);
\node[draw=poppurple, fill=poppinkL, rounded corners, align=center, text width=4.5cm, below=1cm of aktiv, xshift=3.2cm] (von) {L'agent : \all{von} + Dativ\\ \all{vom Täter} (souvent omis)};
\draw[->, thick, poppurple, dashed] (von) -- (passiv.south);
\end{tikzpicture}
\legende{Schéma de transformation : de l'actif au passif}
\end{center}
\end{popfigure}

\begin{center}
\begin{tabularx}{\linewidth}{|X|X|}
\hline
\rowcolor{popblueL} Aktiv (actif) & Passiv (passif) \\ \hline
\all{Der Täter beschädigt das Auto.} & \all{Das Auto wird (vom Täter) beschädigt.} \\ \hline
\all{Die Polizei verhaftet den Mann.} & \all{Der Mann wird (von der Polizei) verhaftet.} \\ \hline
\all{Jugendliche beschmierten die Wände.} & \all{Die Wände wurden (von Jugendlichen) beschmiert.} \\ \hline
\all{Die Stadt installiert Kameras.} & \all{Kameras werden (von der Stadt) installiert.} \\ \hline
\end{tabularx}
\end{center}

\courssub{Emploi du passif}

\begin{definition}[Quand utilise-t-on le passif (Vorgangspassiv) ?]
On emploie le passif de processus quand :
\begin{itemize}
\item l'\textbf{auteur de l'action est inconnu} : \all{Die Scheiben wurden eingeschlagen.} (on ne sait pas par qui) ;
\item l'auteur est \textbf{sans importance} : \all{Die Schäden werden bezahlt.} (peu importe qui paie) ;
\item on veut \textbf{insister sur l'action} et non sur l'auteur — c'est le style typique du \textbf{journalisme} et des textes administratifs : \all{Die Sicherheit wird verstärkt.}
\end{itemize}
\end{definition}

\courssub{Comparaison avec le français}

\begin{center}
\begin{tabularx}{\linewidth}{|X|X|X|}
\hline
\rowcolor{popblueL} Français & Deutsch & Remarque \\ \hline
Le coupable est arrêté. & \all{Der Täter wird verhaftet.} & \all{werden} = auxiliaire du passif \\ \hline
Les murs ont été barbouillés. & \all{Die Wände wurden beschmiert.} & prétérit : \all{wurden} \\ \hline
Il est blessé (action). & \all{Er wird verletzt.} & passif : \all{werden} + Partizip II \\ \hline
Il blessera (futur). & \all{Er wird verletzen.} & futur : \all{werden} + infinitif \\ \hline
\end{tabularx}
\end{center}

\begin{attention}[Erreur n° 1 des francophones : « ist verhaftet »]
En français, le passif se forme avec « être » : « il \textbf{est} arrêté ». Les élèves francophones traduisent donc souvent par \all{Er \textbf{ist} verhaftet} — c'est une \textbf{erreur} ! En allemand, le passif de processus se forme \textbf{toujours avec \all{werden}}, jamais avec \all{sein} :
\begin{itemize}
\item Correct : \all{Der Täter \textbf{wird} verhaftet.} « Le coupable est arrêté (action en cours). »
\item Faux : \all{Der Täter \textbf{ist} verhaftet.}
\end{itemize}
\all{sein} + Partizip II correspond au \textbf{passif d'état (Zustandspassiv)} : \all{Der Täter \textbf{ist} verhaftet.} « Le coupable \emph{est} arrêté (état résultant) » ou au \textbf{parfait passif} : \all{Er \textbf{ist} verhaftet \textbf{worden}.} « Il a été arrêté. »
\end{attention}

\begin{attention}[Ne pas confondre passif et futur]
\all{werden} sert aussi à former le futur. La différence se voit au verbe final :
\begin{itemize}
\item Passif : \all{Er wird \textbf{verletzt}.} « Il est blessé. » (\all{werden} + \textbf{Partizip II})
\item Futur : \all{Er wird \textbf{verletzen}.} « Il blessera. » (\all{werden} + \textbf{infinitif})
\end{itemize}
Regardez toujours la fin de la phrase : participe passé (\emph{ge-...-t/en}) = passif ; infinitif (\emph{-en}) = futur.
\end{attention}

\courssub{Exercice résolu}

\begin{exoresolu}[Transformation actif $\rightarrow$ passif]
Énoncé : Mettez les phrases actives suivantes au passif (présent ou prétérit selon le temps de la phrase active).
\begin{enumerate}
\item \all{Der Täter zerstört das Auto.}
\item \all{Die Polizei befragte den Zeugen.}
\item \all{Die Stadt verstärkt die Sicherheit.}
\end{enumerate}
\tcblower
Solution détaillée :
\begin{enumerate}
\item L'accusatif est \all{das Auto} ; il devient nominatif. Le présent \all{zerstört} devient \all{wird} + participe \all{zerstört} à la fin : \all{\textbf{Das Auto wird (vom Täter) zerstört.}} « La voiture est détruite (par le coupable). »
\item L'accusatif est \all{den Zeugen} ; le prétérit \all{befragte} devient \all{wurde} + \all{befragt} : \all{\textbf{Der Zeuge wurde (von der Polizei) befragt.}} « Le témoin a été interrogé (par la police). »
\item L'accusatif est \all{die Sicherheit} ; présent : \all{\textbf{Die Sicherheit wird (von der Stadt) verstärkt.}} « La sécurité est renforcée (par la ville). »
\end{enumerate}
\end{exoresolu}

\begin{exercice}[À vous de jouer]
Mettez au passif : 
\begin{enumerate}
\item \all{Jugendliche beschmieren die Wände.}
\item \all{Man installiert Kameras.}
\item \all{Die Polizei verhaftete zwei Täter.}
\end{enumerate}
\end{exercice}

\begin{corrige}[Exercice « À vous de jouer »]
\begin{enumerate}
\item \all{Die Wände werden (von Jugendlichen) beschmiert.}
\item \all{Kameras werden installiert.} (pas d'agent avec \all{Man})
\item \all{Zwei Täter wurden (von der Polizei) verhaftet.}
\end{enumerate}
\end{corrige}

\begin{aretenir}[L'essentiel du passif]
\begin{itemize}
\item Passif (Vorgangspassiv) = \all{werden} conjugué + \textbf{Partizip II en fin de phrase}.
\item Présent : \all{wird/werden} ; prétérit : \all{wurde/wurden}.
\item Actif $\rightarrow$ passif : l'accusatif devient nominatif ; l'ancien sujet devient \all{von} + datif (ou disparaît).
\item Jamais \all{sein} au passif de processus ; attention à ne pas confondre avec le futur \all{werden} + infinitif.
\end{itemize}
\end{aretenir}

\coursec{Grammaire 2 : la cause et la conséquence}

Dans la section précédente, nous avons étudié le \cle{passif simple} : nous avons vu comment dire « les vitres sont brisées » (\all{die Fenster werden zerbrochen}) en mettant l'accent sur l'action et non sur son auteur. Mais un bon citoyen ne se contente pas de constater les dégâts : il explique \emph{pourquoi} ils se produisent et \emph{quelles} en sont les conséquences. Pourquoi la sécurité est-elle renforcée ? Que demandent les habitants parce qu'ils ont peur ? C'est exactement ce que permettent les \cle{connecteurs de cause et de conséquence}, indispensables pour commenter notre texte sur le vandalisme et le terrorisme et pour rédiger un paragraphe argumenté.

\courssub{Observation : deux phrases du texte}

Reprenons deux phrases tirées de notre thème et observons-les attentivement.

\begin{exemplebox}[Observation]
\all{Wegen dieser Anschläge wird die Sicherheit verstärkt.}\\
« À cause de ces attentats, la sécurité est renforcée. »

\all{Viele Bürger haben Angst, deshalb fordern sie mehr Polizisten.}\\
« Beaucoup de citoyens ont peur, c'est pourquoi ils demandent plus de policiers. »
\end{exemplebox}

Que remarque-t-on ? Dans la première phrase, la cause est exprimée par un \emph{groupe nominal} introduit par \all{wegen} : \all{wegen dieser Anschläge}. Dans la seconde, la conséquence est introduite par \all{deshalb}, et le sujet \all{sie} passe \textbf{après} le verbe \all{fordern} : c'est l'inversion. En allemand, exprimer la cause ou la conséquence, c'est d'abord choisir un connecteur, puis appliquer la \textbf{place du verbe} qui lui correspond. C'est toute la difficulté — et toute la logique — de cette leçon.

\courssub{La cause : weil, denn, wegen}

\textbf{1. La subordonnée de cause avec \all{weil}}

\all{Weil} signifie « parce que ». C'est une \cle{conjonction de subordination} : elle renvoie le verbe conjugué \textbf{à la fin} de la proposition.

\begin{propriete}[La cause avec weil]
Après \all{weil}, le verbe conjugué va \textbf{à la dernière place} de la subordonnée. Si la subordonnée est placée en tête de phrase, elle occupe la position 1 : le verbe de la principale vient immédiatement après la virgule (inversion du sujet).

\all{Die Bürger fordern mehr Kameras, weil sie Angst \textbf{haben}.}\\
\all{\textbf{Weil} die Bürger Angst \textbf{haben}, \textbf{fordern} sie mehr Kameras.}\\
« Parce que les citoyens ont peur, ils demandent plus de caméras. »
\end{propriete}

\textbf{2. La coordination de cause avec \all{denn}}

\all{Denn} signifie aussi « parce que / car », mais c'est une \cle{conjonction de coordination} : le verbe reste en \textbf{deuxième position}, comme dans une phrase normale. \all{Denn} ne peut jamais se placer en tête de phrase ; il introduit toujours la deuxième proposition.

\begin{propriete}[La cause avec denn]
Après \all{denn}, l'ordre des mots est celui d'une phrase principale : sujet + verbe en position 2.

\all{Sie fordern mehr Kameras, denn sie \textbf{haben} Angst.}\\
« Elles demandent plus de caméras, car elles ont peur. »
\end{propriete}

\textbf{3. La cause avec \all{wegen} + nom}

\all{Wegen} (« à cause de ») est une \cle{préposition} : elle est suivie d'un nom, pas d'une phrase avec verbe conjugué. À l'écrit, elle exige le \textbf{génitif}.

\begin{propriete}[La cause avec wegen]
\all{Wegen} + génitif : \all{wegen \textbf{des} Terrors}, \all{wegen \textbf{der} Anschläge}, \all{wegen \textbf{des} Vandalismus}.\\
\all{Wegen dieser Anschläge wird die Sicherheit verstärkt.}\\
« À cause de ces attentats, la sécurité est renforcée. »
\end{propriete}

\begin{attention}[Wegen + génitif, pas + datif]
À l'oral familier, les Allemands disent souvent \all{wegen dem Vandalismus} (datif). À l'écrit — et surtout au Bac — cette forme est une faute : écrivez \all{wegen \textbf{des} Vandalismus}, \all{wegen \textbf{der} Gewalt}, \all{wegen \textbf{des} Anschlags}. Rappel du génitif : masculin et neutre \all{des} (+ s au nom : \all{des Anschlag\textbf{s}}), féminin et pluriel \all{der}.
\end{attention}

\begin{exemplebox}[La cause : trois constructions pour une même idée]
\all{\textbf{Weil} Bushaltestellen beschädigt \textbf{werden}, kostet das viel Geld.}\\
« Parce que les arrêts de bus sont endommagés, cela coûte cher. »

\all{Das kostet viel Geld, \textbf{denn} Bushaltestellen \textbf{werden} beschädigt.}\\
« Cela coûte cher, car les arrêts de bus sont endommagés. »

\all{\textbf{Wegen der} beschädigten Bushaltestellen kostet das viel Geld.}\\
« À cause des arrêts de bus endommagés, cela coûte cher. »
\end{exemplebox}

\courssub{La conséquence : deshalb, darum, deswegen, so dass}

\textbf{1. \all{Deshalb / darum / deswegen} : « c'est pourquoi »}

Ces trois adverbes sont synonymes et s'emploient exactement de la même façon. Ils occupent la \textbf{première position} de la phrase de conséquence : le verbe conjugué reste en position 2, et le sujet passe \emph{après} le verbe (inversion).

\begin{propriete}[La conséquence avec deshalb]
\all{Deshalb} (ainsi que \all{darum} et \all{deswegen}) = position 1 ; verbe en position 2 ; sujet après le verbe.

\all{Viele Bürger haben Angst, \textbf{deshalb fordern sie} mehr Polizisten.}\\
« Beaucoup de citoyens ont peur, c'est pourquoi ils demandent plus de policiers. »
\end{propriete}

\begin{attention}[Inversion obligatoire après deshalb]
Le francophone écrit spontanément : \emph{\all{deshalb die Polizei fordert mehr Kontrollen}} — faute grave ! Après \all{deshalb}, le verbe doit venir en deuxième position : \all{\textbf{deshalb fordert} die Polizei mehr Kontrollen.} La règle est absolue : en allemand, le verbe conjugué est \emph{toujours} en deuxième position dans une phrase principale, quelle que soit la première position occupée (sujet, adverbe, complément).
\end{attention}

\textbf{2. La subordonnée de conséquence avec \all{so dass}}

\all{So dass} (« si bien que, de sorte que ») est une conjonction de subordination : comme \all{weil}, elle renvoie le verbe conjugué \textbf{à la fin}.

\begin{propriete}[La conséquence avec so dass]
Après \all{so dass}, le verbe conjugué va à la dernière place.

\all{Die Gewalt nimmt zu, \textbf{so dass} die Polizei mehr Kontrollen \textbf{macht}.}\\
« La violence augmente, si bien que la police fait plus de contrôles. »
\end{propriete}

\begin{exemplebox}[La conséquence : exemples traduits]
\all{Es gab einen Anschlag, \textbf{deshalb wurde} der Bahnhof geschlossen.}\\
« Il y a eu un attentat, c'est pourquoi la gare a été fermée. »

\all{Die Schule wurde beschädigt, \textbf{darum bleibt} sie heute geschlossen.}\\
« L'école a été endommagée, c'est pourquoi elle reste fermée aujourd'hui. »

\all{Der Terror verbreitet Angst, \textbf{so dass} viele Menschen zu Hause \textbf{bleiben}.}\\
« La terreur répand la peur, si bien que beaucoup de gens restent chez eux. »
\end{exemplebox}

\courssub{Tableaux récapitulatifs et comparaison avec le français}

\begin{center}
\begin{tabularx}{\linewidth}{|l|l|X|}
\hline
\rowcolor{popblueL} Connecteur & Place du verbe & Exemple \\
\hline
\all{weil} (parce que) & verbe à la fin & \all{Weil sie Angst haben, fordern sie Kameras.} \\
\hline
\all{denn} (car) & verbe en position 2 & \all{Sie fordern Kameras, denn sie haben Angst.} \\
\hline
\all{wegen} + Gén. (à cause de) & pas de verbe (nom) & \all{wegen der Anschläge} \\
\hline
\all{deshalb/darum/deswegen} (c'est pourquoi) & verbe en position 2, inversion & \all{Deshalb fordern sie Kameras.} \\
\hline
\all{so dass} (si bien que) & verbe à la fin & \all{..., so dass sie Kameras fordern.} \\
\hline
\end{tabularx}
\end{center}

\begin{center}
\begin{tabularx}{\linewidth}{|X|X|X|}
\hline
\rowcolor{popblueL} Français & Deutsch & Remarque \\
\hline
parce que + phrase & \all{weil} / \all{denn} & \all{weil} : verbe à la fin ; \all{denn} : verbe en position 2 \\
\hline
à cause de + nom & \all{wegen} + génitif & génitif obligatoire à l'écrit \\
\hline
c'est pourquoi / donc & \all{deshalb}, \all{darum}, \all{deswegen} & inversion du sujet \\
\hline
si bien que / de sorte que & \all{so dass} & verbe à la fin \\
\hline
\end{tabularx}
\end{center}

\begin{popfigure}
\begin{tikzpicture}[
  col1/.style={draw=popblue, fill=popblueL, rounded corners, align=center, font=\small, minimum width=3.1cm, minimum height=1.5cm},
  col2/.style={draw=poporange, fill=poporangeL, rounded corners, align=center, font=\small, minimum width=3.4cm, minimum height=1.5cm},
  col3/.style={draw=popgreen, fill=popgreenL, rounded corners, align=center, font=\small, minimum width=4.6cm, minimum height=1.5cm},
  head/.style={font=\bfseries\small, align=center}
]
\node[head, text=popblue] at (0,1.2) {Connecteur};
\node[head, text=poporange] at (3.9,1.2) {Type de phrase};
\node[head, text=popgreen] at (8.6,1.2) {Exemple traduit};
\node[col1] at (0,0) {\all{weil}};
\node[col2, align=center] at (3.9,0){subordonnée\\verbe à la fin};
\node[col3, align=center] at (8.6,0){\all{Weil sie Angst haben,}\\\all{fordern sie Kameras.}\\« Parce qu'ils ont peur,\\ils demandent des caméras. »};
\node[col1] at (0,-2.2) {\all{denn}};
\node[col2, align=center] at (3.9,-2.2){coordinée\\verbe en position 2};
\node[col3, align=center] at (8.6,-2.2){\all{Sie fordern Kameras,}\\\all{denn sie haben Angst.}\\« Ils demandent des caméras,\\car ils ont peur. »};
\node[col1] at (0,-4.4) {\all{deshalb}};
\node[col2, align=center] at (3.9,-4.4){principale\\inversion du sujet};
\node[col3, align=center] at (8.6,-4.4){\all{Sie haben Angst,}\\\all{deshalb fordern sie Kameras.}\\« Ils ont peur, c'est pourquoi\\ils demandent des caméras. »};
\end{tikzpicture}
\legende{Une même idée (peur $\rightarrow$ demande de caméras) exprimée par trois constructions différentes : seuls le connecteur et la place du verbe changent.}
\end{popfigure}

\courssub{Texte d'application : lire et repérer les connecteurs}

\begin{textebox}[Mehr Sicherheit in der Stadt]
Seit einigen Monaten gibt es in unserer Stadt mehr Vandalismus. Bushaltestellen werden beschädigt, Wände werden beschmiert, und nachts haben viele Bürger Angst. Weil die Gewalt zunimmt, reagieren die Politiker jetzt. Der Bürgermeister sagt: „Wegen der vielen Zwischenfälle verstärken wir die Polizeipräsenz. Wir installieren auch mehr Kameras, denn die Sicherheit der Bürger ist sehr wichtig.“ Viele Einwohner begrüßen diese Maßnahmen, deshalb gibt es jetzt mehr Kontrollen am Bahnhof und in den Schulen. Ein Lehrer erklärt: „Die Jugendlichen müssen verstehen, dass Vandalismus keine Lösung ist. Sie zerstören öffentliches Eigentum, so dass die Stadt viel Geld für Reparaturen ausgeben muss.“ Die Polizei organisiert deswegen auch Informationstage in den Schulen. Dort sprechen Polizisten mit den Schülern über Gewalt und ihre Folgen. Weil Prävention besser ist als Strafe, wollen die Schulen dieses Projekt fortsetzen. Die Bürger hoffen, dass ihre Stadt bald wieder sicher und sauber wird.
\end{textebox}

\textbf{Glossaire.} \all{beschmieren} : salir, couvrir de graffitis ; \all{der Zwischenfall, die Zwischenfälle} : l'incident ; \all{die Polizeipräsenz} : la présence policière ; \all{die Maßnahme, -n} : la mesure ; \all{begrüßen} : accueillir favorablement ; \all{das öffentliche Eigentum} : le bien public ; \all{die Reparatur, -en} : la réparation ; \all{die Strafe, -n} : la punition ; \all{die Prävention} : la prévention ; \all{fortsetzen} : continuer.

\textbf{Questions de compréhension.}
\begin{enumerate}
\item Relevez dans le texte tous les connecteurs de cause et de conséquence et classez-les : verbe à la fin, verbe en position 2, ou préposition.
\item \all{Warum verstärkt der Bürgermeister die Polizeipräsenz?} (Pourquoi le maire renforce-t-il la présence policière ?)
\item \all{Was passiert, weil die Jugendlichen öffentliches Eigentum zerstören?} (Que se passe-t-il parce que les jeunes détruisent le bien public ?)
\end{enumerate}

\begin{corrige}[Questions de compréhension]
\begin{enumerate}
\item Verbe à la fin : \all{weil} (\all{Weil die Gewalt zunimmt}), \all{so dass} (\all{so dass die Stadt viel Geld ausgeben muss}). Verbe en position 2 : \all{denn} (\all{denn die Sicherheit ist sehr wichtig}), \all{deshalb} (\all{deshalb gibt es mehr Kontrollen}), \all{deswegen} (\all{Die Polizei organisiert deswegen Informationstage}). Préposition : \all{wegen} (\all{Wegen der vielen Zwischenfälle}).
\item \all{Er verstärkt die Polizeipräsenz wegen der vielen Zwischenfälle.} / \all{Weil es mehr Vandalismus gibt, verstärkt er die Polizeipräsenz.}
\item \all{Die Stadt muss viel Geld für Reparaturen ausgeben.}
\end{enumerate}
\end{corrige}

\courssub{Méthode : relier deux idées par la cause ou la conséquence}

\begin{methode}[Bien choisir son connecteur en quatre étapes]
\begin{enumerate}
\item \textbf{Identifier la relation logique.} La deuxième idée est-elle la \emph{cause} de la première (réponse à « pourquoi ? ») ou sa \emph{conséquence} (réponse à « avec quel résultat ? ») ?
\item \textbf{Choisir le type de construction.} Cause avec un verbe conjugué : \all{weil} ou \all{denn}. Cause avec un simple nom : \all{wegen} + génitif. Conséquence : \all{deshalb/darum/deswegen} ou \all{so dass}.
\item \textbf{Placer le verbe correctement.} \all{weil} et \all{so dass} : verbe à la fin. \all{denn} et \all{deshalb} : verbe en position 2 (avec inversion après \all{deshalb}).
\item \textbf{Vérifier la phrase.} Relire en contrôlant : génitif après \all{wegen}, virgule avant \all{denn}/\all{so dass}, sujet après le verbe avec \all{deshalb}.
\end{enumerate}
\end{methode}

\begin{exoresolu}[Relier deux phrases]
Énoncé. Reliez les deux phrases de trois façons différentes (avec \all{weil}, avec \all{denn}, avec \all{deshalb}) : \all{Die Polizei verstärkt die Kontrollen. / Es gab einen Anschlag.}
\tcblower
Solution détaillée. La deuxième phrase (\all{Es gab einen Anschlag}) est la \emph{cause} ; la première est la \emph{conséquence}.
\begin{enumerate}
\item Avec \all{weil} (cause en subordonnée, verbe \all{gegeben hat} à la fin ; la subordonnée en tête provoque l'inversion dans la principale) :\\
\all{\textbf{Weil} es einen Anschlag \textbf{gegeben hat}, \textbf{verstärkt} die Polizei die Kontrollen.}\\
« Parce qu'il y a eu un attentat, la police renforce les contrôles. »
\item Avec \all{denn} (cause en coordinée, verbe en position 2) :\\
\all{Die Polizei \textbf{verstärkt} die Kontrollen, \textbf{denn} es \textbf{hat} einen Anschlag gegeben.}\\
« La police renforce les contrôles, car il y a eu un attentat. »
\item Avec \all{deshalb} (conséquence, inversion du sujet) :\\
\all{Es gab einen Anschlag, \textbf{deshalb verstärkt} die Polizei die Kontrollen.}\\
« Il y a eu un attentat, c'est pourquoi la police renforce les contrôles. »
\end{enumerate}
On pouvait aussi utiliser \all{wegen} + génitif : \all{\textbf{Wegen eines Anschlags} verstärkt die Polizei die Kontrollen.}
\end{exoresolu}

\begin{attention}[Les trois erreurs typiques du francophone]
\begin{enumerate}
\item Mettre le verbe en position 2 après \all{weil} : \emph{\all{weil sie haben Angst}} est faux. Correct : \all{weil sie Angst \textbf{haben}}.
\item Oublier l'inversion après \all{deshalb} : \emph{\all{deshalb die Bürger fordern}} est faux. Correct : \all{deshalb \textbf{fordern} die Bürger}.
\item Traduire « parce que » toujours par \all{weil} en tête de phrase : \all{denn} est souvent plus naturel quand la cause vient après, et \all{wegen} + nom permet une formulation plus courte et élégante.
\end{enumerate}
\end{attention}

\begin{exercice}[À vous de jouer]
\begin{enumerate}
\item Complétez avec \all{weil}, \all{denn}, \all{wegen} ou \all{deshalb} :\\
a) \all{............ der Anschläge bleiben viele Menschen zu Hause.}\\
b) \all{Die Schüler haben Angst, ............ wird ein Polizist an der Schule eingesetzt.}\\
c) \all{Der Bahnhof wurde geschlossen, ............ es dort einen Anschlag gab.}\\
d) \all{............ die Bushaltestelle beschädigt wurde, müssen die Schüler zu Fuß gehen.}
\item Traduisez en allemand : « Parce que la violence augmente, les citoyens demandent plus de sécurité. » (deux versions : \all{weil} et \all{denn})
\end{enumerate}
\end{exercice}

\begin{corrige}[Exercice : À vous de jouer]
\begin{enumerate}
\item a) \all{\textbf{Wegen} der Anschläge bleiben viele Menschen zu Hause.} (préposition + génitif pluriel \all{der})\\
b) \all{Die Schüler haben Angst, \textbf{deshalb wird} ein Polizist an der Schule eingesetzt.} (conséquence, inversion)\\
c) \all{Der Bahnhof wurde geschlossen, \textbf{denn} es gab dort einen Anschlag.} (coordinée, verbe en position 2)\\
d) \all{\textbf{Weil} die Bushaltestelle beschädigt wurde, müssen die Schüler zu Fuß gehen.} (subordonnée en tête, verbe \all{wurde} à la fin)
\item Avec \all{weil} : \all{\textbf{Weil} die Gewalt zunimmt, fordern die Bürger mehr Sicherheit.} Avec \all{denn} : \all{Die Bürger fordern mehr Sicherheit, \textbf{denn} die Gewalt nimmt zu.}
\end{enumerate}
\end{corrige}

\begin{aretenir}[L'essentiel de la section]
\begin{itemize}
\item Cause avec phrase : \all{weil} (verbe à la fin) ou \all{denn} (verbe en position 2).
\item Cause avec nom : \all{wegen} + \textbf{génitif} (\all{wegen des Terrors}, \all{wegen der Anschläge}).
\item Conséquence : \all{deshalb / darum / deswegen} + \textbf{inversion} du sujet ; \all{so dass} + verbe à la fin.
\item Règle d'or : le choix du connecteur détermine la place du verbe. Toujours vérifier cette place avant de rendre sa copie.
\end{itemize}
\end{aretenir}

\coursec{Méthode du commentaire et production guidée}

Après avoir maîtrisé le passif simple et les connecteurs de cause et de conséquence, nous disposons maintenant des outils grammaticaux indispensables pour structurer un commentaire de texte rigoureux et une production écrite argumentée. Le commentaire de texte est l'exercice roi du Baccalauréat : il évalue à la fois votre compréhension fine, votre capacité à synthétiser et votre aisance à vous exprimer en allemand correct. Cette section vous guide pas à pas, du texte source à la rédaction finale, en passant par l'analyse structurée.

\courssub{Les étapes du commentaire de texte au Bac}

Un commentaire réussi ne se contente pas de paraphraser le texte. Il suit une architecture en quatre temps, codifiée aux examens officiels. La transition entre l'analyse grammaticale (sections 4 et 5) et l'expression (cette section) est naturelle : \textbf{le passif permet de rapporter les faits objectivement} (ex. \all{„Die Scheiben wurden eingeschlagen.“}), \textbf{les connecteurs de cause/conséquence permettent d'articuler l'analyse} (ex. \all{„Weil die Jugendlichen gelangweilt sind, zerstören sie Eigentum.“}).

\begin{methode}[Les 4 étapes du commentaire de texte (Bac LV2)]
\begin{enumerate}
    \item \textbf{Présenter le texte (2 phrases maximum).} Indiquer le type de document (article de presse, rapport de police, témoignage, extrait de roman), la source (journal, auteur, date si connue), le thème général (le vandalisme, le terrorisme, la prévention) et l'intention de l'auteur (informer, alerter, dénoncer, proposer des solutions).
    \item \textbf{Résumer le contenu (3 à 5 phrases).} Restituer les idées principales \emph{avec ses propres mots} (reformulation), dans l'ordre logique du texte, sans copier-coller de longs segments. Utiliser des connecteurs (\all{zuerst, dann, schließlich, deshalb}).
    \item \textbf{Analyser les idées principales (2 à 3 idées).} Relever les arguments clés, les illustrer par de \emph{courtes citations} intégrées dans la phrase (entre guillemets allemands „…“), et expliquer leur portée. C'est ici qu'on mobilise le vocabulaire thématique et la grammaire (passif, cause/conséquence).
    \item \textbf{Donner son avis personnel argumenté (3 à 4 phrases).} Répondre à la question implicite : „Qu'en pensez-vous ?“. Prendre position, justifier avec des exemples (contexte camerounais, actualité allemande, expérience personnelle) et ouvrir sur une perspective (solution, prévention, rôle de l'école).
\end{enumerate}
\end{methode}

\begin{popfigure}
\begin{tikzpicture}[
    node distance=1.2cm and 1.5cm,
    box/.style={rectangle, rounded corners, draw=popblue, thick, fill=popblueL, text width=3.2cm, align=center, font=\small\bfseries, minimum height=1.8cm},
    arrow/.style={-Stealth, thick, popdark}
]
\node[box, align=center] (pres){1. PRÉSENTATION\\Type, source, thème,\\intention};
\node[box, right=of pres, align=center] (res){2. RÉSUMÉ\\Idées principales,\\reformulation,\\connecteurs};
\node[box, right=of res, align=center] (ana){3. ANALYSE\\Arguments, citations,\\vocabulaire, grammaire};
\node[box, right=of ana, align=center] (avis){4. AVIS PERSONNEL\\Position, arguments,\\exemples, ouverture};

\draw[arrow] (pres) -- (res) node[midway, above, font=\tiny, sloped] {structure};
\draw[arrow] (res) -- (ana) node[midway, above, font=\tiny, sloped] {approfondit};
\draw[arrow] (ana) -- (avis) node[midway, above, font=\tiny, sloped] {évalue};
\end{tikzpicture}
\legende{Organigramme de la structure du commentaire de texte au Baccalauréat}
\end{popfigure}

\courssub{Commentaire modèle : « Gewalt in der Stadt »}

Avant de rédiger, il faut un texte support. Voici un article de presse original, niveau A2/B1, traitant d'un fait de vandalisme dans une ville allemande moyenne, avec une dimension préventive.

\begin{textebox}[Gewalt in der Stadt — Artikel aus der „Rheinischen Post“, 12. März 2024]
In der nordrhein-westfälischen Stadt Hagen kam es am vergangenen Wochenende zu schweren Vandalismus. Eine Gruppe von etwa fünfzehn Jugendlichen zerstörte in der Innenstadt mehrere Bushaltestellen und parkte Fahrräder. Die Scheiben der Wartehäuschen wurden eingeschlagen, Graffiti an Hauswände gesprüht und Mülltonnen umgekippt. Der Sachschaden wird von der Polizei auf etwa 25.000 Euro geschätzt.

Anwohner alarmierten die Polizei in der Nacht zum Sonntag. Als die Beamten eintrafen, waren die Täter bereits geflohen. Dank Videoüberwachung konnten jedoch drei Personen identifiziert und vorläufig festgenommen werden. Sie sind zwischen 16 und 18 Jahre alt. Nach der Vernehmung wurden sie an ihre Eltern übergeben.

Der Bürgermeister der Stadt, Herr Müller, zeigte sich entsetzt: „Diese sinnlose Zerstörung trifft uns alle. Die Kosten muss die Allgemeinheit tragen.“ Er kündigte an, dass die Stadt die Videoüberwachung ausbauen und mehr Streetworker in den Vierteln einsetzen wolle, wo sich Jugendliche oft langweilen. „Prävention ist besser als Repression“, betonte er. Die Schulen sollen künftig stärker in Projekte gegen Gewalt eingebunden werden.
\end{textebox}

\begin{definition}[Glossaire : mots difficiles du texte]
\begin{itemize}
    \item \cle{die nordrhein-westfälische Stadt Hagen} : la ville de Hagen (Rhénanie-du-Nord-Westphalie).
    \item \cle{die Bushaltestelle, -n} : l'arrêt de bus.
    \item \cle{das Wartehäuschen, -} : l'abri d'attente / la salle d'attente.
    \item \cle{einschlagen (schlug ein, eingeschlagen)} : briser (vitre), enfoncer (porte).
    \item \cle{der Sachschaden} : les dégâts matériels.
    \item \cle{der Anwohner, - (die Anwohner)} : le riverain, l'habitant du quartier.
    \item \cle{vorläufig festnehmen} : interpeller / placer en garde à vue provisoirement.
    \item \cle{die Vernehmung, -en} : l'audition, l'interrogatoire.
    \item \cle{die Allgemeinheit} : la collectivité, le public.
    \item \cle{der Streetworker, -s} : l'éducateur de rue.
    \item \cle{die Prävention} : la prévention ; \cle{die Repression} : la répression.
\end{itemize}
\end{definition}

Voici un commentaire modèle respectant les quatre étapes, rédigé \textbf{en allemand} (langue de l'examen), d'une longueur idéale (environ 100 mots). La traduction française suit pour compréhension.

\begin{exemplebox}[Commentaire modèle / Modellkommentar (niveau attendu Bac)]
\textbf{Einleitung (Présentation) :} \all{Der Zeitungsartikel „Gewalt in der Stadt“ aus der „Rheinischen Post“ vom 12. März 2024 berichtet über Vandalismus in Hagen. Der Autor informiert über die Taten, die Polizei und die Reaktion der Politik.}

\textbf{Zusammenfassung (Résumé) :} \all{In der Nacht zum Sonntag zerstörte eine Jugendgruppe Bushaltestellen und Fahrräder. Der Schaden beträgt 25.000 Euro. Die Polizei nahm dank Videoüberwachung drei Täter fest. Der Bürgermeister verurteilt die Taten und kündigt mehr Kameras und Streetworker an.}

\textbf{Analyse (Analyse) :} \all{Ein zentrales Argument ist die Langeweile der Jugendlichen als Ursache: „wo sich Jugendliche oft langweilen“. Der Text nutzt Passivkonstruktionen („wurden eingeschlagen“, „wurde geschätzt“), um den Schaden objektiv darzustellen. Die Lösung verbindet Repression und Prävention.}

\textbf{Stellungnahme (Avis personnel) :} \all{Meiner Meinung nach ist der Ansatz richtig, weil er Ursachen bekämpft. In Kamerun gibt es ähnliche Probleme in Yaoundé oder Douala. Deshalb sind Jugendarbeit und Bildung wichtig, damit keine Gewalt entsteht.}
\end{exemplebox}

\begin{exemplebox}[Banque d'expressions pour le commentaire (à mémoriser)]
\begin{center}
\begin{tabularx}{\linewidth}{|X|X|}\hline
\rowcolor{popblueL} \textbf{Fonction} & \textbf{Formules types (Deutsch)} \\ \hline
Présenter le texte & \all{Der Text handelt von ...} (Le texte traite de ...) \\
& \all{Es geht um ...} (Il s'agit de ...) \\
& \all{Der Artikel wurde in ... veröffentlicht.} (L'article a été publié dans ...) \\
& \all{Der Autor informiert über / warnt vor / kritisiert ...} (L'auteur informe sur / met en garde contre / critique ...) \\ \hline
Résumer & \all{Zuerst wird berichtet, dass ...} (D'abord il est rapporté que ...) \\
& \all{Anschließend / Danach ...} (Ensuite / Après ...) \\
& \all{Schließlich stellt der Autor fest, dass ...} (Finalement l'auteur constate que ...) \\ \hline
Analyser / Citer & \all{Ein zentrales Argument ist ...} (Un argument central est ...) \\
& \all{Der Autor schreibt: „...“} (L'auteur écrit : „...“) \\
& \all{Dies zeigt, dass ...} (Cela montre que ...) \\
& \all{Aufgrund von ... / Deshalb ...} (En raison de ... / C'est pourquoi ...) \\ \hline
Donner son avis & \all{Meiner Meinung nach ...} (À mon avis ...) \\
& \all{Ich finde, dass ... / Ich meine, dass ...} (Je trouve que / Je pense que ...) \\
& \all{Man sollte ... / Es ist wichtig, dass ...} (On devrait / Il est important que ...) \\
& \all{In Kamerun ist die Situation ähnlich / anders, weil ...} (Au Cameroun la situation est similaire / différente, car ...) \\ \hline
\end{tabularx}
\end{center}
\end{exemplebox}

\courssub{Production guidée : rédiger un paragraphe sur un fait de vandalisme}

\emph{Consigne :} Rédigez un paragraphe cohérent de 5 à 6 phrases en allemand. Décrivez un fait de vandalisme réel ou imaginé dans votre ville (Yaoundé, Douala, Bafoussam, etc.). \textbf{Contraintes obligatoires :}
\begin{enumerate}
    \item Utiliser \textbf{au moins 2 formes du passif} (Présent ou Prétérit : \all{werden / wurden + Partizip II}).
    \item Utiliser \textbf{au moins 1 connecteur de cause ou de conséquence} (\all{weil, da, deshalb, daher, sodass}).
    \item Employer au moins \textbf{5 mots du vocabulaire thématique} (voir Section 3 : \all{der Vandalismus, die Gewalt, beschädigen, die Polizei, die Sicherheit, der Anschlag, die Zerstörung, der Täter, das Opfer, verhindern}).
    \item Respecter l'ordre des mots (verbe en position 2, verbe à la fin en subordonnée).
\end{enumerate}

\begin{exoresolu}[Modèle de production guidée avec analyse]
\textbf{Énoncé :} Rédiger le paragraphe demandé ci-dessus.

\textbf{Proposition de rédaction :}
\all{Letzten Samstag wurde in Douala ein neues Jugendzentrum Opfer von Vandalismus. Die Fensterscheiben wurden beschädigt und an die Wände wurden Graffiti gesprüht. Weil die Sicherheitskameras defekt waren, konnte die Polizei die Täter nicht sofort identifizieren. Deshalb wurde die Zerstörung erst am Montag bemerkt. Der Bürgermeister fordert jetzt mehr Sicherheit und bessere Prävention in den Vierteln.}

\textbf{Analyse grammaticale (vérification des contraintes) :}
\begin{itemize}
    \item \textbf{Passif 1 (Prétérit) :} \all{wurde ... Opfer} (tournure passive « devenir victime »).
    \item \textbf{Passif 2 (Prétérit) :} \all{wurden ... beschädigt} (ont été endommagées — \all{beschädigen}).
    \item \textbf{Passif 3 (Prétérit) :} \all{wurden ... gesprüht} (ont été tagués — \all{sprühen}).
    \item \textbf{Passif 4 (Prétérit) :} \all{wurde ... bemerkt} (a été remarquée — \all{bemerken}).
    \item \textbf{Connecteur de cause :} \all{Weil die Sicherheitskameras defekt waren,} (verbe à la fin !).
    \item \textbf{Connecteur de conséquence :} \all{Deshalb wurde ...} (verbe en position 2 !).
    \item \textbf{Vocabulaire thématique (8 mots) :} \all{Vandalismus, Opfer, beschädigen, Polizei, Täter, Sicherheit, Zerstörung, Prävention}.
\end{itemize}
\textbf{Traduction :} Samedi dernier, un nouveau centre pour jeunes à Douala a été victime de vandalisme. Les vitres ont été endommagées et les murs ont été tagués. Comme les caméras de surveillance étaient en panne, la police n'a pas pu identifier les auteurs immédiatement. C'est pourquoi la destruction n'a été constatée que lundi. Le maire réclame maintenant plus de sécurité et une meilleure prévention dans les quartiers.
\end{exoresolu}

\begin{attention}[Pièges à éviter pour les francophones]
\begin{itemize}
    \item \textbf{Ordre des mots avec \all{weil / da} :} Ne JAMAIS écrire \all{*Weil die Kameras waren defekt}. \rightarrow \all{Weil die Kameras defekt \textbf{waren}.} (Verbe À LA FIN).
    \item \textbf{Ordre des mots avec \all{deshalb / daher / deshalb} :} Ce sont des adverbes conjonctifs, ils occupent la \textbf{place 1}. \rightarrow \all{Deshalb \textbf{wurde} die Zerstörung bemerkt.} (Verbe en position 2).
    \item \textbf{Passif : auxiliaire \all{werden} + Partizip II.} Ne pas confondre avec \all{sein} (état) ou \all{haben} (actif). \all{Die Scheibe ist zerbrochen} (état, \all{sein}) vs \all{Die Scheibe wird zerbrochen} (action, \all{werden}).
    \item \textbf{Genre des noms :} \all{der Vandalismus, der Terrorismus, der Anschlag, die Gewalt, die Sicherheit, die Zerstörung, die Prävention}. Apprenez l'article \emph{avec} le nom.
\end{itemize}
\end{attention}

\courssub{Grille d'auto-évaluation et correction collective}

Avant de rendre votre copie, utilisez cette grille pour vous auto-corriger. C'est une compétence transversale (méthodologie) évaluée au Bac.

\begin{definition}[Grille d'auto-évaluation — Production écrite (Vandalisme / Terrorisme)]
\begin{center}
\begin{tabularx}{\linewidth}{|>{\bfseries}X|X|X|X|}\hline
\rowcolor{popblueL} \textbf{Critère} & \textbf{Non acquis (0)} & \textbf{En cours (1)} & \textbf{Acquis (2)} \\ \hline
Vocabulaire thématique (5 mots min.) & 0--1 mot utilisé & 2--4 mots utilisés & 5+ mots utilisés correctement (article + pluriel connu) \\ \hline
Passif simple (2 occurrences min.) & Absent ou mal formé (ex. *ist beschädigt) & 1 passif correct ou 2 avec erreurs mineures & 2 passifs corrects (Présent/Prétérit, \all{werden} + Partizip II) \\ \hline
Connecteurs cause/conséquence (1 min.) & Absent & Présent mais ordre des mots faux (ex. *weil verb pos 2) & Correct : \all{weil/da} (verbe fin) OU \all{deshalb/daher} (verbe pos 2) \\ \hline
Ordre des mots (Hauptsatz / Nebensatz) & Erreurs fréquentes & Quelques erreurs & Respecté partout (V2 / Verbe final) \\ \hline
Orthographe / Majuscules (Noms) & Beaucoup d'erreurs & Quelques oublis de majuscules & Parfait : tous les noms en majuscule, Umlauts/ß corrects \\ \hline
Cohérence / Contenu (5-6 phrases) & Incompréhensible / HS & Compréhensible mais décousu & Logique, fluide, répond à la consigne \\ \hline
\end{tabularx}
\end{center}
\textbf{Score :} /12. Objectif Bac : $\ge$ 10/12.
\end{definition}

\begin{experience}[Correction collective en classe (Activité orale guidée)]
\textbf{Déroulement (15 min) :}
\begin{enumerate}
    \item \textbf{Échange de copies :} Chaque élève donne sa production à son voisin.
    \item \textbf{Lecture à voix haute :} Le correcteur lit le texte de son camarade à haute voix (entraîne la phonétique).
    \item \textbf{Repérage ciblé :} Le correcteur surligne au stylo vert : (1) les passifs, (2) les connecteurs, (3) le vocabulaire thème. Il entoure au rouge les erreurs d'ordre des mots ou de majuscules.
    \item \textbf{Feedback oral structuré :} « Ton passif \all{wurden gesprüht} est correct. Par contre, attention : \all{*weil die Polizei kam} $\rightarrow$ \all{weil die Polizei \textbf{kam}} (verbe à la fin). Tu as utilisé 4 mots du thème, essaie d'en mettre un 5e. »
    \item \textbf{Rédaction d'une version améliorée :} L'auteur réécrit son paragraphe en intégrant les corrections (5 min).
    \item \textbf{Mise en commun :} 2-3 volontaires lisent leur version finale au tableau ; la classe valide avec la grille.
\end{enumerate}
\end{experience}

\begin{savaistu}[L'astuce du Bac : l'avis personnel rapporte gros]
Au Baccalauréat camerounais (LV2 Allemand), la partie « Avis personnel / Expression libre » vaut souvent \textbf{4 à 6 points sur 20} à elle seule. Beaucoup d'élèves la bâclent par manque de temps ou d'idées.
\cle{Stratégie gagnante :} Préparez \textbf{2 à 3 « blocs d'opinion » prêts à l'emploi} (templates) que vous adaptez au sujet du jour.
\begin{itemize}
    \item \textbf{Bloc 1 (Rôle de l'école/éducation) :} \all{Meiner Meinung nach muss die Schule mehr Prävention machen. Jugendliche brauchen Perspektiven, damit sie nicht gewalttätig werden.}
    \item \textbf{Bloc 2 (Rôle de la famille/société) :} \all{Ich finde, dass die Eltern mehr Verantwortung tragen müssen. Auch die Gesellschaft muss hinsehen, nicht wegsehen.}
    \item \textbf{Bloc 3 (Comparaison Cameroun/Allemagne) :} \all{In Kamerun sehen wir ähnliche Probleme. Die Arbeitslosigkeit der Jugend führt oft zu Gewalt. Wir brauchen mehr Ausbildungsplätze.}
\end{itemize}
Apprenez ces phrases par cœur (prononciation comprise). Le jour J, vous n'aurez plus qu'à les insérer à l'étape 4. Cela vous garantit des points faciles et montre une maîtrise syntaxique (subordonnées \all{dass/weil}, modaux \all{müssen/sollen}, passif \all{werden}).
\end{savaistu}

\begin{exercice}[Entraînement final : Commentaire complet]
Sur la base du texte \all{„Gewalt in der Stadt“} (page 2), rédigez un commentaire complet structuré en 4 paragraphes (Présentation, Résumé, Analyse, Avis). Longueur visée : 120--150 mots. Utilisez la banque d'expressions (page 3) et la grille d'auto-évaluation (page 5).
\end{exercice}


\begin{corrige}[Exercice : Commentaire complet]
\textbf{Einleitung :} \all{Der Zeitungsartikel „Gewalt in der Stadt“ aus der „Rheinischen Post“ vom 12. März 2024 informiert über Vandalismus in Hagen und die Reaktion der Politik.}

\textbf{Zusammenfassung :} \all{Zuerst wird beschrieben, wie eine Jugendgruppe Bushaltestellen und Fahrräder zerstörte. Der Schaden beträgt 25.000 Euro. Danach wird berichtet, dass die Polizei dank Videos drei Täter festnahm. Schließlich erklärt der Bürgermeister, dass die Stadt mehr Kameras und Streetworker einsetzen will.}

\textbf{Analyse :} \all{Ein Hauptargument ist die Langeweile der Jugendlichen als Ursache: „wo sich Jugendliche oft langweilen“. Der Autor zeigt durch Passivsätze („wurden eingeschlagen“, „wurde geschätzt“), dass die Täter erst sekundär wichtig sind, der Schaden primär. Die Lösung verbindet Repression (Polizei, Video) und Prävention (Streetworker, Schule).}

\textbf{Stellungnahme :} \all{Meiner Meinung nach ist der Ansatz des Bürgermeisters richtig. Man muss die Ursachen bekämpfen: Jugendarbeit und Bildung. In Kamerun gibt es ähnliche Probleme in Yaoundé oder Douala. Deshalb sind Projekte für die Jugend sehr wichtig, damit sie keine Gewalt anwenden.}
\end{corrige}

\newpage

\coursec{Activité d'intégration}

\coursec{Le vandalisme et le terrorisme}

\courssub{Objectifs de l'activité}

À la fin de cette activité d'intégration, tu seras capable de :

\begin{itemize}
\item mobiliser dans une production écrite authentique le \cle{vocabulaire} du thème « le vandalisme et le terrorisme » (\all{der Vandalismus, der Terror, der Anschlag, die Gewalt, beschädigen, zerstören, die Polizei, die Sicherheit, evakuieren, der Bahnhof}) ;
\item employer correctement le \cle{passif simple} (\all{Vorgangspassiv}) au \emph{Présent}, au \emph{Präteritum} et au \emph{Perfekt} pour rapporter des faits sans nommer les auteurs ;
\item articuler ton texte avec les \cle{connecteurs de cause et de conséquence} (\all{weil, wegen, deshalb, darum}) ;
\item rédiger en groupe un court \cle{article de presse} structuré (titre, introduction, faits, réaction, conclusion) ;
\item lire ton article à voix haute, répondre aux questions de la classe et évaluer les productions des autres groupes selon des critères précis.
\end{itemize}

\courssub{Situation}

Tu es journaliste stagiaire au journal de ton lycée, \emph{Die Schülerzeitung}. Hier soir, un incident s'est produit dans ta ville : des abribus ont été dégradés près du marché, ou bien un colis suspect a provoqué une alerte de sécurité à la gare. Le rédacteur en chef te demande, avec ton équipe, un court article pour la page « Faits divers » de demain. L'article doit informer les lecteurs clairement, sans exagération, et respecter les règles de la langue allemande.

Voici le modèle de dépêche que le rédacteur en chef vous distribue comme exemple de style journalistique :

\begin{textebox}[Modèle de dépêche — \all{Kurzmeldung}]
\all{Am Montagabend wurde in der Innenstadt eine Bushaltestelle beschädigt. Die Scheiben wurden zerstört und die Sitzbänke wurden beschmiert. Die Polizei wurde sofort informiert. Wegen des Vorfalls wurde der Verkehr eine Stunde lang gestoppt. Ein Zeuge wurde befragt, aber der Täter wurde noch nicht gefunden. Die Sicherheit in der Stadt muss verbessert werden, sagt der Bürgermeister.}\par
\emph{Traduction : Lundi soir, un arrêt de bus a été dégradé au centre-ville. Les vitres ont été détruites et les bancs ont été couverts de graffitis. La police a été informée immédiatement. À cause de l'incident, la circulation a été interrompue pendant une heure. Un témoin a été interrogé, mais le coupable n'a pas encore été retrouvé. La sécurité en ville doit être améliorée, dit le maire.}
\end{textebox}

\begin{savaistu}[La presse locale en Allemagne]
En Allemagne, presque chaque ville possède son journal local (\all{die Lokalzeitung, -en}). Les articles de faits divers (\all{die Kurzmeldung, -en}) utilisent très souvent le passif, car on ignore qui a commis l'acte : \all{Die Bank wurde überfallen} (la banque a été attaquée). Au Cameroun aussi, les journaux comme \emph{Cameroon Tribune} rapportent chaque jour des incidents à Yaoundé ou à Douala : à toi de faire le même travail... en allemand !
\end{savaistu}

\courssub{Consignes}

Travail en groupes de 3 ou 4 élèves. Durée totale : 50 minutes. Suivez les étapes dans l'ordre.

\begin{enumerate}
\item \textbf{Choix du sujet (5 min).} Choisissez ensemble UN des deux sujets : (a) un acte de vandalisme imaginaire (abribus dégradé, mur d'école tagué, lampadaires cassés au marché de Mokolo ou de Deïdo) ; (b) une alerte de sécurité imaginaire (colis suspect à la gare, sac oublié au stade). Notez : \emph{qui ? quoi ? où ? quand ? conséquences ?}
\item \textbf{Collecte du lexique (5 min).} Listez au moins 8 mots du vocabulaire thématique de la leçon que vous pourrez utiliser, avec leur article et leur pluriel.
\item \textbf{Plan de l'article (5 min).} Organisez vos idées en quatre parties : titre accrocheur ; phrase d'introduction (où et quand ?) ; description des faits au passif ; réaction (police, maire, habitants) et conclusion.
\item \textbf{Rédaction (20 min).} Rédigez l'article \all{Vorfall in unserer Stadt} en 10 à 12 lignes. Contraintes obligatoires :
\begin{itemize}
\item au moins \textbf{4 phrases au passif} (présent, \emph{Präteritum} ou \emph{Perfekt}) ;
\item au moins \textbf{5 mots du lexique thématique} de la leçon ;
\item au moins \textbf{1 phrase avec \all{weil} ou \all{wegen}} (cause) ;
\item au moins \textbf{1 phrase avec \all{deshalb}} (conséquence).
\end{itemize}
\item \textbf{Relecture (5 min).} Vérifiez : majuscules aux noms, place du verbe conjugué en deuxième position, verbe à la fin après \all{weil}, inversion du sujet après \all{deshalb}, accord de \all{werden} au passif, génitif après \all{wegen}.
\item \textbf{Lecture publique et questions (10 min).} Un membre du groupe lit l'article à la classe. Les autres élèves posent 2 questions de compréhension (par exemple : \all{Wann ist der Vorfall passiert? Wer wurde informiert?}).
\item \textbf{Vote.} La classe vote pour l'article le plus clair et le mieux construit, selon la grille : vocabulaire (2 pts), passif correct (2 pts), connecteurs (2 pts), clarté et prononciation (2 pts).
\end{enumerate}

\courssub{Ressources à mobiliser}

\begin{aretenir}[Boîte à outils pour l'article]
\textbf{Lexique (noms, verbes, articles, pluriels) :}\par
\all{der Vandalismus} (le vandalisme), \all{der Terror} (le terrorisme), \all{der Anschlag, die Anschläge} (l'attentat), \all{die Gewalt} (la violence), \all{beschädigen} (endommager), \all{zerstören} (détruire), \all{evakuieren} (évacuer), \all{der Bahnhof, die Bahnhöfe} (la gare), \all{die Polizei} (la police), \all{die Sicherheit} (la sécurité), \all{der Zeuge, -n} (le témoin), \all{der Täter, die Täter} (l'auteur, le coupable), \all{informieren} (informer), \all{festnehmen} (arrêter), \all{der Verdacht, die Verdächte} (le soupçon), \all{der Sprengsatz, die Sprengsätze} (l'engin explosif).\par
\textbf{Passif Présent :} \all{werden + Partizip II} → \all{Die Scheibe wird zerstört.} (La vitre est détruite.)\par
\textbf{Passif Präteritum :} \all{wurden + Partizip II} → \all{Die Scheiben wurden zerstört.} (Les vitres ont été détruites — temps du récit journalistique.)\par
\textbf{Passif Perfekt :} \all{sein + Partizip II + worden} → \all{Die Scheibe ist zerstört worden.} (La vitre a été détruite.)\par
\textbf{Cause :} \all{weil + verbe à la fin} → \all{..., weil der Täter nicht gefunden wurde.} / \all{wegen + génitif} → \all{Wegen des Vorfalls...}\par
\textbf{Conséquence :} \all{deshalb + verbe en 2e position} → \all{Deshalb wurde die Polizei informiert.} / \all{darum + verbe en 2e position} → \all{Darum patrouilliert die Polizei jetzt.}
\end{aretenir}

\begin{attention}[Pièges à éviter]
\begin{itemize}
\item Ne dis jamais \all{Die Scheibe wurde zerstört geworden} : au \emph{Perfekt} passif, on écrit \all{worden} (sans \all{ge-}).
\item Après \all{deshalb} ou \all{darum}, le verbe conjugué reste en deuxième position : \all{Deshalb kam die Police.} et non \all{Deshalb die Police kam.}
\item Après \all{weil}, le verbe va à la fin : \all{..., weil viele Menschen Angst hatten.}
\item \all{wegen} + génitif (écrit) / datif (oral) : \all{wegen des Vorfalls}, pas \all{wegen der Vorfall}.
\item \all{der Täter} (masculin) → pluriel \all{die Täter} (Umlaut) ; \all{der Zeuge} (faible) → \all{den Zeugen} (accusatif/datif).
\item Majuscule obligatoire à \textbf{tous} les noms : \all{die Sicherheit}, \all{der Vandalismus}, \all{der Bahnhof}.
\end{itemize}
\end{attention}

\courssub{Proposition de production et corrigé commenté}

\begin{exoresolu}[Article modèle : \all{Vorfall in unserer Stadt}]
\textbf{Énoncé.} Rédigez un article de 10 à 12 lignes respectant les quatre contraintes de la consigne.\par
\tcblower
\textbf{Production modèle :}\par
\all{\textbf{Vorfall in unserer Stadt}}\par
\all{Am Dienstagabend wurde in Yaoundé, nahe dem Markt Mokolo, eine Bushaltestelle beschädigt. Die Scheiben wurden zerstört und zwei Sitzbänke wurden beschmiert. Ein Zeuge hat die Polizei sofort informiert. Wegen des Vorfalls wurde der Verkehr zwei Stunden lang gestoppt. Viele Menschen hatten Angst, weil ein solcher Akt von Vandalismus neu in diesem Viertel ist. Deshalb wurde die Sicherheit am Markt verstärkt: Jetzt patrouillieren zwei Polizisten jeden Abend. Der Täter wurde noch nicht gefunden, aber die Polizei sucht weiter. Der Bürgermeister sagt: „Die Gewalt muss gestoppt werden. Unsere Stadt muss sicher bleiben."}\par
\emph{Traduction : Mardi soir, à Yaoundé, près du marché Mokolo, un abribus a été dégradé. Les vitres ont été détruites et deux bancs ont été couverts de graffitis. Un témoin a informé la police immédiatement. À cause de l'incident, la circulation a été interrompue pendant deux heures. Beaucoup de gens ont eu peur, parce qu'un tel acte de vandalisme est nouveau dans ce quartier. C'est pourquoi la sécurité au marché a été renforcée : désormais, deux policiers patrouillent chaque soir. Le coupable n'a pas encore été retrouvé, mais la police continue ses recherches. Le maire dit : « La violence doit être stoppée. Notre ville doit rester sûre. »}\par
\textbf{Commentaire des choix de langue :}
\begin{itemize}
\item \textbf{Passif (6 occurrences) :} \all{wurde beschädigt} (Präteritum), \all{wurden zerstört} (Präteritum), \all{wurden beschmiert} (Präteritum), \all{wurde gestoppt} (Präteritum), \all{wurde verstärkt} (Präteritum), \all{wurde nicht gefunden} (Präteritum) ; \all{muss ... gestoppt werden} (Passif modal présent). Le \emph{Präteritum} est le temps typique du reportage écrit ; il permet de ne pas nommer l'auteur inconnu.
\item \textbf{Lexique thématique (7 mots) :} \all{beschädigen, zerstören, die Polizei, der Vandalismus, die Sicherheit, die Gewalt, der Täter}.
\item \textbf{Cause :} \all{weil ein solcher Akt ... neu ist} (verbe \all{ist} en fin de subordonnée) et \all{Wegen des Vorfalls} (génitif correct : \all{des} = \all{von dem}).
\item \textbf{Conséquence :} \all{Deshalb wurde die Sicherheit verstärkt} (inversion correcte : verbe \all{wurde} en position 2, sujet \all{die Sicherheit} en position 3).
\item \textbf{Précision contextuelle :} \all{nahe dem Markt Mokolo} (Datif après \all{nahe}), \all{Bushaltestelle} (terme allemand standard pour « abribus / arrêt de bus »).
\item \textbf{Style journalistique :} titre nominal, faits précis (où, quand, quoi), citation directe du maire entre guillemets allemands „…“, conclusion ouverte sur l'avenir.
\end{itemize}
\end{exoresolu}

\courssub{Bilan de l'activité}

Cette activité t'a permis de réinvestir l'ensemble des savoir-faire de la leçon dans une tâche authentique et motivante :

\begin{itemize}
\item \textbf{Compréhension :} tu as lu et analysé un modèle de dépêche avant de produire, et tu as écouté les articles des autres groupes pour poser des questions.
\item \textbf{Lexique :} tu as utilisé au moins cinq mots du champ lexical du vandalisme, du terrorisme et de la sécurité, avec leurs articles et pluriels.
\item \textbf{Grammaire :} tu as conjugué le passif au \emph{Präteritum} (temps du récit), au \emph{Perfekt} et au \emph{Présent} (modal), et tu as articulé ton texte avec \all{weil}, \all{wegen} et \all{deshalb}/\all{darum} en respectant la place du verbe.
\item \textbf{Expression écrite et orale :} tu as rédigé un article structuré, relu ton texte, lu à voix haute et répondu aux questions de la classe.
\end{itemize}

\begin{methode}[Pour réussir ce type de production à l'avenir]
\begin{enumerate}
\item Lis d'abord un modèle et souligne les formes utiles (passif, connecteurs, déclinaisons).
\item Fais un plan simple : titre → introduction (où/quand) → faits (passif) → réaction → conclusion.
\item Écris des phrases courtes ; vérifie chaque verbe (place, temps, auxiliaire \all{werden}/\all{sein}, participe passé).
\item Relis à voix haute : si une phrase sonne faux, corrige-la (ordre des mots, cas après préposition).
\item Garde une liste personnelle des erreurs corrigées (ex. \all{wegen + Genitiv}, \all{worden} sans \all{ge-}) pour ne pas les refaire.
\end{enumerate}
\end{methode}

Le vote final et la correction collective au tableau permettront de retenir, pour toute la classe, les meilleures formulations : recopie dans ton cahier les deux ou trois phrases modèles élues par le groupe.

\newpage

\coursec{Méthodes et exercices résolus}

\coursec{Le vandalisme et le terrorisme}

\courssub{Texte support : Vandalisme dans une école de Yaoundé}

\begin{textebox}[Vandalismus in der Schule]
Am Montagmorgen fanden die Schüler des Gymnasiums in Yaoundé eine böse Überraschung: In der Nacht hatten unbekannte Personen zehn Fenster zerstört und die Wände der Bibliothek mit Farbe beschmiert. Auch zwei Computer wurden beschädigt. Der Direktor informierte sofort die Polizei. „Diese \cle{Gewalt} ist ein großes Problem für unsere Schule“, sagte er. „Die \cle{Reparatur} wird viel Geld kosten.“ Die \cle{Polizei} sucht jetzt nach den \cle{Tätern}. Viele Schüler sind schockiert. „Wir brauchen mehr \cle{Sicherheit} in unserer Schule“, meint eine Schülerin. Die Lehrer wollen mit den Schülern über \cle{Respekt} und \cle{Zusammenleben} diskutieren.
\end{textebox}

\begin{definition}[Mots-clés du texte (avec article et pluriel)]
\begin{itemize}
\item \all{der Anschlag, die Anschläge} — attentat, acte de violence ciblé.
\item \all{die Gewalt} (singulier seulement) — la violence.
\item \all{beschädigen} (verbe régulier : beschädigte, hat beschädigt) — endommager, abîmer (+ accusatif).
\item \all{zerstören} (verbe régulier : zerstörte, hat zerstört) — détruire (+ accusatif).
\item \all{die Polizei} (singulier seulement) — la police (verbe au singulier !).
\item \all{der Täter, die Täter} — l'auteur, le coupable (masculin) ; féminin : \all{die Täterin, die Täterinnen}.
\item \all{das Opfer, die Opfer} — la victime.
\item \all{die Sicherheit} — la sécurité, la sûreté.
\item \all{suchen nach + Dativ} — chercher, rechercher.
\item \all{informieren + Akkusativ} — informer.
\end{itemize}
\end{definition}

\courssub{Méthode 1 : Comprendre un texte sur le vandalisme et le terrorisme}

\begin{methode}[Lire et comprendre un texte argumentatif / factuel]
\begin{enumerate}
\item \textbf{Lire le titre et la source} : ils annoncent le thème (fait divers, article de presse). Souligner les mots-clés : \cle{der Anschlag}, \cle{die Gewalt}, \cle{beschädigen}, \cle{der Täter}.
\item \textbf{Première lecture globale} : repérer les \cle{5 W} — \all{wer?} (qui), \all{was?} (quoi), \all{wo?} (où), \all{wann?} (quand), \all{warum?} (pourquoi).
\item \textbf{Deuxième lecture détaillée} : distinguer les \cle{Tatsachen} (faits vérifiables) des \cle{Meinungen} (opinions, sentiments). Repérer les chiffres, lieux, noms propres.
\item \textbf{Deviner les mots inconnus} par le contexte et les familles de mots : \all{der Terror} $\rightarrow$ \all{der Terrorismus, der Terrorist, terroristisch} ; \all{die Gewalt} $\rightarrow$ \all{gewalttätig, die Gewalttat} ; \all{sicher} $\rightarrow$ \all{die Sicherheit, sicherlich, sichern}.
\item \textbf{Répondre aux questions} en reformulant avec ses propres mots, en phrases \textbf{complètes} (sujet + verbe conjugué en 2\up{e} position), jamais en recopiant le texte mot à mot.
\item \textbf{Vérifier} : majuscules à \textbf{tous} les noms, place du verbe (2\up{e} position en principale, fin en subordonnée), orthographe des mots du texte (Umlaute, ß).
\end{enumerate}
\end{methode}

\begin{exoresolu}[Compréhension de texte : un acte de vandalisme]
Lisez le texte ci-dessus puis répondez en allemand par des phrases complètes.

\textbf{Questions :}
1. Was haben die unbekannten Personen in der Nacht gemacht?
2. Wann hat der Direktor die Polizei informiert?
3. Was meint die Schülerin zur Situation?
4. Was wollen die Lehrer jetzt tun?

\tcblower
\textbf{Raisonnement (repérage des 5 W) :}
\emph{Wer} = unbekannte Personen ; \emph{Was} = 10 Fenster zerstört, Wände beschmiert, 2 Computer beschädigt ; \emph{Wo} = Gymnasium in Yaoundé ; \emph{Wann} = in der Nacht auf Montag / Montagmorgen ; \emph{Warum} (opinion) = Respekt fehlt, mehr Sicherheit nötig.

\textbf{Réponses rédigées :}

1. \all{Die unbekannten Personen haben zehn Fenster zerstört, die Wände der Bibliothek mit Farbe beschmiert und zwei Computer beschädigt.} « Les personnes inconnues ont détruit dix fenêtres, barbouillé les murs de la bibliothèque avec de la peinture et endommagé deux ordinateurs. » — On reprend les trois faits, verbe \all{haben} (2\up{e} pos.) + participes à la fin (Perfekt).

2. \all{Der Direktor hat die Polizei sofort informiert.} « Le directeur a informé la police immédiatement. » — \all{sofort} (adverbe de temps) en position 3, après le sujet.

3. \all{Die Schülerin meint, dass sie mehr Sicherheit in der Schule brauchen.} « L'élève pense qu'ils/elles ont besoin de plus de sécurité dans l'école. » — Subordonnée avec \all{dass} $\rightarrow$ verbe \all{brauchen} à la fin. \all{sie} reprend \all{Wir} du texte (discours indirect).

4. \all{Die Lehrer wollen mit den Schülern über Respekt und Zusammenleben diskutieren.} « Les professeurs veulent discuter avec les élèves du respect et de la vie commune. » — Infinitif \all{diskutieren} à la fin (verbe modal \all{wollen} en 2\up{e} pos.).

\textbf{Conclusion :} Chaque réponse est une phrase autonome, verbe conjugué en 2\up{e} position (principale) ou en fin (subordonnée \all{dass}), noms avec majuscule.
\end{exoresolu}

\begin{exercice}[Compréhension : Entraînement personnel]
Répondez en allemand (phrases complètes) :
1. Wie viele Fenster wurden zerstört?
2. Wer hat die Wände beschmiert?
3. Warum ist der Direktor schockiert? (Formulez une hypothèse avec \all{weil}.)
4. Trouvez dans le texte le synonyme de \all{reparieren} (réparer) au nom.
\end{exercice}

\begin{corrige}[Exercice Compréhension]
1. \all{Zehn Fenster wurden zerstört.} (Passif Präteritum)
2. \all{Unbekannte Personen haben die Wände beschmiert.} (Actif Perfekt)
3. \all{Der Direktor ist schockiert, weil die Schäden sehr groß sind.} (Cause \all{weil} + verbe final)
4. \all{die Reparatur} (nom féminin, majuscule).
\end{corrige}

\courssub{Lexique thématique : Le vandalisme et le terrorisme}

\begin{methode}[Construire et utiliser le lexique du thème]
\begin{enumerate}
\item \textbf{Apprendre chaque nom avec son article et son pluriel} :
\begin{center}
\begin{tabular}{|l|l|l|}
\hline
\rowcolor{popblueL} Singulier & Pluriel & Français \\
\hline
\cle{der Terror} & — & le terrorisme (concept) \\
\cle{der Anschlag} & \cle{die Anschläge} & l'attentat \\
\cle{die Gewalt} & — & la violence \\
\cle{der Vandalismus} & — & le vandalisme \\
\cle{die Polizei} & — & la police \\
\cle{die Sicherheit} & — & la sécurité \\
\cle{der Täter} & \cle{die Täter} & l'auteur (m.) \\
\cle{das Opfer} & \cle{die Opfer} & la victime \\
\cle{der Vandale} & \cle{die Vandalen} & le vandale \\
\cle{die Reparatur} & \cle{die Reparaturen} & la réparation \\
\hline
\end{tabular}
\end{center}
\item \textbf{Regrouper par familles de mots (Wortfamilien)} :
\begin{itemize}
\item \all{Terror} $\rightarrow$ \all{der Terrorismus, der Terrorist, die Terroristin, terroristisch, terrorisieren}.
\item \all{Gewalt} $\rightarrow$ \all{gewalttätig, die Gewalttat, gewaltsam}.
\item \all{Sicherheit} $\rightarrow$ \all{sicher, sichern, die Sicherung, versichern}.
\item \all{Schutz} $\rightarrow$ \all{schützen, der Schutz, schützenswert}.
\end{itemize}
\item \textbf{Mémoriser les verbes avec leur construction (Rektion)} :
\begin{center}
\begin{tabular}{|l|l|l|}
\hline
\rowcolor{popblueL} Verbe & Construction & Exemple \\
\hline
\cle{beschädigen} & + Akkusativ & \all{Die Täter beschädigen das Auto.} \\
\cle{zerstören} & + Akkusativ & \all{Der Sturm hat das Haus zerstört.} \\
\cle{schützen} & + Akkusativ & \cle{Die Polizei schützt die Bürger.} \\
\cle{suchen} & \cle{nach + Dativ} & \all{Wir suchen nach den Tätern.} \\
\cle{informieren} & + Akkusativ & \all{Der Lehrer informiert die Eltern.} \\
\cle{verüben} & + Akkusativ & \all{Einen Anschlag verüben.} \\
\hline
\end{tabular}
\end{center}
\item \textbf{Collocations figées (à apprendre par cœur)} :
\all{einen Anschlag verüben}, \all{Gewalt anwenden}, \all{die Polizei rufen / verständigen}, \all{für Sicherheit sorgen}, \all{Opfer fordern}, \all{Schaden anrichten}.
\item \textbf{En production} : réemployer au moins 5 mots du tableau pour enrichir le texte.
\end{enumerate}
\end{methode}

\begin{exoresolu}[Lexique : Texte à trous]
Complétez avec les mots de la liste (forme correcte : cas, nombre) : \\
\all{die Polizei, der Anschlag, beschädigt, die Sicherheit, die Gewalt, die Täter, die Vandalen, verübt, suchen, informiert}

1. Gestern gab es in der Stadt \_\_\_\_\_\_ auf einen Bus.
2. Viele Fenster wurden \_\_\_\_\_\_.
3. \_\_\_\_\_\_ kam schnell und \_\_\_\_\_\_ zwei Personen.
4. \_\_\_\_\_\_ waren zwei junge Männer.
5. Die Bürger wollen mehr \_\_\_\_\_\_ in den Straßen.
6. \_\_\_\_\_\_ ist keine Lösung, sagen die Politiker.
7. Wer hat den \_\_\_\_\_\_ \_\_\_\_\_\_? (Partizip II)
8. Die \_\_\_\_\_\_ haben großen Schaden angerichtet.

\tcblower
\textbf{Solutions et justifications :}

1. \all{einen Anschlag} (masc., Akk. après \all{es gab}).
2. \all{beschädigt} (Passiv Präteritum : \all{wurden beschädigt} — Partizip II à la fin).
3. \all{Die Polizei} (Nom. sg., verbe \all{kam} sg.) \all{verhaftete} (Präteritum actif, 3\up{e} pers. sg.).
4. \all{Die Täter} (Nom. pl., verbe \all{waren} pl.).
5. \all{Sicherheit} (Akk. sg. sans article après \all{mehr}).
6. \all{Die Gewalt} (Nom. sg.).
7. \all{Anschlag verübt} (Partizip II de \all{verüben} : \all{verübt}, verbe à particule inséparable $\rightarrow$ pas de \all{ge-}).
8. \all{Die Vandalen} (Nom. pl., sujet de \all{haben}).

\textbf{Conclusion :} Toujours analyser le cas (Nom./Akk./Dat.), le nombre (sg./pl.) et le temps avant d'insérer le mot.
\end{exoresolu}

\begin{exercice}[Lexique : Entraînement]
1. Formez le féminin de : \all{der Täter, der Vandale, der Terrorist}.
2. Mettez au pluriel : \all{der Anschlag, das Opfer, die Reparatur, die Polizei}.
3. Complétez avec la préposition correcte : \all{suchen \_\_\_\_\_\_ den Tätern}, \all{schützen \_\_\_\_\_\_ die Bürger}, \all{informieren \_\_\_\_\_\_ die Eltern}.
4. Traduisez : « Les vandales ont causé d'énormes dégâts. » (Indice : \all{Schaden anrichten} / \all{riesig}).
\end{exercice}

\begin{corrige}[Exercice Lexique]
1. \all{die Täterin, die Vandalinin (rare, plutôt die Vandale), die Terroristin}.
2. \all{die Anschläge, die Opfer, die Reparaturen, — (pas de pluriel)}.
3. \all{nach + Dat.}, \all{+ Akk. (sans préposition)}, \all{+ Akk. (sans préposition)}.
4. \all{Die Vandalen haben riesigen Schaden angerichtet.} (ou \all{angeregt} ? Non, \all{anrichten} : \all{angerichtet}). Partizip II : \all{angerichtet}.
\end{corrige}

\courssub{Grammaire 1 : Le passif simple (Vorgangspassiv)}

\begin{propriete}[Formation du passif de procès (Vorgangspassiv)]
Le passif met l'accent sur l'\textbf{action} ou la \textbf{victime/objet}, l'agent (le "faiseur") passe au second plan ou est omis.
\begin{center}
\begin{tabular}{|l|l|l|}
\hline
\rowcolor{popblueL} Temps & Structure & Exemple (verbe \all{zerstören}) \\
\hline
\cle{Präsens} & \all{werden} (conjugué) + \textbf{Partizip II} & \all{Die Fenster \textbf{werden zerstört}.} \\
\cle{Präteritum} & \all{wurden} + \textbf{Partizip II} & \all{Die Fenster \textbf{wurden zerstört}.} \\
\cle{Perfekt} & \all{sind} + \textbf{Partizip II} + \cle{worden} & \all{Die Fenster \textbf{sind zerstört worden}.} \\
\cle{Futur I} & \all{werden} + \textbf{Partizip II} + \all{werden} & \all{Die Fenster \textbf{werden zerstört werden}.} \\
\hline
\end{tabular}
\end{center}

\textbf{Règles clés :}
\begin{itemize}
\item Le COD de l'actif devient \textbf{sujet (Nominatif)} du passif $\rightarrow$ accord du verbe \all{werden} / \all{sein} avec ce nouveau sujet.
\item L'agent (sujet actif) devient \textbf{complément d'agent} : \cle{von + Datif} (\all{von den Tätern}, \all{von der Polizei}). Souvent omis si inconnu/général.
\item \textbf{Perfekt passif} : auxiliaire \all{sein} (pas \all{haben}) + Partizip II du verbe principal + \cle{worden} (jamais \all{geworden} !).
\item Verbes intransitifs (sans COD) : pas de passif standard (sauf passif impersonnel \all{Es wird getanzt} — hors programme A1+/A2).
\end{itemize}
\end{propriete}

\begin{exemplebox}[Transformation Actif $\rightarrow$ Passif (même temps)]
\begin{tabular}{ll}
\textbf{Actif (Präsens)} & \all{Unbekannte Personen zerstören die Fenster.} \\
\textbf{Passif (Präsens)} & \all{Die Fenster werden (von unbekannten Personen) zerstört.} \\[4pt]
\textbf{Actif (Präteritum)} & \all{Die Polizei verhaftete den Täter.} \\
\textbf{Passif (Präteritum)} & \all{Der Täter wurde (von der Polizei) verhaftet.} \\[4pt]
\textbf{Actif (Perfekt)} & \all{Die Schüler haben die Wände beschmiert.} \\
\textbf{Passif (Perfekt)} & \all{Die Wände sind (von den Schülern) beschmiert worden.} \\[4pt]
\textbf{Actif (Futur I)} & \all{Die Polizei wird den Bereich absperren.} \\
\textbf{Passif (Futur I)} & \all{Der Bereich wird (von der Polizei) abgesperrt werden.}
\end{tabular}
\end{exemplebox}

\begin{attention}[Pièges majeurs pour francophones]
\begin{enumerate}
\item \textbf{Perfekt passif} : \all{sind zerstört \textbf{worden}} (pas \all{geworden}). \all{geworden} = "devenu" (verbe \all{werden} actif). \all{worden} = marqueur passif.
\item \textbf{Accord du verbe} : \all{Das Fenster \textbf{wird} zerstört} (sg.) vs \all{Die Fenster \textbf{werden} zerstört} (pl.).
\item \textbf{Complément d'agent} : \all{von \textbf{dem} Täter} (Dat. sg.), \all{von \textbf{den} Tätern} (Dat. pl.), \all{von \textbf{der} Polizei} (Dat. sg. fém.). Jamais \all{von die Täter} !
\item \textbf{Sujet passif} = Nominatif : \all{\textbf{Der} Täter wurde verhaftet} (pas \all{den Täter}).
\item \textbf{Faux ami} : \all{die Polizei} est \textbf{singulier} $\rightarrow$ \all{die Polizei \textbf{kommt}} (pas \all{kommen}).
\end{enumerate}
\end{attention}

\begin{exoresolu}[Grammaire : Transformation Actif $\rightarrow$ Passif]
Mettez les phrases suivantes au passif (conservez le temps). Omettez l'agent si "unbekannt" ou "man".

1. \all{Unbekannte Personen zerstören die Fenster.} (Präsens)
2. \all{Die Polizei verhaftete den Täter.} (Präteritum)
3. \all{Die Schüler haben die Wände beschmiert.} (Perfekt)
4. \all{Man wird die Reparatur teuer bezahlen müssen.} (Futur I, sujet impersonnel \all{man} $\rightarrow$ passif sans agent)
5. Traduisez : « Les ordinateurs ont été endommagés par les vandales. »

\tcblower
\textbf{Raisonnement pas à pas :} Identifier COD $\rightarrow$ Nouveau sujet (Nominatif) $\rightarrow$ Conjuguer \all{werden}/\all{sein} $\rightarrow$ Partizip II à la fin $\rightarrow$ (Agent : \all{von + Dat.}).

1. COD = \all{die Fenster} (pl.) $\rightarrow$ Subj. \all{Die Fenster}. Présens \all{werden} (pl.) = \all{werden}. Part. II \all{zerstört}.
\all{Die Fenster werden zerstört.}

2. COD = \all{den Täter} (masc. sg.) $\rightarrow$ Subj. \all{Der Täter}. Präteritum \all{werden} (sg.) = \all{wurde}. Part. II \all{verhaftet}.
\all{Der Täter wurde verhaftet.}

3. COD = \all{die Wände} (pl.) $\rightarrow$ Subj. \all{Die Wände}. Perfekt passif : aux \all{sein} (pl.) = \all{sind} + Part. II \all{beschmiert} + \all{worden}.
\all{Die Wände sind beschmiert worden.}

4. Sujet \all{man} $\rightarrow$ passif impersonnel ou sujet = \all{die Reparatur} (sg. fém.). Futur I : \all{werden} + Part. II + \all{werden}.
\all{Die Reparatur wird teuer bezahlt werden müssen.} (Passif modal, niveau supérieur, mais correct). Version simple : \all{Die Reparatur wird teuer bezahlt werden.}

5. « Les ordinateurs » = \all{Die Computer} (pl., sujet). « ont été endommagés » = Perfekt passif \all{sind beschädigt worden}. « par les vandales » = \all{von den Vandalen} (Dat. pl.).
\all{Die Computer sind von den Vandalen beschädigt worden.}

\textbf{Conclusion :} La formule magique du Perfekt passif : \cle{sein + Partizip II + worden}.
\end{exoresolu}

\begin{exercice}[Passif : Entraînement]
Transformez au passif (même temps) :
1. \all{Der Terrorist verübt den Anschlag.} (Präsens)
2. \all{Die Gewalt hat viele Opfer gefordert.} (Perfekt — attention : \all{fordern} + Akk.)
3. \all{Die Polizei sperrte die Straße ab.} (Präteritum — verbe à particule \all{absperren})
4. \all{Man muss die Sicherheit erhöhen.} (Präsens modal — passif modal : \all{müssen} + \all{werden} + Part. II)
\end{exercice}

\begin{corrige}[Exercice Passif]
1. \all{Der Anschlag wird verübt.}
2. \all{Viele Opfer sind gefordert worden.} (Perfekt passif : \all{sind ... gefordert worden})
3. \all{Die Straße wurde abgesperrt.} (Part. II \all{abgesperrt}, particule \all{ab-} $\rightarrow$ \all{ge-} inséré).
4. \all{Die Sicherheit muss erhöht werden.} (Passif modal : modal \all{müssen} + \all{werden} + Part. II \all{erhöht}).
\end{corrige}

\courssub{Grammaire 2 : Exprimer la cause et la conséquence}

\begin{propriete}[Connecteurs de cause et de conséquence (Satzverknüpfung)]
\begin{center}
\begin{tabularx}{\linewidth}{|X|X|X|X|}
\hline
\rowcolor{popblueL} Fonction & Structure & Place du verbe & Exemple \\
\hline
\cle{Cause} (subordonnée) & \cle{weil} / \cle{da} + sujet ... verbe \textbf{à la fin} & Verbe final & \all{Die Schule blieb closed, \textbf{weil die Fenster zerstört waren.}} \\
\cle{Cause} (préposition) & \cle{wegen} + \textbf{Génitif} (formel) / Datif (oral) & Pas de verbe (groupe nominal) & \all{\textbf{Wegen des Vandalismus} blieb die Schule geschlossen.} \\
\cle{Cause} (coordination) & \cle{denn} + sujet + verbe \textbf{2\up{e} pos.} & Verbe 2\up{e} pos. & \all{Die Schule blieb geschlossen, \textbf{denn} die Fenster \textbf{waren} zerstört.} \\
\hline
\cle{Conséquence} (adverbe) & \cle{deshalb} / \cle{deswegen} / \cle{darum} + \textbf{verbe 2\up{e} pos.} & \textbf{Inversion} (verbe juste après) & \all{Es gab einen Anschlag. \textbf{Deshalb kam} die Polizei.} \\
\cle{Conséquence} (subordonnée) & \cle{so ... dass} + sujet ... verbe \textbf{à la fin} & Verbe final & \all{Der Schaden war \textbf{so} groß, \textbf{dass} die Schule \textbf{geschlossen wurde.}} \\
\hline
\end{tabularx}
\end{center}

\textbf{Rappel ordre des mots (Topologie) :}
\begin{itemize}
\item Principal : [Connecteur/Sujet] \textbf{Verbe} ... [Compléments] ... [Participe/Infinitif].
\item Subordonnée (\all{weil, dass, so ... dass}) : \textbf{Connecteur} Sujet ... Compléments ... \textbf{Verbe final}.
\end{itemize}
\end{propriete}

\begin{exoresolu}[Grammaire : Cause et conséquence]
1. Reliez les deux phrases par \all{weil} (cause) : \all{Die Schule blieb geschlossen. Die Fenster waren zerstört.}
2. Reliez par \all{deshalb} (conséquence) : \all{Es gab einen Anschlag. Die Polizei sperrte die Straße.}
3. Reformulez la cause avec \all{wegen + Génitif} : \all{Weil es Vandalismus gab, haben die Schüler Angst.}
4. Traduisez : « Les dégâts étaient si importants que l'école a dû fermer. » (\all{so ... dass} + passif modal \all{müssen}).

\tcblower
\textbf{Raisonnement :}

1. Cause = fenêtres détruites. \all{Die Schule blieb geschlossen, weil die Fenster zerstört \textbf{waren}.} (Verbe \all{waren} à la fin). Note : \all{zerstört waren} = Zustandspassiv (état) au Präteritum.

2. Conséquence = police bloque. \all{Es gab einen Anschlag. \textbf{Deshalb sperrte} die Polizei die Straße.} (Inversion : \all{deshalb} = pos 1, \all{sperrte} = pos 2).

3. \all{Wegen des Vandalismus haben die Schüler Angst.} (\all{wegen} + Génitif \all{des Vandalismus} ; verbe \all{haben} en 2\up{e} pos. car phrase principale).

4. \all{Die Schäden waren so groß, dass die Schule geschlossen werden \textbf{musste}.} (\all{so ... dass} + subordonnée verbe final ; \all{müssen} + passif infinitif \all{geschlossen werden} $\rightarrow$ \all{musste} final).

\textbf{Conclusion :} Mémoriser le trio infernal : \cle{weil} (verbe final), \cle{wegen + Gen.} (nominal), \cle{deshalb} (inversion). Ne jamais écrire \all{*Deshalb die Polizei kam} ni \all{*Weil die Fenster waren zerstört}.
\end{exoresolu}

\begin{exercice}[Cause / Conséquence : Entraînement]
1. Choisissez le bon connecteur : \all{Die Polizei kam \_\_\_\_\_\_ (weil / deshalb) der Alarm losging.} \all{Die Polizei kam \_\_\_\_\_\_ (weil / deshalb) der Alarm losging.}
2. Transformez avec \all{wegen} : \all{Weil der Terrorismus zunimmt, gibt es mehr Kontrollen.}
3. Complétez : \all{Der Schaden war \_\_\_\_\_\_ groß, \_\_\_\_\_\_ die Reparatur \_\_\_\_\_\_ teuer.} (\all{so / dass / wird / wurde / ist})
4. Traduisez : « À cause de l'attentat, le match est annulé. » (\all{absagen} — Part. II \all{abgesagt}).
\end{exercice}

\begin{corrige}[Exercice Cause/Conséquence]
1. \all{weil} (subordonnée, verbe final \all{losging}) / \all{deshalb} (principale, inversion \all{kam}).
   Phrase 1 : \all{Die Polizei kam, weil der Alarm losging.}
   Phrase 2 : \all{Der Alarm losging. Deshalb kam die Polizei.}
2. \all{Wegen der Zunahme des Terrorismus gibt es mehr Kontrollen.} (Génitif : \all{der Zunahme}, \all{des Terrorismus}).
3. \all{so, dass, wurde} (Präteritum passif \all{wurde ... geschlossen} ou actif \all{wurde teuer} — adjectif). Correct : \all{Der Schaden war \textbf{so} groß, \textbf{dass} die Reparatur \textbf{teuer wurde}.} (devenir cher).
4. \all{Wegen des Anschlags wurde das Spiel abgesagt.} (Passif Präteritum : \all{wurde abgesagt}).
\end{corrige}

\courssub{Méthode de rédaction : Produire un paragraphe d'opinion}

\begin{methode}[Rédiger un court texte structuré (6–8 phrases) — Niveau A2]
\begin{enumerate}
\item \textbf{Introduction (1 phrase)} : Présenter le fait (qui, quoi, où, quand). \all{Letzte Woche gab es in unserer Schule Vandalismus.}
\item \textbf{Faits (2–3 phrases)} : Raconter au \cle{Präteritum} (récit) ou \cle{Perfekt} (résultat). Utiliser le \cle{Passif} si agent inconnu : \all{Fenster wurden zerstört. Wände wurden beschmiert.}
\item \textbf{Lien Cause / Conséquence (1–2 phrases)} : \all{Weil die Schäden groß waren, musste die Polizei kommen. Deshalb diskutieren wir jetzt.}
\item \textbf{Opinion personnelle (1–2 phrases)} : \cle{Ich meine, dass ...} / \cle{Meiner Meinung nach ...} / \cle{Ich finde es \textbf{wichtig / falsch / gut}, dass ...} (verbe à la fin après \all{dass}).
\item \textbf{Proposition / Conclusion (1 phrase)} : \all{Wir brauchen mehr Sicherheit und Respekt. / Die Täter müssen bestraft werden.} (Passif modal).
\item \textbf{Relecture checklist} : \cle{Majuscules noms} ? \cle{Verbe 2\up{e} pos.} ? \cle{Verbe final (weil/dass)} ? \cle{Passif correct (worden)} ? \cle{5 mots du thème} ? \cle{6–8 phrases} ?
\end{enumerate}
\end{methode}

\begin{exoresolu}[Expression écrite : « Vandalismus an Schulen — Ihre Meinung »]
Consigne : Rédigez 6–8 phrases. Mots obligatoires : \cle{die Gewalt}, \cle{beschädigen}, \cle{die Polizei}, \cle{die Sicherheit}, \cle{weil}, \cle{deshalb}.

\tcblower
\textbf{Production modèle :}

\all{In vielen Schulen gibt es heute Vandalismus. Letzte Woche wurden in unserer Schule Fenster zerstört und Wände beschmiert. Der Direktor hat die Polizei informiert, \textbf{weil} die Schäden sehr groß waren. \textbf{Deshalb} diskutieren jetzt alle Schüler über dieses Problem. \textbf{Ich meine, dass} \cle{Gewalt} niemals eine Lösung ist. Wir brauchen mehr \cle{Sicherheit} in den Schulen, aber auch mehr Respekt. \textbf{Meiner Meinung nach} müssen die Täter bestraft werden. Zusammen können wir unsere Schulen schützen.}

\textbf{Traduction :} « Dans de nombreuses écoles, il y a aujourd'hui du vandalisme. La semaine dernière, des fenêtres ont été détruites et des murs barbouillés dans notre école. Le directeur a informé la police \textbf{parce que} les dégâts étaient très importants. \textbf{C'est pourquoi} tous les élèves discutent maintenant de ce problème. \textbf{Je pense que} la violence n'est jamais une solution. Nous avons besoin de plus de \textbf{sécurité} dans les écoles, mais aussi de plus de respect. \textbf{À mon avis}, les auteurs doivent être punis. Ensemble, nous pouvons protéger nos écoles. »

\textbf{Analyse grammaticale du modèle :}
\begin{itemize}
\item Phrases 1, 4, 6, 7, 8 : Principales (verbe 2\up{e} pos. : \all{gibt, diskutieren, brauchen, müssen, können}).
\item Phrase 2 : Passif Präteritum \all{wurden zerstört / beschmiert} (agent inconnu).
\item Phrase 3 : Subordonnée \all{weil} $\rightarrow$ verbe \all{waren} final.
\item Phrase 4 : \all{Deshalb} (pos 1) $\rightarrow$ inversion \all{diskutieren} (pos 2).
\item Phrase 5 : \all{Ich meine, dass ... ist} (verbe \all{ist} final dans \all{dass}).
\item Phrase 7 : Passif modal \all{müssen ... bestraft werden}.
\item Lexique imposé : \all{Gewalt, beschädigt (Part. II), Polizei, Sicherheit, weil, deshalb} — tous employés.
\end{itemize}
\end{exoresolu}

\begin{exercice}[Expression écrite : Entraînement personnel]
Rédigez un paragraphe (6–8 phrases) sur le thème : \all{„Terrorismus und Sicherheit in unserer Stadt“}.
Utilisez obligatoirement : \cle{der Anschlag}, \cle{die Opfer}, \cle{suchen nach}, \cle{wegen}, \cle{deshalb}, \cle{meiner Meinung nach}.
Structure : Fait $\rightarrow$ Conséquences $\rightarrow$ Opinion $\rightarrow$ Solution.
\end{exercice}

\begin{corrige}[Exercice Expression écrite] (Proposition de correction)
\all{Letzten Monat gab es in unserer Stadt einen Anschlag auf einen Markt. Viele Opfer wurden verletzt. Die Polizei sucht nach den Tätern. Wegen der Gefahr bleiben viele Menschen zu Hause. Deshalb ist die Angst groß. Meiner Meinung nach muss die Regierung mehr für die Sicherheit tun. Wir dürfen uns nicht von Terror einschüchtern lassen. Solidarität ist die beste Antwort.}
\end{corrige}

\courssub{Activité orale : Signaler un acte de vandalisme / attentat}

\begin{experience}[Jeu de rôle : Appel d'urgence / Témoignage]
\textbf{Situation :} Vous êtes témoin d'un acte de vandalisme (fenêtres brisées, graffitis) ou vous entendez une alerte attentat. Appelez la police (\all{Notruf 110} en Allemagne, \all{117} au Cameroun) ou témoignez devant un officier.

\textbf{Rôles :} A = Témoin / Victime ; B = Policier / Opératrice 110.

\textbf{Phrases utiles (Redemittel) :}
\begin{itemize}
\item \all{Ich möchte einen Vandalismus / einen Anschlag melden.}
\item \all{Es passiert gerade ... / Es ist passiert ...} (Présens / Perfekt).
\item \all{Der Ort ist ... (Straße, Platz, Schule).}
\item \all{Die Täter sind ... (Anzahl, Beschreibung: jung, alt, Kleidung, Fluchtweg).}
\item \all{Es gibt Verletzte / Schäden.}
\item \all{Was soll ich tun? / Bleiben Sie am Apparat!}
\end{itemize}

\textbf{Déroulement :}
1. Préparation 5 min (notes mots-clés, pas phrases toutes faites).
2. Jeu 3 min par binôme.
3. Échange de rôles.
4. Retour classe : vocabulaire manquant, correction prononciation (\all{ü, ö, ä, ch, sch, z}), structure \all{weil/deshalb}.
\end{experience}

\courssub{Le saviez-vous ? : Sécurité scolaire en Allemagne et au Cameroun}

\begin{savaistu}[Prévention et chiffres clés]
\begin{itemize}
\item \textbf{Allemagne} : La \all{Schulpflicht} (obligation scolaire) s'accompagne d'un devoir de protection (\all{Schutzpflicht}). Depuis les attentats d'Erfurt (2002) et Winnenden (2009), les écoles ont des \all{Sicherheitskonzepte} (plans de sécurité) : portiques, agents de sécurité (\all{Sicherheitsdienste}), exercices d'alerte (\all{Amokalarm-Übungen}). La police fait de la prévention (\all{Prävention}) dans les classes (projet \all{„Gewaltfrei Lernen“}).
\item \textbf{Cameroun} : Le \cle{MINESEC} et le \cle{MINAT} (Ministère de l'Administration Territoriale) coordonnent la sécurité des établissements. Après les crises dans les régions du Nord-Ouest/Sud-Ouest (2017–), les écoles ont mis en place des \all{Comités de vigilance} et des barrières. La police et la gendarmerie assurent des patrouilles. Le vandalisme (cassse de bancs, tags) est sanctionné par le règlement intérieur et le Code pénal (Art. 322 ss. Code pénal : \all{Sachbeschädigung}).
\item \textbf{Vocabulaire comparé} :
\begin{center}
\begin{tabular}{ll}
\all{der Amokalarm} & alerte attentat / tireur actif (école) \\
\all{der Sicherheitsdienst} & service de sécurité / agents \\
\all{die Prävention} & la prévention \\
\all{der Notruf 110} & numéro police (Allemagne) \\
\all{der Vandalismus} & le vandalisme (terme commun)
\end{tabular}
\end{center}
\end{itemize}
\end{savaistu}

\courssub{Erreurs fréquentes à éviter (Check-list Bac)}

\begin{attention}[Les 5 erreurs qui coûtent des points à l'examen]
\begin{enumerate}
\item \textbf{Majuscules oubliées aux noms} : \all{die polizei}, \all{die gewalt}, \all{der anschlag} $\rightarrow$ \textbf{Faute grave}. Écrire \all{die \textbf{P}olizei}, \all{die \textbf{G}ewalt}, \all{der \textbf{A}nschlag}. \emph{Règle : Tout nom allemand prend une majuscule, sans exception.}
\item \textbf{Perfekt du passif mal formé} : \all{*sind beschädigt geworden} $\rightarrow$ \textbf{Faux}. Correct : \all{sind beschädigt \textbf{worden}}. Rappel : \all{geworden} = "devenu" (verbe \all{werden} actif) ; \all{worden} = marqueur passif.
\item \textbf{Place du verbe après connecteurs} :
\begin{itemize}
\item \all{weil / dass / so ... dass} $\rightarrow$ \textbf{Verbe À LA FIN}. \all{..., weil die Fenster zerstört \textbf{waren}.}
\item \all{deshalb / deswegen / darum / dann} $\rightarrow$ \textbf{Verbe EN 2\up{e} POSITION (inversion)}. \all{\textbf{Deshalb kam} die Polizei.}
\end{itemize}
Ne jamais écrire : \all{*Deshalb die Polizei kam} ou \all{*Weil die Fenster waren zerstört}.
\item \textbf{Faux amis et calques du français} :
\begin{itemize}
\item « La police » = \all{die Polizei} (singulier) $\rightarrow$ \all{die Polizei \textbf{kommt}} (pas \all{kommen}).
\item « Avoir peur » = \all{\textbf{Angst haben}} (pas \all{Angst sein}, pas \all{*haben Furcht} au sens courant).
\item « La violence » = \all{die Gewalt} (pas \all{die Macht} = le pouvoir/la force).
\item « Un attentat » = \all{der Anschlag} (pas \all{das Attentat} — existe mais plus formel/historique).
\end{itemize}
\item \textbf{Cas du complément d'agent (passif)} : « Par » = \cle{von + DATIF}.
\begin{itemize}
\item \all{von \textbf{dem} Täter} (masc. sg.), \all{von \textbf{der} Polizei} (fém. sg.), \all{von \textbf{den} Tätern} (pl.).
\item Sujet passif = \textbf{Nominatif} : \all{\textbf{Der} Täter wurde verhaftet} (pas \all{*den Täter}).
\end{itemize}
\end{enumerate}
\end{attention}

\begin{aretenir}[Synthèse de la leçon ALA25]
\begin{itemize}
\item \textbf{Texte} : Comprendre un fait divers (vandalisme/terrorisme) — repérer 5 W, faits/opinions, reformuler.
\item \textbf{Lexique} : Noms (article + pluriel), familles (\all{Terror/Terrorismus}, \all{Gewalt/gewalttätig}), verbes + construction (\all{nach+Dat}, \all{Akk}), collocations (\all{Anschlag verüben}).
\item \textbf{Grammaire 1} : \cle{Vorgangspassiv} (Präsens \all{werden + Part.II}, Präteritum \all{wurden + Part.II}, Perfekt \all{sind + Part.II + worden}, Futur \all{werden + Part.II + werden}). Agent : \all{von + Dat.}
\item \textbf{Grammaire 2} : Cause (\all{weil/da} + verbe final, \all{wegen+Gen.}, \all{denn} + V2) ; Conséquence (\all{deshalb} + inversion, \all{so...dass} + verbe final).
\item \textbf{Expression} : Rédiger 6–8 phrases structurées (Intro, Faits/Passif, Cause/Conséquence, Opinion/\all{dass}, Solution/Passif modal). Check-list relecture.
\item \textbf{Oral} : Signaler un incident (numéros d'urgence, description lieu/auteurs/dégâts).
\end{itemize}
\end{aretenir}

\newpage

\coursec{Exercices}

\courssub{Je vérifie mes connaissances}

\begin{exercice}[QCM : vocabulaire du thème]
Entoure la bonne réponse.
\begin{enumerate}
\item \all{der Anschlag} signifie :
\begin{itemize}
\item[a)] la manifestation
\item[b)] l'attentat
\item[c)] la réparation
\end{itemize}
\item Le pluriel de \all{die Gewalt} est :
\begin{itemize}
\item[a)] \all{die Gewalte}
\item[b)] \all{die Gewalten}
\item[c)] ce nom n'a pas de pluriel
\end{itemize}
\item \all{beschädigen} signifie :
\begin{itemize}
\item[a)] réparer
\item[b)] endommager
\item[c)] construire
\end{itemize}
\item \all{die Sicherheit} signifie :
\begin{itemize}
\item[a)] la sécurité
\item[b)] la certitude seulement
\item[c)] la police
\end{itemize}
\end{enumerate}
\end{exercice}

\begin{exercice}[Vrai ou faux ? Justifie ta réponse]
Réponds par vrai ou faux et justifie en une phrase.
\begin{enumerate}
\item Au passif, le verbe \all{werden} se conjugue et \all{sein} reste à l'infinitif.
\item Dans la phrase \all{Die Bushaltestelle wird beschädigt}, le sujet est celui qui fait l'action.
\item \all{wegen} est suivi du génitif (ou du datif dans la langue courante).
\item \all{deshalb} est une conjonction de coordination qui n'influence pas la place du verbe.
\end{enumerate}
\end{exercice}

\begin{exercice}[Familles de mots]
Complète le tableau avec le mot de la même famille.

\begin{center}
\begin{tabularx}{\linewidth}{|X|X|X|}
\hline
\rowcolor{popblueL} Verbe & Nom & Traduction du nom \\
\hline
\all{beschädigen} & die \ldots\ldots\ldots\ldots & \\
\hline
\all{verhaften} & die \ldots\ldots\ldots\ldots & \\
\hline
\all{schützen} & der \ldots\ldots\ldots\ldots & \\
\hline
\all{sichern} & die \ldots\ldots\ldots\ldots & \\
\hline
\end{tabularx}
\end{center}
\end{exercice}

\begin{exercice}[Conjugaison : le passif à compléter]
Complète avec la forme correcte de \all{werden} au temps indiqué.
\begin{enumerate}
\item (Präsens) Der Täter \ldots\ldots\ldots\ldots von der Polizei verhaftet.
\item (Präsens) Die Fenster der Schule \ldots\ldots\ldots\ldots zerstört.
\item (Präteritum) Das Denkmal \ldots\ldots\ldots\ldots gestern beschmiert.
\item (Präteritum) Viele Häuser \ldots\ldots\ldots\ldots durch den Sturm beschädigt.
\item (Perfekt) Der Bahnhof \ldots\ldots\ldots\ldots \ldots\ldots\ldots\ldots gesichert.
\end{enumerate}
\end{exercice}

\courssub{Je m'entraîne}

\begin{exercice}[Transformation : de l'actif au passif]
Mets les phrases suivantes au passif, d'abord au présent, puis au prétérit.
\begin{enumerate}
\item \all{Jugendliche beschädigen die Bushaltestelle.}
\item \all{Die Polizei verhaftet den Täter.}
\item \all{Die Schüler streichen Graffiti an die Wand.}
\item \all{Die Regierung verstärkt die Sicherheit.}
\end{enumerate}
\end{exercice}

\begin{exercice}[La cause et la conséquence]
Complète avec \all{weil}, \all{denn}, \all{wegen} ou \all{deshalb}. Attention à la place du verbe !
\begin{enumerate}
\item \ldots\ldots\ldots\ldots der Anschläge wird die Sicherheit verstärkt.
\item Die Bürger haben Angst, \ldots\ldots\ldots\ldots fordern sie mehr Polizisten.
\item Der Täter wurde verhaftet, \ldots\ldots\ldots\ldots er die Bushaltestelle beschädigt hatte.
\item Die Schule bleibt heute geschlossen, \ldots\ldots\ldots\ldots die Fenster zerstört wurden.
\item \ldots\ldots\ldots\ldots des Vandalismus installiert die Stadt Kameras.
\end{enumerate}
\end{exercice}

\begin{exercice}[Texte à trous]
Complète le texte avec les mots suivants : \all{die Gewalt}, \all{beschädigt}, \all{die Polizei}, \all{der Anschlag}, \all{die Sicherheit}, \all{verhaftet}.

\begin{textebox}[Ein Vorfall in der Stadt]
Gestern Abend wurde eine Bushaltestelle in der Innenstadt \ldots\ldots\ldots\ldots. Die Fensterscheiben waren kaputt. Ein Zeuge rief sofort \ldots\ldots\ldots\ldots. Zwei Beamte kamen schnell und der junge Täter wurde noch am Ort \ldots\ldots\ldots\ldots. Der Bürgermeister sagte: „\ldots\ldots\ldots\ldots gegen öffentliche Gebäude ist ein großes Problem. Wir müssen \ldots\ldots\ldots\ldots in unserer Stadt verbessern.“ Nach dem \ldots\ldots\ldots\ldots auf den Bahnhof im letzten Monat sind die Bürger sehr besorgt.
\end{textebox}
\end{exercice}

\begin{exercice}[Appariements]
Relie chaque début de phrase à la bonne fin.

\begin{center}
\begin{tabularx}{\linewidth}{|X|X|}
\hline
\rowcolor{popblueL} Début & Fin \\
\hline
1. Wegen des Terroranschlags & a) wurde der Täter von der Polizei verhaftet. \\
\hline
2. Weil er das Denkmal beschmiert hatte, & b) wurde der Flughafen geschlossen. \\
\hline
3. Die Sicherheit wird verstärkt, & c) deshalb bleiben viele Touristen zu Hause. \\
\hline
4. Die Menschen haben Angst vor Anschlägen, & d) denn die Stadt will die Bürger schützen. \\
\hline
\end{tabularx}
\end{center}
\end{exercice}

\courssub{Je me prépare au Bac}

\begin{exercice}[Compréhension de l'écrit et questions de langue]
Lis le texte suivant, puis réponds aux questions en allemand, par des phrases complètes.

\begin{textebox}[Vandalismus an deutschen Schulen]
In vielen deutschen Städten wird über ein wachsendes Problem gesprochen: Vandalismus an Schulen. Fenster werden zerstört, Wände werden mit Graffiti beschmiert und Schultische werden beschädigt. Allein in einer Großstadt wurden im letzten Jahr über zweihundert Fälle gemeldet. Die Reparaturen kosten die Kommunen viel Geld — Geld, das für Bücher und Computer fehlt.

Warum passiert das? Viele Täter sind Jugendliche zwischen vierzehn und siebzehn Jahren. Manche wollen Langeweile bekämpfen, andere protestieren gegen die Schule oder suchen Anerkennung in ihrer Gruppe. Wegen der vielen Vorfälle wird die Sicherheit an den Schulen verstärkt: Kameras werden installiert und Wachpersonal wird eingestellt.

Aber reicht das? Einige Lehrer meinen, dass Prävention wichtiger ist als Strafe. „Die Schüler müssen lernen, dass die Schule ihr eigenes Haus ist“, sagt eine Schulleiterin aus Hamburg. Deshalb werden in vielen Schulen Projekte organisiert, bei denen die Jugendlichen selbst ihre Klassenräume renovieren. Wer die Schule mitgestaltet, zerstört sie nicht — so die Hoffnung der Pädagogen.
\end{textebox}

\textbf{Comprehension questions (beantworte auf Deutsch) :}
\begin{enumerate}
\item Was wird an deutschen Schulen zerstört oder beschädigt? Nenne drei Beispiele.
\item Wie alt sind die meisten Täter?
\item Nenne zwei Gründe für den Vandalismus der Jugendlichen.
\item Welche Maßnahmen werden wegen der Vorfälle ergriffen?
\item Was meint die Schulleiterin aus Hamburg zur Prävention?
\item \textbf{Deine Meinung :} Glaubst du, dass Kameras an Schulen eine gute Lösung sind? Begründe deine Antwort in zwei oder drei Sätzen.
\end{enumerate}

\textbf{Fragen zur Sprache :}
\begin{enumerate}
\item Suche im Text drei Passivsätze und schreibe sie ab.
\item Setze den folgenden Satz ins Präteritum-Passiv : \all{Fenster werden zerstört.}
\item Ergänze : \all{Wegen der vielen Vorfälle \ldots\ldots\ldots\ldots die Sicherheit verstärkt.} (Zeitform und Verbform erklären)
\end{enumerate}
\end{exercice}

\begin{exercice}[Thème et version]
\textbf{Version.} Traduis le paragraphe suivant en français.

\begin{textebox}[Ein Anschlag wird verhindert]
Gestern wurde in Berlin ein Anschlag von der Polizei verhindert. Ein verdächtiger Koffer war am Bahnhof gefunden worden. Sofort wurde der Bahnhof gesperrt und die Reisenden wurden evakuiert. Spezialisten wurden gerufen, und der Koffer wurde geöffnet: Es war nur eine Tasche mit Kleidern. Trotzdem lobte der Innenminister die Arbeit der Beamten. „Wegen der schnellen Reaktion der Polizei waren die Bürger nie in Gefahr“, sagte er. Die Sicherheit an allen Bahnhöfen wird jetzt noch einmal überprüft.
\end{textebox}

\textbf{Thème.} Traduis les phrases suivantes en allemand.
\begin{enumerate}
\item Les murs ont été barbouillés parce que la sécurité n'était pas renforcée.
\item L'auteur (le coupable) a été arrêté par la police hier soir.
\item À cause des attentats, beaucoup de gens ont peur.
\item Les vitres de l'école sont détruites chaque année ; c'est pourquoi la ville installe des caméras.
\end{enumerate}
\end{exercice}

\begin{exercice}[Expression écrite guidée]
Choisis \textbf{un} des deux sujets suivants et rédige un texte en allemand.

\textbf{Sujet A — Résumé :} Résume le texte \all{Vandalismus an deutschen Schulen} (exercice de compréhension) en \textbf{quatre phrases} (environ 40 à 50 mots). Tu dois employer \textbf{au moins trois formes passives} et \textbf{un connecteur de cause ou de conséquence} (\all{weil}, \all{denn}, \all{wegen}, \all{deshalb}).

\textbf{Sujet B — Lettre au maire :} Dans ton quartier à Yaoundé, des abribus et des bancs publics ont été endommagés. Écris une lettre courte (\textbf{60 à 80 mots}) au maire : décris les dégâts, explique les conséquences pour les habitants et propose deux solutions (par exemple plus de police, des caméras, des projets avec les jeunes). Utilise au moins deux passifs et deux expressions de cause ou de conséquence. N'oublie pas la formule d'appel et la formule de politesse.
\end{exercice}

\newpage

\coursec{Corrigés des exercices}

\coursec{Corrigés des exercices — Leçon ALA25 : Le vandalisme et le terrorisme}

\begin{corrige}[Exercice 1 — QCM : vocabulaire du thème]
\begin{enumerate}
\item \textbf{b)} l'attentat. \all{der Anschlag, die Anschläge} désigne un attentat, une attaque. On dit aussi \all{der Terroranschlag} (l'attentat terroriste).
\item \textbf{c)} ce nom n'a pas de pluriel. \all{die Gewalt} (la violence) est un nom abstrait employé seulement au singulier, comme en français.
\item \textbf{b)} endommager. \all{beschädigen} = endommager, détériorer ; le contraire est \all{reparieren} (réparer). \cle{Attention} : \all{beschädigen} est un verbe faible à particule inséparable \all{be-} : \all{er beschädigt, er beschädigte, er hat beschädigt} (pas de \all{ge-} au participe passé).
\item \textbf{a)} la sécurité. \all{die Sicherheit} signifie d'abord la sécurité ; le sens de « certitude » existe mais n'est pas le sens du thème. La police se dit \all{die Polizei} (toujours au singulier).
\end{enumerate}
\end{corrige}

\begin{corrige}[Exercice 2 — Vrai ou faux ? Justifie ta réponse]
\begin{enumerate}
\item \textbf{Vrai.} Au passif de procès (\emph{Prozesspassiv}), l'auxiliaire \all{werden} se conjugue et porte la marque du temps et de la personne ; le verbe lexical est au participe passé : \all{Der Täter wird verhaftet.} « Le coupable est arrêté. » (\cle{Attention} : ce n'est pas \all{sein} qui intervient, sauf au passif d'état \emph{Zustandspassiv} : \all{Das Fenster ist zerstört.} « La fenêtre est détruite. »)
\item \textbf{Faux.} Au passif, le sujet grammatical \emph{subit} l'action : c'est l'abribus qui est endommagé. Celui qui fait l'action (l'agent) est introduit par \all{von} + datif : \all{Die Bushaltestelle wird von Jugendlichen beschädigt.} « L'abribus est endommagé par des jeunes. »
\item \textbf{Vrai.} \all{wegen} + génitif : \all{wegen des Vandalismus} « à cause du vandalisme ». Dans la langue courante et parlée, le datif est fréquent : \all{wegen dem Vandalismus}. À l'écrit soigné, on préfère le génitif.
\item \textbf{Faux.} \all{deshalb} n'est pas une conjonction de coordination mais un \emph{adverbe} de conséquence : placé en tête, il occupe la première position et provoque l'inversion du sujet : \all{Die Fenster wurden zerstört, deshalb bleibt die Schule geschlossen.} « Les fenêtres ont été détruites, c'est pourquoi l'école reste fermée. » (verbe en 2\textsuperscript{e} position, sujet après le verbe).
\end{enumerate}
\end{corrige}

\begin{corrige}[Exercice 3 — Familles de mots]
\begin{center}
\begin{tabularx}{\linewidth}{|X|X|X|}
\hline
\rowcolor{popblueL} Verbe & Nom & Traduction du nom \\
\hline
\all{beschädigen} & die \textbf{Beschädigung}, -en & l'endommagement, la détérioration \\
\hline
\all{verhaften} & die \textbf{Verhaftung}, -en & l'arrestation \\
\hline
\all{schützen} & der \textbf{Schutz} (sans pluriel) & la protection \\
\hline
\all{sichern} & die \textbf{Sicherung}, -en & la sécurisation, la protection \\
\hline
\end{tabularx}
\end{center}
\textbf{Remarques :} les noms en \all{-ung} sont toujours féminins (\all{die}) et forment leur pluriel en \all{-en}. \all{der Schutz} est un nom abstrait sans pluriel ; ne pas confondre avec \all{die Sicherheit} (la sécurité / la sûreté).
\end{corrige}

\begin{corrige}[Exercice 4 — Conjugaison : le passif à compléter]
\begin{enumerate}
\item Der Täter \textbf{wird} von der Polizei verhaftet. « Le coupable est arrêté par la police. » (\all{der Täter} = 3\textsuperscript{e} pers. sg., Präsens)
\item Die Fenster der Schule \textbf{werden} zerstört. « Les fenêtres de l'école sont détruites. » (Pluriel, Präsens)
\item Das Denkmal \textbf{wurde} gestern beschmiert. « Le monument a été barbouillé hier. » (Präteritum, sg. : \all{wurde})
\item Viele Häuser \textbf{wurden} durch den Sturm beschädigt. « Beaucoup de maisons ont été endommagées par la tempête. » (Präteritum, pl. : \all{wurden} ; \all{durch} + accusatif pour la cause/l'agent naturel)
\item \textbf{Erreur fréquente :} \all{Der Bahnhof ist worden gesichert.} \\
\textbf{Forme correcte :} \all{Der Bahnhof \textbf{ist} gesichert \textbf{worden}.} « La gare a été sécurisée. » (Perfekt passif : \all{sein} conjugué + participe passé + \all{worden}, \cle{sans \all{ge-}} !)
\end{enumerate}
\textbf{Point de vigilance :} au Perfekt passif, le participe de \all{werden} est \all{worden} (et non \all{geworden}) : \all{Es ist repariert worden.} « Cela a été réparé. »
\end{corrige}

\begin{corrige}[Exercice 5 — Transformation : de l'actif au passif]
\textbf{Méthode :} le COD de la phrase active devient sujet (nominatif) ; le verbe passe à \all{werden} + participe passé ; le sujet actif devient complément d'agent introduit par \all{von} + datif.
\begin{enumerate}
\item Présent : \all{Die Bushaltestelle wird (von Jugendlichen) beschädigt.} « L'abribus est endommagé (par des jeunes). »\\
Prétérit : \all{Die Bushaltestelle wurde (von Jugendlichen) beschädigt.}
\item Présent : \all{Der Täter wird von der Polizei verhaftet.} « Le coupable est arrêté par la police. »\\
Prétérit : \all{Der Täter wurde von der Polizei verhaftet.}
\item Présent : \all{Graffiti werden (von den Schülern) an die Wand gestrichen.} « Des graffitis sont peints sur le mur. »\\
Prétérit : \all{Graffiti wurden (von den Schülern) an die Wand gestrichen.}\\
(\all{streichen} est un verbe fort : participe \all{gestrichen}.)
\item Présent : \all{Die Sicherheit wird (von der Regierung) verstärkt.} « La sécurité est renforcée (par le gouvernement). »\\
Prétérit : \all{Die Sicherheit wurde (von der Regierung) verstärkt.}
\end{enumerate}
\textbf{Point de vigilance :} ne pas traduire « par » par \all{bei} ou \all{mit} ; l'agent au passif s'introduit avec \all{von} (+ datif).
\end{corrige}

\begin{corrige}[Exercice 6 — La cause et la conséquence]
\begin{enumerate}
\item \textbf{Wegen} der Anschläge wird die Sicherheit verstärkt. « À cause des attentats, la sécurité est renforcée. » (\all{wegen} + génitif pluriel \all{der Anschläge} ; inversion du verbe.)
\item Die Bürger haben Angst, \textbf{deshalb} fordern sie mehr Polizisten. « Les citoyens ont peur, c'est pourquoi ils demandent plus de policiers. » (\all{deshalb} adverbe en tête → inversion : \all{fordern sie}.)
\item Der Täter wurde verhaftet, \textbf{weil} er die Bushaltestelle beschädigt \textbf{hatte}. « Le coupable a été arrêté parce qu'il avait endommagé l'abribus. » (\all{weil} : subordonnée, verbe conjugué \all{hatte} à la fin ; Plusquamperfekt pour l'antériorité.)
\item Die Schule bleibt heute geschlossen, \textbf{denn} die Fenster wurden zerstört. « L'école reste fermée aujourd'hui, car les fenêtres ont été détruites. » (\all{denn} : coordination, pas d'inversion, verbe en 2\textsuperscript{e} position.)
\item \textbf{Wegen} des Vandalismus installiert die Stadt Kameras. « À cause du vandalisme, la ville installe des caméras. » (\all{wegen} + génitif sg. \all{des Vandalismus} ; le groupe prépositionnel en tête entraîne l'inversion : \all{installiert die Stadt}.)
\end{enumerate}
\textbf{Rappel synthétique :}
\begin{center}
\begin{tabular}{|l|l|l|}
\hline
\rowcolor{popblueL} Connecteur & Structure & Place du verbe \\
\hline
\all{weil} & Subordonnée & À la fin \\
\hline
\all{denn} & Coordination (phrase principale) & 2\textsuperscript{e} position \\
\hline
\all{deshalb} / \all{darum} & Adverbe (phrase principale) & 2\textsuperscript{e} position (inversion si en tête) \\
\hline
\all{wegen} + Nom (Génitif) & Groupe prépositionnel & 2\textsuperscript{e} position (inversion si en tête) \\
\hline
\end{tabular}
\end{center}
\end{corrige}

\begin{corrige}[Exercice 7 — Texte à trous]
\begin{textebox}[Texte à compléter]
Gestern Abend wurde eine Bushaltestelle in der Innenstadt \textbf{beschädigt}. Die Fensterscheiben waren kaputt. Ein Zeuge rief sofort \textbf{die Polizei}. Zwei Beamte kamen schnell und der junge Täter wurde noch am Ort \textbf{verhaftet}. Der Bürgermeister sagte: „\textbf{Die Gewalt} gegen öffentliche Gebäude ist ein großes Problem. Wir müssen \textbf{die Sicherheit} in unserer Stadt verbessern.“ Nach dem \textbf{Anschlag} auf den Bahnhof im letzten Monat sind die Bürger sehr besorgt.
\end{textebox}

\textbf{Traduction :} « Hier soir, un abribus du centre-ville a été endommagé. Les vitres étaient cassées. Un témoin a immédiatement appelé la police. Deux agents sont vite arrivés et le jeune coupable a été arrêté sur place. Le maire a dit : "La violence contre les bâtiments publics est un grand problème. Nous devons améliorer la sécurité dans notre ville." Après l'attentat contre la gare le mois dernier, les citoyens sont très inquiets. »

\textbf{Justifications :}
\begin{itemize}
\item \all{beschädigt}, \all{verhaftet} : participes passés du passif Präteritum (\all{wurde ... beschädigt/verhaftet}).
\item \all{die Polizei} : nom féminin, toujours au singulier.
\item \all{Die Gewalt}, \all{die Sicherheit} : noms du thème, majuscule obligatoire.
\item \all{nach dem Anschlag} : \all{nach} + datif singulier (\all{dem}).
\end{itemize}
\end{corrige}

\begin{corrige}[Exercice 8 — Appariements cause / conséquence]
\begin{center}
\begin{tabularx}{\linewidth}{|X|X|}
\hline
\rowcolor{popblueL} Début & Fin \\
\hline
1. Wegen des Terroranschlags & \textbf{b)} wurde der Flughafen geschlossen. \\
\hline
2. Weil er das Denkmal beschmiert hatte, & \textbf{a)} wurde der Täter von der Polizei verhaftet. \\
\hline
3. Die Sicherheit wird verstärkt, & \textbf{d)} denn die Stadt will die Bürger schützen. \\
\hline
4. Die Menschen haben Angst vor Anschlägen, & \textbf{c)} deshalb bleiben viele Touristen zu Hause. \\
\hline
\end{tabularx}
\end{center}
\textbf{Traductions :}
\begin{enumerate}
\item « À cause de l'attentat terroriste, l'aéroport a été fermé. »
\item « Parce qu'il avait barbouillé le monument, le coupable a été arrêté par la police. »
\item « La sécurité est renforcée, car la ville veut protéger les citoyens. »
\item « Les gens ont peur des attentats, c'est pourquoi beaucoup de touristes restent chez eux. »
\end{enumerate}
\textbf{Justifications grammaticales :}
\begin{itemize}
\item 1 : \all{wegen} + génitif (\all{des Terroranschlags}) en tête → inversion (\all{wurde der Flughafen}).
\item 2 : Subordonnée \all{weil}... \all{hatte} (verbe à la fin) en position 1 → inversion dans la principale (\all{wurde der Täter}).
\item 3 : \all{denn} coordonne deux principales → pas d'inversion (\all{die Stadt will}).
\item 4 : \all{deshalb} adverbe en tête → inversion (\all{bleiben viele Touristen}).
\end{itemize}
\end{corrige}

\begin{corrige}[Exercice 9 — Compréhension de l'écrit et questions de langue]
\textbf{Barème indicatif (20 points) :} compréhension 6 × 1,5 = 9 pts ; questions de langue 3 × 2 = 6 pts ; qualité de la langue (ordre des mots, orthographe, majuscules) 5 pts.

\textbf{Comprehension questions — réponses modèles :}
\begin{enumerate}
\item \all{An deutschen Schulen werden Fenster zerstört, Wände werden mit Graffiti beschmiert und Schultische werden beschädigt.} « Des fenêtres sont détruites, des murs sont barbouillés de graffitis et des tables d'école sont endommagées. »
\item \all{Die meisten Täter sind Jugendliche zwischen vierzehn und siebzehn Jahren.} « La plupart des auteurs sont des jeunes entre quatorze et dix-sept ans. »
\item \all{Manche wollen Langeweile bekämpfen, andere protestieren gegen die Schule oder suchen Anerkennung in ihrer Gruppe.} « Certains veulent combattre l'ennui, d'autres protestent contre l'école ou cherchent la reconnaissance dans leur groupe. » (Deux raisons suffisent.)
\item \all{Wegen der vielen Vorfälle wird die Sicherheit verstärkt: Kameras werden installiert und Wachpersonal wird eingestellt.} « À cause des nombreux incidents, la sécurité est renforcée : des caméras sont installées et du personnel de surveillance est engagé. »
\item \all{Sie meint, dass Prävention wichtiger ist als Strafe, und dass die Schüler lernen müssen, dass die Schule ihr eigenes Haus ist.} « Elle pense que la prévention est plus importante que la punition et que les élèves doivent apprendre que l'école est leur propre maison. »
\item \textbf{Exemple de réponse personnelle :} \all{Ich glaube, dass Kameras keine gute Lösung sind, weil sie die Schüler kontrollieren, aber nicht erziehen. Ich finde Projekte besser, denn wer die Schule mitgestaltet, zerstört sie nicht.} « Je pense que les caméras ne sont pas une bonne solution, parce qu'elles contrôlent les élèves mais ne les éduquent pas. Je trouve les projets meilleurs, car celui qui participe à l'aménagement de l'école ne la détruit pas. »
\end{enumerate}

\textbf{Fragen zur Sprache :}
\begin{enumerate}
\item Trois passifs possibles (au choix) : \all{Fenster werden zerstört.} / \all{Wände werden mit Graffiti beschmiert.} / \all{Schultische werden beschädigt.} / \all{Kameras werden installiert.} / \all{Wachpersonal wird eingestellt.} / \all{Allein in einer Großstadt wurden im letzten Jahr über zweihundert Fälle gemeldet.}
\item \all{Fenster wurden zerstört.} « Des fenêtres furent détruites / ont été détruites. » (Präteritum de \all{werden}, pluriel : \all{wurden}.)
\item \all{Wegen der vielen Vorfälle \textbf{wird} die Sicherheit verstärkt.} Il s'agit du \textbf{Präsens passif} : \all{wird} (3\textsuperscript{e} pers. sg. de \all{werden}, accord avec \all{die Sicherheit}) + participe \all{verstärkt}. Le groupe causal \all{wegen der vielen Vorfälle} occupe la première position, donc le verbe reste en deuxième position et le sujet suit (inversion).
\end{enumerate}
\textbf{Points de vigilance :} majuscule à tous les noms ; verbe en 2\textsuperscript{e} position dans les réponses ; ne pas oublier l'inversion après un complément placé en tête.
\end{corrige}

\begin{corrige}[Exercice 10 — Thème et version]
\textbf{Barème indicatif (20 points) :} version 10 pts (fidélité 5, correction du français 3, passifs bien rendus 2) ; thème 10 pts (2,5 pts par phrase : morphologie du passif 1 pt, connecteurs et place du verbe 1 pt, vocabulaire et orthographe 0,5 pt).

\textbf{Version — Allemand → Français}
\textbf{Texte source (allemand) :}
\all{Gestern wurde ein Anschlag in Berlin von der Polizei verhindert. Ein verdächtiger Koffer war am Bahnhof gefunden worden. Sofort wurde der Bahnhof gesperrt und die Reisenden wurden evakuiert. Spezialisten wurden gerufen, und der Koffer wurde geöffnet: es war nur eine Tasche mit Kleidung. Trotzdem lobte der Innenminister die Arbeit der Beamten. „Dank der schnellen Reaktion der Polizei waren die Bürger nie in Gefahr“, sagte er. Die Sicherheit in allen Bahnhöfen wird jetzt noch einmal überprüft.}

\textbf{Traduction proposée :}
« Hier, un attentat a été empêché à Berlin par la police. Une valise suspecte avait été trouvée à la gare. Aussitôt, la gare a été fermée et les voyageurs ont été évacués. Des spécialistes ont été appelés, et la valise a été ouverte : ce n'était qu'un sac avec des vêtements. Néanmoins, le ministre de l'Intérieur a salué le travail des agents. "Grâce à la réaction rapide de la police, les citoyens n'ont jamais été en danger", a-t-il dit. La sécurité dans toutes les gares va maintenant être vérifiée une nouvelle fois. »

\textbf{Points de vigilance (version) :}
\begin{itemize}
\item \all{war ... gefunden worden} : Plusquamperfekt passif → « avait été trouvée ».
\item \all{verhindern} = empêcher (faux ami : ce n'est pas « venir »).
\item \all{sperren} = fermer, bloquer (une route, un lieu).
\item \all{trotzdem} = néanmoins, pourtant (adverbe, pas conjonction).
\item \all{der Innenminister} = le ministre de l'Intérieur.
\end{itemize}

\textbf{Thème — Français → Allemand}
\begin{enumerate}
\item \textbf{Source :} « Les murs ont été barbouillés parce que la sécurité n'avait pas été renforcée. »\\
\textbf{Cible :} \all{Die Wände wurden beschmiert, weil die Sicherheit nicht verstärkt worden war.} (ou : \all{..., denn die Sicherheit war nicht verstärkt worden.}) — Plusquamperfekt passif : \all{war ... worden} ; verbe à la fin après \all{weil}.
\item \textbf{Source :} « Le coupable a été arrêté hier soir par la police. »\\
\textbf{Cible :} \all{Der Täter wurde gestern Abend von der Polizei verhaftet.} — Präteritum passif : \all{wurde verhaftet} ; agent avec \all{von} + datif.
\item \textbf{Source :} « À cause des attentats, beaucoup de gens ont peur. »\\
\textbf{Cible :} \all{Wegen der Anschläge haben viele Menschen Angst.} — \all{wegen} + génitif pluriel ; inversion après le groupe en tête (\all{haben viele}).
\item \textbf{Source :} « Les vitres de l'école sont détruites chaque année ; c'est pourquoi la ville installe des caméras. »\\
\textbf{Cible :} \all{Die Fensterscheiben der Schule werden jedes Jahr zerstört; deshalb installiert die Stadt Kameras.} — Präsens passif au pluriel : \all{werden zerstört} ; \all{deshalb} + inversion (\all{installiert die Stadt}).
\end{enumerate}
\textbf{Points de vigilance (thème) :}
\begin{itemize}
\item Ne pas écrire \all{geworden} au lieu de \all{worden} (Perfekt/Plusquamperfekt passif).
\item « Les vitres » se dit \all{die Fensterscheiben} (pluriel de \all{die Fensterscheibe}).
\item « C'est pourquoi » = \all{deshalb} (avec inversion) et non \all{*das ist warum}.
\item \all{installieren} : participe passé \all{installiert} (verbe en \all{-ieren} → pas de \all{ge-}).
\end{itemize}
\end{corrige}

\begin{corrige}[Exercice 11 — Expression écrite guidée]
\textbf{Barème indicatif (20 points) :} respect de la consigne et du nombre de mots 4 pts ; contenu et cohérence 4 pts ; passifs correctement formés (au moins 2 ou 3 selon le sujet) 4 pts ; connecteurs de cause/conséquence (au moins 1 ou 2) 4 pts ; orthographe, majuscules, ordre des mots 4 pts.

\textbf{Sujet A — Proposition de résumé modèle (environ 45 mots) :}
\begin{textebox}[Résumé]
\all{An deutschen Schulen werden Fenster zerstört und Wände beschmiert. Die Täter sind meist Jugendliche, weil sie Langeweile haben oder Anerkennung suchen. Wegen der vielen Vorfälle werden Kameras installiert. Aber Prävention ist wichtiger: Die Schüler sollen ihre Schule selbst mitgestalten.}
\end{textebox}
« Dans les écoles allemandes, des fenêtres sont détruites et des murs barbouillés. Les auteurs sont surtout des jeunes, parce qu'ils s'ennuient ou cherchent la reconnaissance. À cause des nombreux incidents, des caméras sont installées. Mais la prévention est plus importante : les élèves doivent participer eux-mêmes à l'aménagement de leur école. »

\textbf{Passifs employés :} \all{werden zerstört}, \all{beschmiert} (coordination de deux participes avec \all{werden} elliptique), \all{werden installiert} ; \textbf{connecteur de cause :} \all{weil}, \all{wegen}.

\textbf{Sujet B — Proposition de lettre modèle (environ 75 mots) :}
\begin{textebox}[Lettre formelle]
\all{Sehr geehrter Herr Bürgermeister,}

\all{in unserem Viertel in Yaoundé wurden letzte Woche zwei Bushaltestellen und mehrere Bänke beschädigt. Die Scheiben sind kaputt, und die Menschen müssen jetzt im Regen warten. Deshalb sind viele Bürger sehr verärgert.}

\all{Ich schlage vor, dass mehr Polizisten in unserem Viertel kontrollieren und dass Kameras installiert werden. Außerdem wäre ein Projekt mit den Jugendlichen sinnvoll, denn wer mitarbeitet, zerstört nicht.}

\all{Mit freundlichen Grüßen}

\all{Amina Mbarga}
\end{textebox}
« Monsieur le Maire, dans notre quartier de Yaoundé, deux abribus et plusieurs bancs ont été endommagés la semaine dernière. Les vitres sont cassées et les gens doivent maintenant attendre sous la pluie. C'est pourquoi beaucoup d'habitants sont très mécontents. Je propose que davantage de policiers patrouillent dans notre quartier et que des caméras soient installées. De plus, un projet avec les jeunes serait utile, car celui qui participe ne détruit pas. Cordialement, Amina Mbarga. »

\textbf{Points de vigilance :}
\begin{itemize}
\item Formule d'appel : \all{Sehr geehrter Herr Bürgermeister,} (minuscule au mot qui suit la virgule) ; formule de politesse : \all{Mit freundlichen Grüßen} (sans virgule finale).
\item Passifs : \all{wurden ... beschädigt} (Präteritum), \all{installiert werden} (infinitif passif après \all{dass}).
\item Connecteurs : \all{deshalb} (inversion : \all{sind viele Bürger}), \all{denn} (pas d'inversion).
\item \cle{Orthographe} : ne pas écrire \all{*geinstalliert} → \all{installieren} garde son participe \all{installiert} sans \all{ge-} (verbe en \all{-ieren}).
\item Contexte camerounais : \all{Yaoundé}, \all{Bushaltestellen}, \all{Regen} (saison des pluies) bien intégrés.
\end{itemize}
\end{corrige}

\newpage

\coursec{Fiche bilan}

\courssub{Fiche bilan}

\begin{popfigure}
\begin{tikzpicture}[
  mindmap,
  grow cyclic,
  every node/.style={concept, execute at begin node=\hskip0pt},
  concept color=popblue,
  level 1/.append style={level distance=4.5cm, sibling angle=360/7},
  level 2/.append style={level distance=3cm, sibling angle=45},
  text=white,
  font=\small\bfseries
]
\node[concept, scale=1.2, align=center]{Le vandalisme\\et le terrorisme}
  child[concept color=poporange] { node[concept, align=center]{Texte support\\Compréhension} }
  child[concept color=popgreen] { node[concept, align=center]{Lexique thématique\\Terror, Anschlag, Gewalt} }
  child[concept color=poppurple] { node[concept, align=center]{Grammaire 1\\Passif (werden + Partizip II)} }
  child[concept color=poppink] { node[concept, align=center]{Grammaire 2\\Cause \& Conséquence} }
  child[concept color=popgold] { node[concept, align=center]{Expression\\Réponses, Résumé, Production} }
  child[concept color=popdark] { node[concept, align=center]{Méthodologie\\Commentaire guidé} }
  child[concept color=popblue] { node[concept, align=center]{Contexte\\Cameroun / D-A-CH} };
\end{tikzpicture}
\legende{Carte mentale de la leçon ALA25 — Chapitre 5 : Citoyenneté}
\end{popfigure}

\textbf{Lexique clé (allemand / français) — articles, pluriels, familles de mots}\par

\begin{center}
\begin{tabularx}{\linewidth}{|>{\bfseries}l|X|X|}\hline
\rowcolor{popblueL} \textbf{Deutsch} & \textbf{Français} & \textbf{Famille / Exemple} \\ \hline
der Terror, - & le terrorisme & terrorisieren (v), terroristisch (adj), der Terrorist, die Terroristin \\ \hline
der Vandalismus, - & le vandalisme & der Vandale, die Vandalen, vandalieren (v), vandalistisch (adj) \\ \hline
der Anschlag, die Anschläge & l'attentat & der Bombenanschlag, verüben (einen Anschlag), der Anschlagsversuch \\ \hline
die Gewalt (kein Pl.) & la violence & gewalttätig (adj), der Gewalttäter, die Gewalttat, gewaltsam \\ \hline
beschädigen (reg.) & endommager & der Schaden, die Beschädigung, der Sachschaden \\ \hline
die Polizei (kein Pl.) & la police & der Polizist, die Polizistin, polizeilich, das Polizeirevier \\ \hline
die Sicherheit, - & la sécurité & sicher (adj), sichern (v), das Sicherheitsgefühl, die Sicherheitskraft \\ \hline
der Täter, die Täter & l'auteur (d'un crime) & die Tat, tätlich, täterlos, der Haupttäter \\ \hline
das Opfer, die Opfer & la victime & das Todesopfer, opfern (v), die Opferhilfe \\ \hline
die Prävention, - & la prévention & präventiv, vorbeugen (v, trennbar), die Präventionsarbeit \\ \hline
der Notruf, die Notrufe & l'appel d'urgence & die Notrufnummer (110 / 112), den Notruf wählen \\ \hline
verletzen (reg.) & blesser & die Verletzung, schwer verletzt, der Verletzte \\ \hline
\end{tabularx}
\end{center}

\vspace{0.5cm}
\textbf{Grammaire 1 : Le passif simple (das Prozessualpassiv) — Règles à connaître par cœur}\par

\begin{center}
\begin{tabularx}{\linewidth}{|>{\bfseries}l|X|X|}\hline
\rowcolor{popblueL} \textbf{Temps} & \textbf{Structure} & \textbf{Exemple (aktiv $\rightarrow$ passiv)} \\ \hline
Präsens & \all{werden} (konjugiert: \all{werde, wirst, wird, werden, werdet, werden}) + \all{Partizip II} & \all{Die Polizei verhaftet den Täter.} $\rightarrow$ \all{Der Täter wird verhaftet.} \\ \hline
Präteritum & \all{werden} (Prät.: \all{wurde, wurdest, wurde, wurden, wurdet, wurden}) + \all{Partizip II} & \all{Man beschädigte das Gebäude.} $\rightarrow$ \all{Das Gebäude wurde beschädigt.} \\ \hline
Perfekt & \all{sein} (konjugiert: \all{ist / sind}) + \all{Partizip II} + \all{worden} & \all{Man hat den Anschlag verhindert.} $\rightarrow$ \all{Der Anschlag ist verhindert worden.} \\ \hline
Futur I & \all{werden} (konjugiert) + \all{Partizip II} + \all{werden} & \all{Man wird die Sicherheit erhöhen.} $\rightarrow$ \all{Die Sicherheit wird erhöht werden.} \\ \hline
\end{tabularx}
\end{center}

\begin{attention}[Pièges fréquents francophones]
\begin{itemize}
\item \textbf{Oublier \all{worden} au Perfekt :} \all{ist verhindert} (Zustandspassif = état) $\neq$ \all{ist verhindert worden} (Vorgangspassif = action). \emph{Erreur fatale au Bac.}
\item \textbf{Accord du participe :} En allemand, le participe passé \textbf{ne s'accorde jamais} en genre ni en nombre (\all{die Tür ist geöffnet}, \all{die Türen sind geöffnet}).
\item \textbf{Agent (qui agit) :} Utiliser \all{von} + \textbf{Datif} (\all{von der Polizei}, \all{von einem Unbekannten}). \all{durch} + Accusatif = moyen/instrument sans intention (\all{durch den Sturm}, \all{durch eine Bombe}).
\item \textbf{Verbes intransitifs (sans COD) :} Pas de passif personnel possible (\all{schlafen, sterben, passieren}). Seul le passif impersonnel existe : \all{Es wird getanzt / Es wurde geschlafen}.
\item \textbf{Verbes à double accusatif (ex. \all{lehren, fragen, kosten}) :} Seul l'accusatif de la personne devient sujet au passif : \all{Er fragt mich den Weg} $\rightarrow$ \all{Ich werde den Weg gefragt}.
\end{itemize}
\end{attention}

\vspace{0.5cm}
\textbf{Grammaire 2 : Cause et Conséquence (Kausalität und Konsequenz)}\par

\begin{center}
\begin{tabularx}{\linewidth}{|>{\bfseries}l|X|X|}\hline
\rowcolor{popblueL} \textbf{Fonction} & \textbf{Moyens linguistiques (Allemand)} & \textbf{Structure / Place du verbe} \\ \hline
\textbf{CAUSE} & \all{weil} + verbe à la \textbf{fin} (subordonnée) & \all{Er bleibt zu Hause, weil er krank ist.} \\ \cline{2-3}
& \all{da} + verbe à la \textbf{fin} (subordonnée, formel/écrit, cause connue) & \all{Da es regnet, bleibe ich zu Hause.} \\ \cline{2-3}
& \all{denn} + verbe en \textbf{2e position} (phrase principale, coordination) & \all{Er bleibt zu Hause, denn er ist krank.} \\ \cline{2-3}
& \all{wegen} + \textbf{Génitif} (nom, langue soutenue) / \all{wegen} + Datif (parlé) & \all{Wegen des Regens...} / \all{Wegen dem Regen...} \\ \cline{2-3}
& \all{durch} + \textbf{Accusatif} (moyen/instrument direct) & \all{Durch den Sturm wurde das Dach abgedeckt.} \\ \hline
\textbf{CONSEQUENCE} & \all{deshalb / darum / deswegen} + verbe en \textbf{2e position} (phrase principale) & \all{Es regnet, deshalb bleibe ich zu Hause.} \\ \cline{2-3}
& \all{so dass} + verbe à la \textbf{fin} (subordonnée consécutive) & \all{Es regnet stark, so dass ich nass werde.} \\ \cline{2-3}
& \all{daher / folglich} (formel, écrit) + verbe en \textbf{2e position} & \all{Es regnet; daher / folglich bleibe ich zu Hause.} \\ \hline
\end{tabularx}
\end{center}

\begin{methode}[Commentaire de texte guidé — 5 étapes]
\begin{enumerate}
\item \textbf{Lecture globale :} Identifier le thème, le type de texte (article, rapport, témoignage), la source, l'intention de l'auteur (informer, alerter, convaincre).
\item \textbf{Lexique :} Surligner les mots-clés du thème (\all{Terror, Gewalt, Polizei, Prävention, Opfer...}), déduire le sens du contexte, noter les familles de mots.
\item \textbf{Structure :} Repérer les paragraphes (introduction, développement, conclusion), les connecteurs logiques (\all{zuerst, dann, weil, deshalb, schließlich, aber}) et les temps de verbe (Präteritum = récit, Präsens = actualité/généralités).
\item \textbf{Analyse détaillée :} Répondre aux questions (W-Fragen : \all{Wer? Was? Wann? Wo? Warum? Wie?}), citer le texte avec guillemets allemands (\all{Zitat: „...“}) et numéro de ligne.
\item \textbf{Production :} Résumer en 5--6 phrases (discours indirect : \all{Konjunktiv I} — \all{er sage, er habe...} si possible, sinon \all{Präteritum/Perfekt}), donner son avis personnel structuré (\all{Meiner Meinung nach... / Ich finde, dass... / Zum einen..., zum anderen...}).
\end{enumerate}
\end{methode}

\begin{aretenir}[Checklist avant le Bac]
\begin{itemize}
\item $\square$ Je connais le vocabulaire clé avec \textbf{articles} et \textbf{pluriels} (\all{der Terror, die Anschläge, die Polizei, der Vandalismus, das Opfer}).
\item $\square$ Je sais former le \textbf{passif} au Présens (\all{werden + Partizip II}) et au Prétérit (\all{wurden + Partizip II}).
\item $\square$ Je sais former le \textbf{Perfekt passif} correctement : \all{ist / sind + Partizip II + worden} (jamais \all{geworden}).
\item $\square$ J'utilise \textbf{\all{von} + Datif} pour l'agent et \textbf{\all{durch} + Accusatif} pour le moyen/instrument.
\item $\square$ Je distingue \textbf{activ} (sujet agit) et \textbf{passiv} (sujet subit) et je sais transformer une phrase active en passive (COD $\rightarrow$ Sujet, Sujet $\rightarrow$ \all{von}-Objet ou suppression).
\item $\square$ J'exprime la \textbf{cause} avec \all{weil / da} (verbe à la fin) ou \all{denn} (verbe 2e position) ou \all{wegen} + \textbf{Génitif}.
\item $\square$ J'exprime la \textbf{conséquence} avec \all{deshalb / darum / deswegen / folglich} (verbe 2e position) ou \all{so dass} (verbe à la fin).
\item $\square$ Je respecte l'\textbf{ordre des mots} : verbe en 2e position (phrase principale), verbe à la fin (subordonnée \all{weil, dass, so dass, da}).
\item $\square$ J'écris les \textbf{noms avec majuscule} et les \textbf{Umlaute} (ä, ö, ü) / \textbf{ß} correctement (\all{Straße, Größe, Vandalismus}).
\item $\square$ Je sais \textbf{résumer} un texte court (5--6 phrases) en utilisant le discours indirect (\all{Konjunktiv I} : \all{er sage, er habe, er werde}).
\item $\square$ Je connais les numéros d'urgence D-A-CH : \textbf{110 (Polizei)}, \textbf{112 (Feuerwehr / Notarzt)}.
\end{itemize}
\end{aretenir}

\findecours

\end{document}
